\documentclass[11pt,a4paper]{article}
\usepackage[T1]{fontenc}
\usepackage[utf8]{inputenc}
\usepackage[dvipsnames]{xcolor}
\usepackage{amsmath,amssymb,amsthm,mathtools}
\usepackage{newpxtext,newpxmath}
\usepackage[a4paper,margin=25mm,headheight=15pt]{geometry}
\usepackage{microtype}
\usepackage{booktabs,longtable,array,multirow,float}
\usepackage{enumitem,fancyhdr,titlesec}
\usepackage{listings}
\usepackage[most]{tcolorbox}
\usepackage{xurl}
\usepackage[colorlinks=true,linkcolor=MidnightBlue,citecolor=MidnightBlue,urlcolor=MidnightBlue]{hyperref}
\usepackage[nameinlink,noabbrev]{cleveref}
\definecolor{Ink}{HTML}{18324C}
\definecolor{Accent}{HTML}{315D66}
\definecolor{Pale}{HTML}{F0F5F7}
\titleformat{\section}{\Large\bfseries\color{Ink}}{\thesection}{0.65em}{}
\titleformat{\subsection}{\large\bfseries\color{Accent}}{\thesubsection}{0.65em}{}
\titleformat{\subsubsection}{\normalsize\bfseries}{\thesubsubsection}{0.65em}{}
\pagestyle{fancy}
\fancyhf{}
\fancyhead[L]{\small Finite-prefix corrections and dual certificates}
\fancyhead[R]{\small ProveIt research report}
\fancyfoot[C]{\thepage}
\renewcommand{\headrulewidth}{0.3pt}
\setlength{\parindent}{1.3em}
\setlength{\parskip}{0.25em}
\setlist{itemsep=0.25em,topsep=0.4em}
\emergencystretch=2em
\allowdisplaybreaks[2]
\newtheorem{theorem}{Theorem}[section]
\newtheorem{lemma}[theorem]{Lemma}
\newtheorem{proposition}[theorem]{Proposition}
\newtheorem{corollary}[theorem]{Corollary}
\theoremstyle{definition}
\newtheorem{definition}[theorem]{Definition}
\newtheorem{example}[theorem]{Example}
\newtheorem{question}[theorem]{Research question}
\theoremstyle{remark}
\newtheorem{remark}[theorem]{Remark}
\newcommand{\R}{\mathbb R}
\newcommand{\Q}{\mathbb Q}
\newcommand{\N}{\mathbb N}
\newcommand{\Z}{\mathbb Z}
\newcommand{\bA}{\boldsymbol a}
\newcommand{\bX}{\boldsymbol X}
\newcommand{\bone}{\boldsymbol 1}
\newcommand{\spr}{\operatorname{spr}}
\newcommand{\supp}{\operatorname{supp}}
\newcommand{\trunc}[2]{[\xi^{\leq #1}]#2}
\newcommand{\tail}[2]{[\xi^{>#1}]#2}
\newcommand{\repo}{\texttt{ProveIt}}
\newcommand{\sha}{\texttt{e21766d04c2b8a9b2cdba0cd43e563024b3bd1b9}}
% [write] Marker for text added when this manuscript was written into ProveIt
% (batch 36, manuscript 06); everything unmarked is the delivered text.
\newcommand{\wtag}{\textsf{\small[write]}}
% [write] Lean declaration and file names, breakable at any character (xurl).
\DeclareUrlCommand\lean{\urlstyle{tt}}
% Part II (batch 40, manuscript 05): the manuscript's one-argument
% coefficient operators, its defect tail R_N, and its norm.
\newcommand{\strunc}[1]{[\xi^{\le #1}]}
\newcommand{\shigh}[1]{[\xi^{>#1}]}
\newcommand{\satR}{\mathcal R}
\newcommand{\norm}[1]{\left\lVert#1\right\rVert}
\newtcolorbox{resultbox}{colback=Pale,colframe=Accent,boxrule=0.6pt,arc=1mm,left=2.5mm,right=2.5mm,top=2mm,bottom=2mm}
\lstset{basicstyle=\ttfamily\footnotesize,breaklines=true,columns=fullflexible,keepspaces=true,frame=single,rulecolor=\color{gray!40},showstringspaces=false}
\hypersetup{pdftitle={Finite-Prefix Corrections and Dual Certificates for Polyomino Growth},pdfauthor={AI-assisted research report prepared for Vladimir Reshetnikov},pdfsubject={An exact 4.498 upper bound and certificates for convolution recurrences; Part II: spectral saturation of the exact-prefix hierarchy}}

\begin{document}
\hypersetup{pageanchor=false}
\begin{titlepage}
\thispagestyle{empty}
\vspace*{0mm}
\noindent{\color{Accent}\large\scshape ProveIt research series}\par
\vspace{5mm}
\noindent{\color{Ink}\Huge\bfseries Finite-Prefix Corrections\par and Dual Certificates\par for Polyomino Growth}\par
\vspace{4mm}
\noindent{\LARGE An exact bound $\lambda\leq 4.498$\par and a certificate theorem for convolution recurrences}\par
\vspace{5mm}
\noindent{\large Research report prepared for Vladimir Reshetnikov}\par
\noindent September 2026 \quad $\cdot$ \quad AI-assisted research draft\par
\vspace{4mm}
\begin{abstract}
We develop two extensions of the seventeen-neighborhood polyomino project in
\repo. First, exact small-size occurrence counts are incorporated into a
positive convolution system by subtracting its finite coefficient defect.
A nonnegative tail formulation makes this correction rigorous without assuming
convergence of any generating function. Enumerating the seventeen marked
neighborhoods through size $18$ gives an exact rational certificate for
\[
 \boxed{\lambda\leq\frac{2249}{500}=4.498<4.5,}
\]
improving the repository's $4.5235$ upper bound.
Second, balanced integer weights on monomials produce finite arithmetic
certificates of divergence. We prove completeness for productive,
strongly connected, nonlinear positive polynomial systems, thereby providing
a certificate-based answer to Bui's Question~2 for this class. An explicit
$53$-weight certificate and the repository's primal certificate bracket the
uncorrected recurrence growth between $4.52349$ and $4.5235$.
We also prove a spectral stability criterion for the prefix hierarchy and give
a counterexample in which arbitrarily long exact prefixes still fail to
recover the true exponential growth. Full proofs, exact data, enumerators,
rational verifiers, and proposed research directions are supplied.
\end{abstract}
\vfill
\begin{resultbox}
\small
\textbf{Status and scope.} The new bound is a computer-assisted mathematical
result, not a newly kernel-checked Lean theorem. The new arguments are proved
in this report; the finite enumeration and arithmetic are reproducible.
The geometric recurrence lemma is imported from Bui and the pinned repository.
The work has not been independently refereed. General geometric-programming
duality and polyomino enumeration algorithms are prior tools, not claimed
inventions of this report.

\smallskip
\wtag\ In \repo\ the geometric recurrence lemma is more than imported: Lean
proves Bui's Lemma~4 recurrences for the actual marked counts
(\lean{geometricPublishedBuiRecurrences}), and the project's Lean-verified
bound is $\lambda\leq 9047/2000=4.5235$. The bound $\lambda\leq 2249/500=4.498$
of this report is computer-assisted and is \emph{not} formalized. See
\cref{kfp:sec:formal-status}.

\smallskip
[Added 29 September 2026, batch 40:] Part~II (\cref{kfp:sat:sec:source}),
also unformalized, proves that no exact prefix takes the fixed map's
majorant growth below $4.3149$: a limit of the method, \emph{not} a lower
bound on $\lambda$ and \emph{not} an improvement on $4.498$.
\emph{Tempting false readings:} ``$\lambda\geq4.3149$''; ``$4.3149<4.498$ is
a better bound''. Both are false.
\end{resultbox}
\end{titlepage}

\pagenumbering{roman}
\hypersetup{pageanchor=true}
{\small\setlength{\parskip}{0pt}\tableofcontents}
\clearpage
\pagenumbering{arabic}

\phantomsection\addcontentsline{toc}{part}{Part I.\texorpdfstring{\quad}{ } Finite-prefix corrections and dual certificates}
\section{Problem selection, sources, and main contributions}
\label{kfp:sec:scope}

\subsection{Why this repository thread}
The \repo\ repository contains a concrete project on fixed square-lattice
polyominoes at
\begin{center}\small
\texttt{Combinatorics/Polyominoes/KlarnerConstant/}.
\end{center}
Its existing endpoint is an upper bound of $9047/2000=4.5235$ for Klarner's
constant. The relevant sources were read at commit
\begin{center}\small\sha.\end{center}
In particular, the project README, the independent rational verifier,
\texttt{Patterns.lean}, and \texttt{GeometricComplete.lean} fix the mathematical
and computational interface used here \cite{repo,reporeadme,repocert,repopatterns,repocomplete}.
The last of these files assembles the seventeen geometric recurrences for
actual marked polyomino counts and then applies the $4.5235$ certificate.
We inspected those source statements; we did not run a fresh Lean build.

The published input is Vuong Bui's \emph{A convolutional approach to bounding
the number of polyominoes}, arXiv:2511.00461v2, dated 6 May 2026
\cite{bui}. Its Lemma~4 and Appendix~B give the seventeen-neighborhood
recurrence system. Its Section~3, Question~2 asks for compact, readily
verifiable lower-growth certificates for convolution recurrences.
This is a particularly suitable target: the geometric inequalities are
already available, the numerical improvement can have an exact proof
boundary, and the certificate question has a precise published formulation.

\paragraph{\wtag\ Earlier upper bounds.}
The delivered text compares its bound only with Bui's $4.5238$ and the
repository's $4.5235$. For context, the repository's research note on the
$4.5235$ bound \cite{repobound} records the upper-bound history as follows:
Eden obtained $6.75$ \cite{eden}; Klarner and Rivest successively obtained
$5$, $4.83$ and $4.649551$, the last by the procedure of \cite{klarnerrivest};
Barequet and Shalah reached $4.5252$ using over twenty-three trillion local
configurations \cite{barequetshalah}; Bui's seventeen neighborhood types then
gave $4.5238$ \cite{bui}, and the repository's exact supersolution for Bui's
unchanged system gave $9047/2000=4.5235$, which is kernel-checked in Lean
(\cref{kfp:sec:formal-status}). The bound $4.498$ of \cref{kfp:thm:main}
continues this list, but it is computer-assisted and not formalized.
None of the earlier methods used exact marked prefixes of the kind
introduced here; this is a statement about the cited sources, not a priority
search.

\subsection{What is proved here}
The first result concerns actual polyominoes. If $A_n$ is the number of
translation classes of $n$-cell polyominoes, then the data and certificate in
this report give
\begin{equation}\label{kfp:eq:headline}
 A_n\leq\frac{927884613}{10^9}\left(\frac{2249}{500}\right)^n
 \quad(n\geq1),\qquad
 \lambda\leq\frac{2249}{500}.
\end{equation}
The factor $927884613/10^9$ is the $G$ coordinate of the supplied rational
supersolution. No numerical eigenvalue or floating-point root appears in the
verification of \eqref{kfp:eq:headline}.

The second result concerns the equality recurrences used as majorants.
We give a theorem of alternatives: under explicit productivity,
irreducibility, and nonlinearity hypotheses, either a positive
supersolution exists or a balanced integer-monomial certificate disproves its
existence. The two alternatives cannot occur together. This supplies the
requested type of lower-growth certificate for a class containing Bui's
seventeen-variable system. The explicit example proves
\begin{equation}\label{kfp:eq:original-bracket}
 4.52349\leq\mu_{\rm original}\leq4.5235.
\end{equation}
Here $\mu_{\rm original}$ is the exponential growth of the \emph{uncorrected
equality recurrence}, not the polyomino growth constant $\lambda$.
In particular, the lower inequality in \eqref{kfp:eq:original-bracket} is
\emph{not} a lower bound on $\lambda$.

The third result describes both the power and a limitation of exact-prefix
corrections. The corrected majorants decrease coefficientwise as more exact
terms are supplied. They eventually admit certificates at any evaluation
point where the linearized system at the true series has spectral radius
less than one. Conversely, a supercritical linearization obstructs the
whole hierarchy under a mild nonzero-tail condition. A scalar example has
true growth $1$ but corrected-majorant growth tending to $3$.
[Added 29 September 2026, batch 40: Part~II classifies the equality case
for an irreducible true-profile Jacobian, identifies the limiting radius of
the hierarchy, and proves that for the fixed Bui map every corrected
majorant has growth at least $4.3149$ (\cref{kfp:sat:thm:critical,kfp:sat:thm:limit,kfp:sat:thm:bui}).
That number is a limit of the method, not a lower bound on $\lambda$.]

\subsection{Novelty and dependency boundaries}
The proposed research contributions are the explicit $4.498$ polyomino bound,
the finite-prefix application and its exact data, the arithmetic response to
Bui's certificate question, and the accompanying analysis of the prefix
hierarchy. The examined repository snapshot and primary sources did not
contain this $4.498$ certificate or this marked-prefix calculation. That is a
limited novelty assessment, not an exhaustive priority determination.

Several ingredients are established mathematics. The geometric inequalities
come from Bui. Enumeration by canonical include/exclude growth is in the
tradition of Redelmeier \cite{redelmeier}. Positive monomial inequalities,
logarithmic convexification, and entropy-like dual expressions belong to
geometric programming \cite{duffin,boyd,bv}. We do not present those tools as
new discoveries. The proof of the certificate alternative is included to
make its hypotheses, strictness, and rational verification completely
explicit in the recurrence setting.

\begin{resultbox}
\textbf{Three distinct levels.}
The geometric recurrence theorem is an imported mathematical input.
The new prefix and dual-certificate theorems have written proofs here.
The headline numerical bound additionally uses the finite enumeration in
Appendix~\ref{kfp:app:data}, whose source programs and exact outputs are supplied.
Arithmetic verification is not a substitute for verifying enumeration.

\smallskip
\wtag\ In \repo\ the first level is stronger than ``imported'': it is a
kernel-checked Lean theorem about the actual counts
(\cref{kfp:sec:formal-status}). The second and third levels are not
formalized.
\end{resultbox}

\subsection{\wtag\ Notation conventions}
\label{kfp:sec:notation}
This subsection was added when the manuscript was written into \repo. No
delivered symbol was renamed and no normalization was changed; the
conventions below only fix how the delivered symbols are to be read.

\paragraph{Coordinate letters.} The seventeen upper-case letters of
\eqref{kfp:eq:order} name the marked neighborhood types and their counting
sequences ($C_n,\ldots,Z_n$, also written $a_{i,n}$ with $i\in\mathcal I$);
the matching lower-case letters $c,\ldots,z$ are the coordinates of the
vector $\bX$ in \eqref{eq:map} and the row names of \cref{kfp:tab:dual}. A
coordinate is always a bare letter or a subscript index ($v_G$, $a_{D,7}$,
$\Phi_P$). Many of these letters are reused for other objects, which are
recognized by a subscript $N$, by calligraphic type, or by the local context
listed here:
\begin{center}\small
\begin{tabular}{@{}>{\raggedright\arraybackslash}p{0.2\textwidth}>{\raggedright\arraybackslash}p{0.76\textwidth}@{}}
\toprule
Letter & Other meanings in this report\\
\midrule
$P$ & $P_N=\trunc{N}{\bA}$, the exact prefix; a polyomino $P$ (\cref{kfp:sec:system,kfp:prop:enumeration})\\
$D$ & $D_N$, the finite defect; the defect polynomial $D$ in \eqref{kfp:eq:sensitivity}\\
$F$ & $F(X)=\Phi(\zeta,X)$ with components $F_i$, and $f_i(y)=\log F_i(e^y)-y_i$ in \eqref{kfp:eq:h} (\cref{kfp:sec:complete})\\
$X$, $d$ & the unknowns $X=(X_1,\ldots,X_d)$ of a general system $\Phi\in\R_{\geq0}[\xi,X_1,\ldots,X_d]^d$ (\cref{kfp:sec:prefix,kfp:sec:dual,kfp:sec:complete}), e.g.\ $X\in\R_{>0}^d$ in the sets $S_K$ of \cref{kfp:lem:compact}; the bold $\bX$ is the vector of \eqref{eq:map}\\
$T$, $Y$ & $T_N=\trunc{N}{\Phi(\xi,P_N)}$; $Y_N=B_N-P_N$\\
$E$ & forbidden offset set; the enumerator's excluded set; the defect improvement in \eqref{kfp:eq:sensitivity}\\
$R$, $S$ & required offset set, radii $R_i$; the enumerator's occupied set, the sets $S_K$ of \cref{kfp:lem:compact}\\
$U$ & the enumerator's frontier; $U=\bA-P_N$ and a generic subsolution (\cref{kfp:lem:formal,kfp:thm:majorant})\\
$H$ & the nonlinear part in $\Phi=L\bX+H$; the half-plane $\mathcal H$\\
$B$ & the offset sets $B$, $B_4$; the majorants $B_N$; a least fixed point $B$ (\cref{kfp:sec:dual})\\
$C$ & the dual product $\mathcal C$; a constant $C$ (\cref{kfp:thm:stability}); the normalizer $C=\sum_i\ell_i\tau_i$\\
$h$, $q$ & $h(y)$ of \eqref{kfp:eq:h}, $h=(I-J)^{-1}\bone$, a row index; the argument of $\log q$, a degree $q\geq2$\\
$v,w,y,p$, $t,u,x,e$, $c,g,r$ & certificate vector $v$, $w=v-P_N(\zeta)$, log-coordinates $y$, $p_N=P_N(\zeta)$, reals $t,u$, scalar unknown $x$, unit vectors $e_i$, coefficients $c_{ik}$, a nearest point $g$, row sums $r_i$\\
$b$ & $b_n$ (proof of \cref{kfp:lem:formal}); exponents $b_{ik}$ in \eqref{kfp:eq:general-map}; forcing vector $b_N$ (\cref{kfp:sec:stability})\\
\bottomrule
\end{tabular}
\end{center}

\paragraph{Collisions with the \repo\ project.} The project's research note
\cite{repobound} and its Lean files use the same $\lambda$, the same
$A(n)=A_n$, the same seventeen upper-case sequence names with lower-case
coordinates, and the same evaluation point $\zeta$; its numeric map $\Phi_\zeta$
is $\Phi(\zeta,\cdot)$ here. Two symbols differ:
\begin{itemize}
\item The project writes $b_N=(c_N,\ldots,z_N)$ for the \emph{$\zeta$-weighted
 prefix profile} and $\mathcal S_N=\sum_{n\leq N}S(n)\zeta^n$ for one sequence.
 That profile is $p_N=P_N(\zeta)$ here. \emph{Tempting false reading:} the
 $b_N$ of \cref{kfp:sec:stability} is the project's $b_N$. It is not; it is
 the high-degree forcing vector $\bigl(\tail{N}{\Phi(\xi,P_N)}\bigr)(\zeta)$.
\item The project writes $a$ for its rational supersolution; that is $v$ here,
 while $a=\bA(\zeta)$ in \cref{kfp:sec:stability}.
\end{itemize}
The symbol $\mu$ (the growth of an equality recurrence, $\mu_{\rm original}$
for Bui's uncorrected system) occurs only in this report. \emph{Tempting false
reading:} $\mu_{\rm original}\geq4.52349$ is a lower bound on $\lambda$. It is
not (\cref{kfp:sec:dual}). The convention $A_0=0$ of \cref{kfp:sec:system}
differs from the value $a(0)=1$ of OEIS A001168; positive-size terms agree.

\subsection{\wtag\ What \repo\ has formalized, and what it has not}
\label{kfp:sec:formal-status}
This report belongs to \repo's research-report collection and continues the
formal project \texttt{Combinatorics/Polyominoes/KlarnerConstant/}; being
filed next to, or citing, a Lean development confers no formal status on it.
That project proves the following in Lean (namespace
\lean{LeanProofs.KlarnerConstant}; files under
\lean{Combinatorics/Polyominoes/KlarnerConstant/Lean/KlarnerConstant/}, cited
below by file name), with no \texttt{sorry}, project axiom or
\texttt{native\_decide}; its assumption audit is \texttt{Audit.lean}. Line
numbers are those of the pin; the project's files are unchanged between the
pin and the placement commit.
\begin{itemize}
\item \emph{Bui's Lemma~4 for the actual counts.}
 \lean{geometricPublishedBuiRecurrences} (\lean{GeometricComplete.lean:63})
 proves \lean{PublishedBuiRecurrences} for
 \lean{geometricCoefficientProfile}: the
 seventeen recurrences for $n\geq2$, the published degree-one values,
 nonnegativity and vanishing at zero, for the finite marked-occurrence counts
 of the patterns in \texttt{Patterns.lean}. Coefficient by coefficient this is
 the inequality \eqref{kfp:eq:imported}, which the delivered text calls
 imported; reading the Lean structure as \eqref{kfp:eq:imported} in
 formal-series notation is not itself a formal step.
\item \emph{The five partitions \eqref{kfp:eq:partitions}} as equalities:
 \lean{buiF_occurrenceCount_eq_g_add_p}, \lean{buiG_occurrenceCount_eq_e_add_q},
 \lean{buiH_occurrenceCount_eq_d_add_s}, \lean{buiR_occurrenceCount_eq_y_add_w}
 and \lean{buiT_occurrenceCount_eq_x_add_v}
 (\lean{GeometricLinear.lean:279-312}).
\item \emph{$G$-domination} $A_n\leq G_n$:
 \lean{fixedPolyominoCount_le_geometricCoefficientProfile_g}
 (\lean{GeometricProfile.lean:371}). The other half $G_n\leq nA_n$ of
 \eqref{kfp:eq:gdom} is not formalized; the upper bound does not need it.
\item \emph{Supermultiplicativity and Fekete} \eqref{kfp:eq:lambda}:
 \lean{fixedPolyominoCount_supermultiplicative} and
 \lean{tendsto_realNthRoot_fixedPolyominoCount_growthSup}
 (\lean{GeometricComplete.lean:97}).
\item \emph{The Lean-verified bound $\lambda\leq9047/2000=4.5235$}:
 \lean{certificate_isSupersolution} (\lean{Certificate.lean:232}),
 \lean{fixedPolyominoCount_le_9047_div_2000_pow} (\lean{GeometricComplete.lean:69}),
 its real cast (\lean{:76}) and
 \lean{growthSup_fixedPolyominoCount_le_9047_div_2000} (\lean{:88}).
\end{itemize}
The one human boundary of that development is the transcription of Bui's
page-15 diagrams into the offset tables of \texttt{Patterns.lean}.

Nothing new in this report is formalized: not \cref{kfp:lem:formal},
\cref{kfp:thm:majorant,kfp:thm:upper,kfp:thm:nested}, not the size-$\leq18$
table, not \cref{kfp:thm:main} ($\lambda\leq2249/500=4.498$, which is
computer-assisted), and none of \cref{kfp:sec:dual,kfp:sec:complete,kfp:sec:dual-example,kfp:sec:stability}.
[Added 29 September 2026, batch 40: nor anything in Part~II, including the
$4.3149$ barrier of \cref{kfp:sat:thm:bui}.]
The project's endpoint theorems cannot simply be re-instantiated: they are
hard-wired to $\zeta=2000/9047$ (\lean{certificateZeta}) in
\lean{PrefixRecurrence.g_le_certificate} and \lean{g_lt_one}
(\lean{Recurrence.lean:321, 327}),
\lean{CoefficientProfile.g_lt_9047_div_2000_pow} (\lean{:457}) and
\lean{dominatedCoefficient_le_9047_div_2000_pow} (\lean{:477}). Only
\lean{buiIterate_le_supersolution} (\lean{:283}) and
\lean{PrefixRecurrence.le_supersolution} (\lean{:309}) are generic in
$\zeta$. A formal $4.498$ bound therefore needs, beyond the plan of
\cref{kfp:sec:verification}, a prefix-corrected variant of that
weighted-prefix engine and a computable bridge to the noncomputable
\lean{BuiNeighborhood.occurrenceCount} (\lean{Patterns.lean:140}).
As a formalization target, not a claim: the certificates for $N=7,\ldots,10$
in \cref{kfp:tab:progress} ($4.5233$ down to $4.5194$) already improve on
$4.5235$ and need exact marked counts only through size $7$ to $10$.

The trust boundary of \cref{kfp:thm:main} is thus: Lean-proved recurrences
(relative to \texttt{Patterns.lean}), plus an exhaustive enumeration that is
not kernel-checked (regenerated at intake, \cref{kfp:app:provenance}), plus
exactly replayed rational arithmetic (Python, not Lean), plus the written and
unrefereed proofs of \cref{kfp:sec:prefix}.

\section{Polyomino counts and the positive polynomial system}
\label{kfp:sec:system}

\subsection{Counting conventions}
A polyomino is a nonempty finite edge-connected subset of $\Z^2$.
Two such subsets are equivalent precisely when one is a translate of the
other. Rotations and reflections are not identified. Holes are allowed.
Let $A_n$ be the number of equivalence classes with $n$ cells, for $n\geq1$.
For the recurrence calculation we put $A_0=0$ and all marked size-zero
counts equal to zero. This convention differs from the optional empty-object
term $A_0=1$ in some sequence tables; positive-size terms are unchanged.

The standard concatenation argument gives $A_{m+n}\geq A_mA_n$ for positive
$m,n$. Fekete's lemma then gives
\begin{equation}\label{kfp:eq:lambda}
 \lambda=\lim_{n\to\infty}A_n^{1/n}=\sup_{n\geq1}A_n^{1/n}.
\end{equation}
These are background inputs also implemented in the repository's
\texttt{Concatenation.lean} and \texttt{Asymptotic.lean}
\cite{bui,reporeadme,repocomplete}.

A neighborhood pattern consists of a required offset set $R$ containing
$(0,0)$ and a disjoint forbidden offset set $E$. An occurrence at a cell $a$
of a polyomino $P$ means
\[
 a+R\subseteq P,\qquad (a+E)\cap P=\varnothing.
\]
The marked count for a pattern is the sum of its occurrences over all fixed
polyominoes of the given size. Counting occurrences, rather than merely
polyominoes which have an occurrence, is essential.

Write $C_n,D_n,\ldots,Z_n$ for the seventeen counts, in the coordinate order
\begin{equation}\label{kfp:eq:order}
 \mathcal I=(C,D,E,F,G,H,P,Q,R,S,T,U,V,W,X,Y,Z).
\end{equation}
To avoid confusing a coordinate name with a function, the polynomial map
will be denoted by $\Phi$.

\subsection{Exact offset specification}
Let
\[
 B=\{(-1,-1),(0,-1),(1,-1)\},\qquad
 B_4=B\cup\{(2,-1)\}.
\]
In \cref{kfp:tab:patterns}, $1,2,3$ in the required column mean respectively
$\{(0,0)\}$, $\{(0,0),(1,0)\}$, and
$\{(0,0),(1,0),(2,0)\}$. A blank lattice position imposes no condition.
These coordinate sets agree with the repository's pattern table and Bui's
page-15 neighborhood diagrams \cite{repopatterns,bui}.

\begin{table}[htbp]
\centering\small
\begin{tabular}{@{}ccl@{}}
\toprule
Type & Required & Forbidden offsets\\
\midrule
$C$ & 1 & $B\cup\{(-1,1),(1,1),(-1,0),(1,0)\}$\\
$D$ & 1 & $B\cup\{(-1,1),(-1,0),(1,0)\}$\\
$E$ & 1 & $B\cup\{(-1,0),(1,0)\}$\\
$F$ & 1 & $B$\\
$G$ & 1 & $B\cup\{(-1,0)\}$\\
$H$ & 1 & $B\cup\{(-1,1),(-1,0)\}$\\
$P$ & 2 & $\{(0,-1),(1,-1),(2,-1)\}$\\
$Q$ & 2 & $B\cup\{(-1,0)\}$\\
$R$ & 2 & $B_4\cup\{(-1,1),(-1,0)\}$\\
$S$ & 2 & $B\cup\{(-1,1),(-1,0)\}$\\
$T$ & 2 & $B_4\cup\{(-1,0)\}$\\
$U$ & 2 & $B_4$\\
$V$ & 3 & $B_4\cup\{(-1,0)\}$\\
$W$ & 3 & $B_4\cup\{(-1,1),(-1,0)\}$\\
$X$ & 2 & $B_4\cup\{(-1,0),(2,0)\}$\\
$Y$ & 2 & $B_4\cup\{(-1,1),(-1,0),(2,0)\}$\\
$Z$ & 2 & $B_4\cup\{(-1,1),(2,1),(-1,0),(2,0)\}$\\
\bottomrule
\end{tabular}
\caption{The seventeen marked neighborhood patterns.}
\label{kfp:tab:patterns}
\end{table}

Every polyomino has a $G$ occurrence at the leftmost cell of its bottom row.
There are at most $n$ possible marked cells. Consequently
\begin{equation}\label{kfp:eq:gdom}
 A_n\leq G_n\leq nA_n.
\end{equation}
The first six size-one counts equal one; the other eleven equal zero.

\subsection{The imported coefficient inequalities}
Use the formal variable $\xi$ and coordinate vector
$\bX=(c,d,e,f,g,h,p,q,r,s,t,u,v,w,x,y,z)$. Define
\begin{equation}\label{eq:map}
\begin{aligned}
\Phi_C&=\xi+\xi e,&\Phi_D&=\xi+\xi g,&\Phi_E&=\xi+\xi f,\\
\Phi_F&=g+p,&\Phi_G&=e+q,&\Phi_H&=d+s,\\
\Phi_P&=eh+qd+xr+vy+uyz,\\
\Phi_Q&=\xi g+\xi ge+\xi^2u+\xi^2tg+\xi^2ru,\\
\Phi_R&=y+w,\\
\Phi_S&=\xi g+\xi e^2+\xi^2t+\xi^2xg+\xi^2yu,\\
\Phi_T&=x+v,\\
\Phi_U&=dh+sd+yr+wy+uz^2,\\
\Phi_V&=\xi s+\xi^2g^2+\xi^2te+\xi^2rt,\\
\Phi_W&=\xi s+\xi^2eg+\xi^2xe+\xi^2yt,\\
\Phi_X&=\xi d+\xi^2g+\xi^2u,\\
\Phi_Y&=\xi c+\xi^2g+\xi^2t,\\
\Phi_Z&=\xi c+\xi^2e+\xi^2x.
\end{aligned}
\end{equation}
There are $53$ monomials, each with coefficient one. The order displayed
within each row is also the order of the dual weights in
\cref{kfp:tab:dual}. In the first two displayed lines, each coordinate has its
own two-term row.

For the true generating series
$\bA_i(\xi)=\sum_{n\geq1}a_{i,n}\xi^n$, Bui's geometric lemma is the
coefficientwise inequality
\begin{equation}\label{kfp:eq:imported}
 \bA\preceq\Phi(\xi,\bA).
\end{equation}
All products in this equation are formal Cauchy products. Thus convergence is
not part of the statement. Equations with a factor $\xi^2$, for example,
encode a two-cell shift in the corresponding convolution index.

The geometric proof of \eqref{kfp:eq:imported} uses occupancy partitions,
injective leaf deletions, and decompositions into connected marked pieces.
We import that lemma rather than silently treating a numerical recurrence
fit as a geometric proof. For example, the $C$ row deletes its marked cell,
whose only possible neighbor is above it. The $P$ and $U$ rows encode
multi-piece decompositions; these are not inferred from the finite table.
The repository supplies their separate finite-injection modules
\cite{bui,repocomplete}.
\wtag\ Those modules prove the lemma in Lean for the actual counts
(\lean{geometricPublishedBuiRecurrences}); so, within \repo, the input
\eqref{kfp:eq:imported} is a kernel-checked theorem relative to the offset
tables of \texttt{Patterns.lean}, not an imported assumption
(\cref{kfp:sec:formal-status}).

The five same-size partitions are in fact exact:
\begin{equation}\label{kfp:eq:partitions}
 F=G+P,\quad G=E+Q,\quad H=D+S,\quad R=Y+W,\quad T=X+V.
\end{equation}
For example, a type-$G$ anchor either has no right neighbor, yielding type
$E$, or has a right neighbor, yielding type $Q$. Similar left/right occupancy
splits give the other identities. They will serve as independent consistency
checks on the occurrence table.
\wtag\ The five identities are Lean theorems in \repo\
(\lean{GeometricLinear.lean:279-312}); the table's agreement with them is
checked by the Python verifier, not by Lean.

\section{Exact-prefix correction: the general theorem}
\label{kfp:sec:prefix}

\subsection{Proper positive polynomial systems}
\begin{definition}
A \emph{proper positive polynomial system} is a map
$\Phi\in\R_{\geq0}[\xi,X_1,\ldots,X_d]^d$ such that
$\Phi(0,0)=0$ and the nonnegative matrix
\[
 L=\Phi_{\bX}(0,0)
\]
is nilpotent. Here and below vector inequalities are coordinatewise.
\end{definition}

Nilpotence handles same-size linear dependencies. In \eqref{eq:map}, the
only such dependencies are the five rows in \eqref{kfp:eq:partitions}; their
directed graph has no cycle. All remaining occurrences of an unknown are
multiplied either by a positive power of $\xi$ or by another unknown.
Thus the system is proper. One possible degree-by-degree evaluation order is
\[
 C,D,E,Q,S,V,W,X,Y,Z,P,U,G,F,H,T,R.
\]
The formal development does not require this particular order.

\begin{lemma}[Formal fixed points and comparison]\label[lemma]{kfp:lem:formal}
A proper positive polynomial system has a unique fixed point in
$\xi\R[[\xi]]^d$. Its coefficients are nonnegative, and the iterates from
zero increase coefficientwise and stabilize at each finite degree.
If $U\in\xi\R_{\geq0}[[\xi]]^d$ satisfies
$U\preceq\Phi(\xi,U)$, then $U$ is coefficientwise at most this fixed point.
\end{lemma}
\begin{proof}
Write $\Phi=L\bX+H$. Because the unknown series have zero constant terms,
the coefficient of degree $n$ in $H(\xi,\bX)$ depends only on unknown
coefficients of degrees strictly less than $n$. If these have already been
chosen, the coefficient vector $b_n$ in a fixed point must satisfy
\[
 (I-L)b_n=h_n.
\]
For some $k$, $L^k=0$, so
\[
 (I-L)^{-1}=I+L+\cdots+L^{k-1}\geq0.
\]
This proves existence, uniqueness, and nonnegativity by induction on $n$.
For a subsolution, the same argument gives
$(I-L)u_n\leq h_n(u_1,\ldots,u_{n-1})$; multiplication by the nonnegative
inverse and induction prove comparison.

The iterates from zero are increasing by monotonicity. If all lower-degree
coefficients have stabilized, the remaining degree-$n$ error is multiplied
by $L$ at each further iteration, and therefore vanishes after at most $k$
steps. This proves finite-degree stabilization.
\end{proof}

The subsolution conclusion uses formal properness. It would be false for
arbitrary numerical subsolutions: $2\leq2^2$ does not imply that $2$ is below
the least nonnegative fixed point of $x\mapsto x^2$.

\subsection{The finite defect and the nonnegative tail map}
Suppose $\bA\in\xi\R_{\geq0}[[\xi]]^d$ is a true counting series satisfying
\eqref{kfp:eq:imported}. For $N\geq1$, set
\begin{equation}\label{kfp:eq:prefix-def}
 P_N=\trunc{N}{\bA},\qquad
 T_N=\trunc{N}{\Phi(\xi,P_N)},\qquad
 D_N=T_N-P_N.
\end{equation}
The coefficient of $\Phi(\xi,\bA)$ at a degree at most $N$ depends only on
$P_N$, so $D_N\succeq0$. It is precisely the finite prefix of the true
recurrence defect $\Phi(\xi,\bA)-\bA$.

The signed equation
\begin{equation}\label{kfp:eq:signed}
 B_N=\Phi(\xi,B_N)-D_N
\end{equation}
is a convenient shorthand, but its right-hand side is not a positive
polynomial in $\xi$. The correct positivity argument uses $Y_N=B_N-P_N$
and the map
\begin{equation}\label{kfp:eq:psi}
\begin{split}
 \Psi_N(\xi,Y)
 &=\Phi(\xi,P_N+Y)-T_N\\
 &=\Phi(\xi,P_N+Y)-\Phi(\xi,P_N)
   +\tail{N}{\Phi(\xi,P_N)}.
\end{split}
\end{equation}
Every coefficient of $\Psi_N$ in $\xi,Y_1,\ldots,Y_d$ is nonnegative.
Indeed, in the difference on the second line, the terms independent of $Y$
cancel, leaving only monomials containing at least one $Y$.
Moreover, $\Psi_N(0,0)=0$ and $(\Psi_N)_Y(0,0)=L$.

\begin{theorem}[Prefix-corrected majorant]\label{kfp:thm:majorant}
The map $\Psi_N$ has a unique fixed point
$Y_N\in\xi\R_{\geq0}[[\xi]]^d$. Its coefficients through degree $N$ vanish.
The series $B_N=P_N+Y_N$ satisfies
\[
 \bA\preceq B_N,\qquad
 \trunc{N}{B_N}=P_N,
\]
and solves \eqref{kfp:eq:signed}.
\end{theorem}
\begin{proof}
The fixed-point statement follows from \cref{kfp:lem:formal}. The constant-in-$Y$
part of $\Psi_N$ begins above degree $N$, so degree induction gives
$[\xi^n]Y_N=0$ for $n\leq N$.
Put $U=\bA-P_N\succeq0$. The low-degree part of
$\Phi(\xi,\bA)-\bA$ is $D_N$, and its remaining coefficients are
nonnegative. Therefore
\[
 U\preceq\Phi(\xi,P_N+U)-T_N=\Psi_N(\xi,U).
\]
Formal comparison yields $U\preceq Y_N$. Adding $P_N$ proves domination.
Substituting the fixed-point equation for $Y_N$ yields \eqref{kfp:eq:signed}.
\end{proof}

\begin{theorem}[Rational upper certificate]\label{kfp:thm:upper}
Let $\zeta>0$. Suppose a vector $v\in\R^d$ satisfies
\begin{equation}\label{kfp:eq:upper-conditions}
 v\geq P_N(\zeta),\qquad
 v\geq\Phi(\zeta,v)-D_N(\zeta).
\end{equation}
Then all coordinates of $B_N$ converge at $\zeta$, and
\[
 B_N(\zeta)\leq v,\qquad
 a_{i,n}\leq v_i\zeta^{-n}.
\]
For the polyomino system this implies
$A_n\leq v_G\zeta^{-n}$ and $\lambda\leq\zeta^{-1}$.
When the input prefix, $\zeta$, and $v$ are rational, verification of
\eqref{kfp:eq:upper-conditions} requires rational arithmetic only.
\end{theorem}
\begin{proof}
Let $w=v-P_N(\zeta)\geq0$. Equations \eqref{kfp:eq:prefix-def} and
\eqref{kfp:eq:psi} show that \eqref{kfp:eq:upper-conditions} is equivalent to
\[
 \Psi_N(\zeta,w)\leq w.
\]
The numerical iterates $w^{(0)}=0$,
$w^{(j+1)}=\Psi_N(\zeta,w^{(j)})$ therefore stay between $0$ and $w$.
The corresponding formal iterates are polynomials with nonnegative
coefficients. Their evaluation at $\zeta$ is exactly $w^{(j)}$.
For every fixed degree $M$, a sufficiently late formal iterate agrees with
$Y_N$ through degree $M$ by \cref{kfp:lem:formal}. Hence
\[
 \sum_{n=0}^M[\xi^n]Y_N\,\zeta^n\leq w.
\]
Monotone convergence of these bounded finite sums proves
$Y_N(\zeta)\leq w$. Adding $P_N(\zeta)$ proves the asserted bound for $B_N$.
Every nonnegative term of a bounded series is bounded by its sum, yielding
the coefficient estimate. Now use \eqref{kfp:eq:gdom} and \eqref{kfp:eq:lambda}.
\end{proof}

\begin{remark}[The tail condition is essential]
A solution of $v\geq\Phi(\zeta,v)-D_N(\zeta)$ alone need not be a valid
certificate. The subtraction can create extraneous branches or vectors below
the known prefix. The separate inequality $v\geq P_N(\zeta)$ is part of the
proof and is checked explicitly by the verifier.
\end{remark}

\subsection{Monotonicity of the hierarchy}
\begin{theorem}[More exact data never worsens the majorant]\label{kfp:thm:nested}
For every $N\geq1$,
\[
 \bA\preceq B_{N+1}\preceq B_N.
\]
Thus the radii of convergence of the corrected majorants are nondecreasing,
and their reciprocal exponential-growth bounds are nonincreasing.
Any real upper certificate for prefix $N$ is also an upper certificate for
prefix $N+1$ at the same evaluation point.
\end{theorem}
\begin{proof}
The first inequality is \cref{kfp:thm:majorant}. Through degree $N$, the two
majorants agree. At degree $N+1$, $B_{N+1}$ has the true coefficient and
$B_N$ has a coefficient at least as large. Above that degree, both obey the
same positive coefficient recurrence, because the defect polynomials have
no terms there. Degree induction using $(I-L)^{-1}\geq0$ proves the second
inequality. Coefficientwise domination implies the radius assertion.

If $v$ is an $N$-certificate, then the preceding theorem and
$\bA\preceq B_N$ show that $P_{N+1}(\zeta)\leq v$.
Also $D_{N+1}-D_N$ has nonnegative coefficients, so
$\Phi(\zeta,v)-D_{N+1}(\zeta)\leq v$.
Both certificate conditions for $N+1$ follow.
\end{proof}

For the present system the defect is zero through size $6$. Its size-$7$
coefficient, in the order \eqref{kfp:eq:order}, is
\begin{equation}\label{kfp:eq:first-defect}
 [\xi^7]D_7=(0,1,2,0,0,0,0,0,0,1,0,0,0,0,0,1,2).
\end{equation}
For example, $a_{D,7}=260$, whereas the deletion bound uses
$a_{G,6}=261$; the corresponding defect coefficient is $1$.
This first discrepancy explains why prefix lengths $1$ through $6$ do not
improve the original equality system.

\section{An exact upper bound below 4.5}
\label{kfp:sec:numeric}

\subsection{The certificate}
Take $N=18$ and
\[
 \zeta=\frac{500}{2249}.
\]
Use the exact marked counts in Appendix~\ref{kfp:app:data} to construct
$P_{18},T_{18},D_{18}$ by \eqref{kfp:eq:prefix-def}.
The vector $v$ is given in \cref{kfp:tab:upper}; every numerator in its second
column has common denominator $10^9$.

Define the residual and tail budget by
\begin{equation}\label{kfp:eq:slacks}
 \sigma_i=v_i-\Phi_i(\zeta,v)+D_{18,i}(\zeta),\qquad
 \eta_i=v_i-P_{18,i}(\zeta).
\end{equation}
The last two columns are conservative lower bounds:
\[
 \sigma_i\geq\frac{\text{third-column entry}}{10^{12}},\qquad
 \eta_i\geq\frac{\text{fourth-column entry}}{10^6}.
\]
They are obtained by flooring the exact rational numbers after scaling;
no printed decimal approximation is being used to decide a sign.

\begin{table}[htbp]
\centering\small
\begin{tabular}{@{}lrrr@{}}
\toprule
Coordinate & Numerator of $v_i$ & $10^{12}$-scaled slack & $10^6$-scaled tail\\
\midrule
$c$ & 347,752,792 & 99,824 & 28,076 \\
$d$ & 427,615,666 & 100,303 & 59,180 \\
$e$ & 564,832,205 & 100,345 & 123,441 \\
$f$ & 1,551,872,652 & 100,000 & 542,918 \\
$g$ & 927,884,613 & 99,000 & 260,176 \\
$h$ & 727,305,121 & 100,000 & 159,091 \\
$p$ & 623,987,939 & 99,749 & 282,741 \\
$q$ & 363,052,309 & 100,459 & 136,735 \\
$r$ & 234,092,913 & 100,000 & 63,133 \\
$s$ & 299,689,355 & 99,657 & 99,910 \\
$t$ & 283,350,363 & 100,000 & 89,179 \\
$u$ & 487,086,219 & 100,056 & 192,950 \\
$v$ & 119,025,613 & 99,875 & 52,209 \\
$w$ & 97,883,000 & 99,767 & 38,611 \\
$x$ & 164,324,650 & 100,195 & 36,970 \\
$y$ & 136,209,813 & 99,861 & 24,521 \\
$z$ & 112,439,657 & 100,134 & 14,724 \\
\bottomrule
\end{tabular}

\caption{The rational size-$18$ upper certificate and exact lower bounds on
its margins. All vector entries have denominator $10^9$.}
\label{kfp:tab:upper}
\end{table}

\begin{theorem}[Computer-assisted polyomino bound]\label{kfp:thm:main}
For the fixed-polyomino counting convention in \cref{kfp:sec:system},
\[
 A_n\leq\frac{927884613}{10^9}
          \left(\frac{2249}{500}\right)^n\quad(n\geq1),
 \qquad \lambda\leq4.498<4.5.
\]
\end{theorem}
\begin{proof}
The finite enumeration described below produces the true coefficients of
$P_{18}$ in Appendix~\ref{kfp:app:data}. Bui's recurrence lemma gives
\eqref{kfp:eq:imported}. Exact integer polynomial convolution constructs $D_{18}$.
Substituting the rational vector of \cref{kfp:tab:upper} into
\eqref{kfp:eq:slacks} gives
\[
 \min_i\sigma_i\geq\frac{99}{10^9}>0,\qquad
 \min_i\eta_i\geq\frac{14724}{10^6}>0.
\]
Thus \cref{kfp:thm:upper} applies. Its $G$ coordinate is
$v_G=927884613/10^9$, yielding the claimed inequalities.
\end{proof}

The strict residuals also show that the certificate is not balanced on a
rounding boundary. By continuity, the same vector works for a sufficiently
small increase of $\zeta$. We retain $2249/500$ as the stated bound because
it is simple and explicitly verified.

\subsection{Discovery is not verification}
The search for a candidate used floating-point solutions of
\[
 \Phi(\zeta,v)-D_N(\zeta)=v,
 \qquad \det(I-\Phi_{\bX}(\zeta,v))=0.
\]
These equations suggest where the equality envelope becomes critical.
For a slightly smaller evaluation point, we solved a perturbed fixed-point
equation with a small positive residual in every coordinate, then rounded the
candidate to denominator $10^9$. Such a solver may converge to an undesired
branch or misestimate a threshold. Neither event can invalidate an accepted
certificate: the verifier checks the two inequalities in
\eqref{kfp:eq:upper-conditions} independently from scratch.

\Cref{kfp:tab:progress} records the observed progression. The middle column
contains numerical critical-point estimates, not proved optimal values or
certified lower bounds. Every entry in the last column has a separate
exact rational certificate in the archive.

\begin{table}[htbp]
\centering\small
\begin{tabular}{@{}rrr@{}}
\toprule
Prefix size $N$ & Estimated equality growth & Certified upper bound\\
\midrule
$1$--$6$ & $4.523499228$ & $4.5235$ (original certificate)\\
7 & $4.523230947$ & $4.5233$\\
8 & $4.522434056$ & $4.5225$\\
9 & $4.521121227$ & $4.5212$\\
10 & $4.519368978$ & $4.5194$\\
11 & $4.517262575$ & $4.5173$\\
12 & $4.514881253$ & $4.5149$\\
13 & $4.512294495$ & $4.5123$\\
14 & $4.509560290$ & $4.5096$\\
15 & $4.506725656$ & $4.5068$\\
16 & $4.503828169$ & $4.5039$\\
17 & $4.500897596$ & $4.5009$\\
18 & $4.497957351$ & $4.4980$\\
\bottomrule
\end{tabular}

\caption{Progress under exact-prefix correction. Numerical estimates and
rigorous upper certificates are deliberately separated.}
\label{kfp:tab:progress}
\end{table}

The original recurrence's growth is already within $10^{-5}$ of the
repository's bound, as \cref{kfp:sec:dual-example} will prove. Consequently,
\cref{kfp:thm:main} does not arise from finding a better supersolution for the
unchanged system. The correction changes the equality envelope while
preserving its validity as an upper bound on the true counts.

\subsection{A reproducible finite enumeration}
Every translation class has a unique representative whose leftmost cell on
the bottom row is $(0,0)$. All its cells then lie in
\[
 \mathcal H=\{(x,y):y>0\}\cup\{(x,0):x\geq0\}.
\]
The main enumerator maintains a connected occupied set $S$, a finite ordered
frontier $U$ of available adjacent cells, and an excluded set $E$.
Initially $S=\varnothing$, $U=\{(0,0)\}$, and $E=\varnothing$.
At a call, take the last frontier cell $p$, include it in $S$, record the
resulting object, and recurse after adding newly available neighbors.
After that branch returns, remove $p$ and exclude it for all later branches
of the same call. The exclusions introduced by a call are removed when that
call returns. No quotient by rotations or reflections is performed.

\begin{proposition}[Enumeration correctness]\label[proposition]{kfp:prop:enumeration}
The include/exclude frontier procedure records each nonempty finite
edge-connected subset of $\mathcal H$ containing $(0,0)$ exactly once, up to
the specified size limit. Thus its totals are precisely the fixed-polyomino
counts.
\end{proposition}
\begin{proof}
Every recorded set is connected because each added cell is adjacent to the
previous occupied set, except for the initial origin. It stays in
$\mathcal H$ by the admissibility test.

Fix a target set $P$ of the required kind. At any intermediate state with
$S\subsetneq P$ and no excluded cell in $P$, connectedness implies that some
cell of $P\setminus S$ lies in the available frontier. In the order processed
by the current call, there is a unique first frontier cell belonging to $P$.
All cells processed before it can be excluded without affecting $P$.
The branch including that unique cell preserves the same invariant for $P$.
Induction reaches $S=P$, which is recorded.

The same argument proves uniqueness: the first processed frontier cell
belonging to $P$ determines the only branch that can contain $P$ without
excluding one of its cells. Applying this observation recursively determines
a unique path. Finally, the chosen anchor gives exactly one representative
of each translation class.
\end{proof}

For an occupied anchor, its membership in each pattern depends only on ten
nearby positions:
\[
 (-1,1),(1,1),(2,1),(-1,0),(1,0),(2,0),
 (-1,-1),(0,-1),(1,-1),(2,-1).
\]
These are encoded by a ten-bit mask. All $2^{10}$ masks are preclassified
against the required and forbidden sets of \cref{kfp:tab:patterns}.
The direct implementation scans all occupied anchors whenever an object is
recorded. The faster implementation maintains the pattern totals under cell
insertion and deletion. Inserting $p$ changes only the contribution of $p$
and those anchors $p-o$ for which $o$ is one of the ten offsets. Removing $p$
reverses exactly these changes. This local invariant proves that incremental
updates give the same occurrence count as a complete scan.

The size-$18$ computation records
\[
 A_{18}=1,540,820,542,\qquad
 \sum_{n=1}^{18}A_n=2,083,404,030
\]
objects across all recorded sizes. The unmarked totals through $18$ agree
with OEIS A001168 \cite{oeis}; those familiar totals are a consistency check,
not a replacement for the marked enumeration.

\subsection{Finite-machine safety and independent checks}
The C++ sources are restricted to size at most $18$. Any cell of an anchored
$n$-cell polyomino satisfies $|x|\leq n-1$ and $0\leq y\leq n-1$.
The padded array is wide and tall enough for these coordinates and for all
mask and neighbor offsets. Thus no boundary wrapping or out-of-range
access is needed by the algorithm.

The counters are unsigned $64$-bit integers. An a priori bound verifies that
overflow is impossible at these sizes. A deterministic depth-first spanning
tree of an anchored $n$-cell set can be encoded by a plane rooted tree and
directions for its $n-1$ forward steps. There are
$\operatorname{Cat}_{n-1}$ plane tree shapes. The first forward step has at
most four choices; every later forward step has at most three, since a
nonroot vertex cannot move back to its parent and the root has already used
one child direction. Consequently, for $n\geq2$,
\[
 A_n\leq4\cdot3^{n-2}\operatorname{Cat}_{n-1}.
\]
This overcounts heavily but is sufficient. Every marked count is at most
$nA_n$, and the largest bound for $n\leq18$ is
\[
 18\cdot4\cdot3^{16}\operatorname{Cat}_{17}
 =401816383504818480<2^{64}.
\]
Incremental counters count occurrences in the currently occupied finite set;
the insertion/deletion invariant also prevents underflow.

Two structurally different enumeration approaches were compared through size
$11$: the C++ include/exclude procedure and a Python procedure that repeatedly
adds each boundary cell, translates each shape to a canonical coordinate
tuple, and deduplicates using a set. All $187$ marked values and all unmarked
totals agreed. Direct scanning and incremental scanning in C++ agreed through
size $16$. The final size-$18$ output also agrees with the previous
size-$16$ table on every earlier entry.

These checks do not amount to independent formal certification of all
size-$18$ values. They locate the remaining computational trust boundary:
the correctness of the supplied enumerator, its implementation, and its
execution. The recurrence inequalities themselves are not learned from this
table. Five exact partition identities and all nonnegative recurrence defects
are checked again by the rational verifier.

\wtag\ \emph{Regeneration at intake.} The delivered validation record says
that no fresh size-$18$ run was made after the enumerators' final
command-line changes. When the manuscript was received by \repo, the
delivered \texttt{code/enumerate.cpp} was compiled afresh
(\texttt{c++ -O3 -std=c++17}, g++ 16.1.0) and run as \texttt{enumerate 18}
(8\,min\,23\,s on one core of the intake machine); after CSV parsing its
output equals the shipped \texttt{data/profiles\_18.csv} in all $19$ columns
of all $18$ rows (the files differ only in line endings). Fresh runs through
sizes $11$--$14$, $16$ and $17$ matched as well, and an independent set-based
Python enumerator, with pattern coordinates taken from the project's
\texttt{Patterns.lean} rather than from this package, matched all $10$
unmarked and $170$ marked entries through size $10$. So the size-$18$ table is
now reproduced from the final source. This does not change the boundary
stated above: the regeneration uses the same program and algorithm, the
rational verifier checks arithmetic and necessary consistency conditions
only, and the correctness of the enumeration rests on
\cref{kfp:prop:enumeration} and the incremental-update argument together
with this regeneration.

\section{Balanced-monomial certificates of divergence}
\label{kfp:sec:dual}

\subsection{The certificate format}
Consider a finite positive polynomial system
\begin{equation}\label{kfp:eq:general-map}
 \Phi_i(\xi,X)=\sum_{k\in K_i}c_{ik}\xi^{b_{ik}}X^{\alpha_{ik}},
 \quad c_{ik}\in\Q_{>0},\quad b_{ik}\in\N,
 \quad \alpha_{ik}\in\N^d.
\end{equation}
Fix a rational $\zeta>0$. We wish to certify that there is no
$X\in\R_{>0}^d$ satisfying $\Phi(\zeta,X)\leq X$.

Assign a nonnegative integer weight $m_{ik}$ to each monomial, not all zero,
and define row sums
\[
 r_i=\sum_{k\in K_i}m_{ik}.
\]
The weights are \emph{balanced} if
\begin{equation}\label{kfp:eq:balance}
 \sum_{i,k}m_{ik}\alpha_{ik,j}=r_j\qquad(1\leq j\leq d).
\end{equation}
Define the positive rational product
\begin{equation}\label{kfp:eq:dual-product}
 \mathcal C(\zeta,m)=
 \prod_{i,k:m_{ik}>0}
 \left(\frac{c_{ik}\zeta^{b_{ik}}r_i}{m_{ik}}\right)^{m_{ik}}.
\end{equation}
A row of weight zero contributes no factor. The entire certificate consists
of $\zeta$, the integer weights, and a verification that
$\mathcal C(\zeta,m)>1$.

\begin{theorem}[Balanced-monomial obstruction]\label{kfp:thm:dual}
If \eqref{kfp:eq:balance} holds and $\mathcal C(\zeta,m)>1$, then no positive
supersolution of $\Phi(\zeta,X)\leq X$ exists.
\end{theorem}
\begin{proof}
Suppose such a supersolution exists. For each row with $r_i>0$, put
$p_{ik}=m_{ik}/r_i$ on the positive-weight terms. Weighted arithmetic--geometric
mean gives
\begin{align*}
 1\geq\frac{\Phi_i(\zeta,X)}{X_i}
 &\geq\sum_{k:m_{ik}>0}
        c_{ik}\zeta^{b_{ik}}X^{\alpha_{ik}-e_i}\\
 &\geq\prod_{k:m_{ik}>0}
 \left(\frac{c_{ik}\zeta^{b_{ik}}
                  X^{\alpha_{ik}-e_i}}{p_{ik}}\right)^{p_{ik}}.
\end{align*}
Raise the inequality to $r_i$ and multiply over rows. The exponent of $X_j$
in the resulting right-hand side is
$\sum_{i,k}m_{ik}\alpha_{ik,j}-r_j=0$ by balance.
All unknown variables cancel, leaving
$1\geq\mathcal C(\zeta,m)>1$, a contradiction.
\end{proof}

The cancellation is the essential point. A certificate does not need to
supply a very large iterate, an approximate divergent orbit, or an
approximate critical point. Its inequalities quantify over every positive
candidate supersolution at once.

\begin{example}[Catalan recurrence]
For $\Phi(\xi,x)=\xi+x^2$, assign weight one to both terms.
Then the incoming exponent is $0+2=2=r$, and
\[
 \mathcal C(\zeta,m)=(2\zeta)(2)=4\zeta.
\]
Thus $\zeta>1/4$ immediately has a two-weight divergence certificate, recovering
the usual Catalan threshold from one arithmetic--geometric mean inequality.
\end{example}

\subsection{From nonexistence to a lower growth bound}
Let $B$ be the least formal fixed point of a proper positive system. If its
coordinates converge at $\zeta$, finite products of nonnegative convergent
series can be evaluated, giving
\[
 B(\zeta)=\Phi(\zeta,B(\zeta)).
\]
If the system is productive in the sense defined in the next section, this
fixed point is strictly positive. Therefore \cref{kfp:thm:dual} implies that at
least one coordinate fails to converge at $\zeta$.

Define the system's exponential growth by
\[
 \mu=\max_i\limsup_{n\to\infty}\bigl([\xi^n]B_i\bigr)^{1/n}.
\]
Equivalently, $1/\mu$ is the minimum of the coordinate radii of convergence.
Cauchy--Hadamard then gives
\begin{equation}\label{kfp:eq:dual-growth}
 \mu\geq\zeta^{-1}.
\end{equation}
We use the limsup formulation so that no unproved assertion about periodicity
or existence of individual root limits is needed.

For productive systems with a strongly connected dependency graph, the
coordinate radii coincide. Indeed, for every dependency edge $i\to j$,
choose one positive coefficient in each of the other factors of a monomial
containing $X_j$. The formal fixed-point equation implies
$B_i\succeq c\xi^kB_j$ for some $c>0$ and integer $k\geq0$.
Thus $R_i\leq R_j$. Following paths in both directions proves equality.

\begin{resultbox}
\textbf{Do not reverse the majorant inequality.}
The true polyomino series is below the equality solution, so an upper bound
on equality growth bounds $\lambda$ from above. A lower bound on equality
growth does not bound $\lambda$ from below. The dual certificate is useful
for certifying the limitation of a chosen recurrence model and for answering
the recurrence-growth question, not for reversing this implication.
\end{resultbox}

\section{Completeness of the certificate alternative}
\label{kfp:sec:complete}

\subsection{The hypotheses}
For this section fix $\zeta>0$ and write $F(X)=\Phi(\zeta,X)$.
The coefficients of $F$ are positive on its listed monomials.

\begin{definition}
The map $F$ is \emph{productive} if some finite iterate of $F$, starting from
zero, has every coordinate strictly positive. Its dependency graph has an
edge $i\to j$ if some monomial of $F_i$ contains $X_j$.
The map is \emph{strongly connected} if this graph is strongly connected,
and \emph{nonlinear} if at least one monomial has total degree at least two
in the unknowns.
\end{definition}

Productivity is stronger than merely having a constant term somewhere.
For example, $F_1=1+X_1X_2$, $F_2=X_1X_2$ has a strongly connected graph,
but iteration from zero never makes $X_2$ positive. This degeneracy is
excluded. All three conditions are finite structural checks on the monomial
list; they hold for \eqref{eq:map} and are replayed by the verifier.

\subsection{Compactness of bounded multiplicative residuals}
\begin{lemma}\label[lemma]{kfp:lem:compact}
If $F$ is productive, strongly connected, and nonlinear, then, for every
$K>0$, the set
\[
 S_K=\{X\in\R_{>0}^d:F(X)\leq KX\}
\]
is contained in a compact box $[a,b]\subset\R_{>0}^d$ and is closed in that
box. In particular, it is compact, possibly empty.
\end{lemma}
\begin{proof}
Every $X\in S_K$ satisfies $X\geq K^{-1}F(X)$. Iteration from zero of the
map $K^{-1}F$ therefore stays below $X$. Multiplication by a positive scalar
does not change which coordinates become positive at any iteration, so
productivity supplies a fixed vector $a>0$ with $X\geq a$ for every
$X\in S_K$.

For each edge $i\to j$, choose a monomial containing $X_j$. Lower-bound all
its other factors, including extra copies of $X_j$ if necessary, by the
corresponding coordinates of $a$. This gives a constant $c_{ij}>0$ such that
\[
 F_i(X)\geq c_{ij}X_j\quad(X\geq a).
\]
On $S_K$ it follows that $X_i\geq(c_{ij}/K)X_j$.
Strong connectivity now gives constants $u_i,v_i>0$ with
\[
 u_iX_1\leq X_i\leq v_iX_1
 \qquad(X\in S_K).
\]
Choose a nonlinear monomial of total degree $q\geq2$ in row $h$.
The lower comparisons give $F_h(X)\geq cX_1^q$ for some $c>0$,
whereas $F_h(X)\leq KX_h\leq Kv_hX_1$ on $S_K$.
Thus $X_1^{q-1}\leq Kv_h/c$, a uniform upper bound.
The other coordinates are then bounded above by their comparisons with
$X_1$. Finally, the polynomial inequalities are closed within this positive
compact box.
\end{proof}

\subsection{An elementary convex proof}
Put $X_j=e^{y_j}$ and define
\begin{equation}\label{kfp:eq:h}
 f_i(y)=\log F_i(e^y)-y_i,
 \qquad h(y)=\max_i f_i(y).
\end{equation}
Each $f_i$ is a log-sum-exp of affine functions and is convex. For
completeness, its Hessian is the covariance matrix of the monomial exponent
vectors under their positive softmax probabilities, and is therefore
positive semidefinite. Sublevel sets of $h$ correspond to the sets $S_K$ of
\cref{kfp:lem:compact}. Hence $h$ attains its global minimum.

\begin{theorem}[Complete rational certificate alternative]\label{kfp:thm:alternative}
Suppose $F$ has rational coefficients and is productive, strongly connected,
and nonlinear. Exactly one of the following holds:
\begin{enumerate}[label=(\roman*)]
 \item there is $X>0$ with $F(X)\leq X$;
 \item there are balanced nonnegative integer monomial weights, not all zero,
       for which the product \eqref{kfp:eq:dual-product} is greater than one.
\end{enumerate}
Here the coefficient $c_{ik}\zeta^{b_{ik}}$ in
\eqref{kfp:eq:dual-product} is simply the corresponding coefficient of $F$.
\end{theorem}
\begin{proof}
Mutual exclusion is \cref{kfp:thm:dual}. Suppose (i) fails. Let $y^*$ minimize
$h$, and put $\delta=h(y^*)$. Since a minimum is attained and no point has
$h\leq0$, we have $\delta>0$.

Let $I$ be the active rows, those with $f_i(y^*)=\delta$.
The origin belongs to the convex hull of their gradients. To see this
without invoking an optimization duality theorem, suppose otherwise and let
$g$ be the point of that compact convex hull nearest the origin.
Then $g\ne0$ and
$g\cdot\nabla f_i(y^*)\geq\|g\|^2$ for every active row.
A sufficiently small step in direction $-g$ strictly decreases every active
$f_i$. The inactive rows remain below the old maximum by continuity and
finiteness. This contradicts minimality of $h$.
Consequently there are numbers $\theta_i\geq0$, supported on $I$, such that
\[
 \sum_i\theta_i=1,\qquad
 \sum_i\theta_i\nabla f_i(y^*)=0.
\]

For each monomial let
\[
 p_{ik}=
 \frac{c_{ik}\zeta^{b_{ik}}e^{\alpha_{ik}\cdot y^*}}
      {F_i(e^{y^*})},\qquad
 \beta_{ik}=\theta_i p_{ik}.
\]
Then $\sum_k\beta_{ik}=\theta_i$, and the gradient identity is exactly
\[
 \sum_{i,k}\beta_{ik}\alpha_{ik,j}=\theta_j.
\]
Thus the real weights $\beta$ are balanced and have total mass one.
Their logarithmic product is
\begin{align*}
 \mathcal L(\beta)
 &=\sum_{i,k:\beta_{ik}>0}
 \beta_{ik}\log\left(
   \frac{c_{ik}\zeta^{b_{ik}}\theta_i}{\beta_{ik}}\right)\\
 &=\sum_i\theta_i\log F_i(e^{y^*})
   -\sum_j\theta_j y_j^*\\
 &=\sum_i\theta_i f_i(y^*)=\delta>0.
\end{align*}

It remains to replace real weights by integers. Restrict to the support of
$\beta$. On this support, balance and total mass one form a consistent
linear system with rational coefficients. Gaussian elimination supplies a
rational particular solution and a rational basis for the homogeneous
solution space. Approximating the real free parameters by rationals gives
balanced rational weights arbitrarily close to $\beta$, positive on the
same support. The displayed logarithmic expression is continuous there,
so a sufficiently close rational approximation still has positive value.
Multiplying by a common denominator gives integer weights. Scaling multiplies
the logarithmic product by that positive denominator, preserving strict
positivity. This proves (ii).
\end{proof}

\subsection{What this settles, and what it does not}
A positive numerical supersolution bounds the zero iterates and hence all
formal partial sums, by the same argument as in \cref{kfp:thm:upper}. Thus it
forces convergence at the evaluation point. For a proper productive system
in the stated class, divergence at a rational point therefore has a finite
balanced-monomial obstruction, and a finite obstruction
implies the lower-growth bound \eqref{kfp:eq:dual-growth}. Thus there is a
certificate alternative, not just a test based on watching a floating-point
iteration grow. This answers the certificate-existence aspect of Bui's
Question~2 for this class, including his seventeen-variable system.

The result does not assert that the original question for every conceivable
convolution, maximum, or signed recurrence has been settled. Nor does it give
a uniform polynomial bound on the bit length of the weights as a target
approaches the exact growth threshold. The explicit certificate below is
small and quickly checked; a worst-case certificate-size theory remains a
separate problem. The underlying convex mechanism is classical geometric
programming, here specialized and proved in an arithmetic certificate form
\cite{duffin,boyd,bv}.

\section{A compact dual certificate for the original recurrence}
\label{kfp:sec:dual-example}

\subsection{The integer weights}
For the uncorrected map \eqref{eq:map}, take
\[
 \zeta_{\rm dual}=\frac{100000}{452349}.
\]
Use the weights in \cref{kfp:tab:dual}. There are $53$ positive integers, with
sum $1,000,000$ and largest entry $150,463$. Direct integer addition verifies
all seventeen balance equations \eqref{kfp:eq:balance}.

\begin{table}[htbp]
\centering\small
\begin{tabular}{@{}lrl@{}}
\toprule
Row & $r_i$ & Monomial weights, in the order in (\ref{eq:map})\\
\midrule
$c$ & 4,857 & 3,080, 1,777\\
$d$ & 57,691 & 29,588, 28,103\\
$e$ & 242,698 & 93,305, 149,393\\
$f$ & 149,393 & 88,617, 60,776\\
$g$ & 248,519 & 150,463, 98,056\\
$h$ & 47,024 & 27,413, 19,611\\
$p$ & 60,776 & 39,670, 15,071, 3,712, 1,589, 734\\
$q$ & 113,127 & 63,384, 36,448, 7,450, 4,066, 1,779\\
$r$ & 6,370 & 3,666, 2,704\\
$s$ & 27,124 & 18,469, 6,430, 1,247, 679, 299\\
$t$ & 6,546 & 3,753, 2,793\\
$u$ & 11,651 & 7,354, 3,067, 759, 322, 149\\
$v$ & 4,382 & 2,413, 1,560, 289, 120\\
$w$ & 3,026 & 2,033, 796, 139, 58\\
$x$ & 8,357 & 4,786, 2,331, 1,240\\
$y$ & 7,427 & 4,155, 2,506, 766\\
$z$ & 1,032 & 702, 256, 74\\
\bottomrule
\end{tabular}

\caption{An explicit balanced-monomial certificate. The coefficient of every
monomial is one. Within each coordinate row, weights follow the term order
in \eqref{eq:map}.}
\label{kfp:tab:dual}
\end{table}

For example, the incoming $C$ exponent is contributed only by the
$\xi c$ terms in rows $Y$ and $Z$. Their weights are $4155$ and $702$,
whose sum is $4857=r_C$. Other balance equations count repeated variables
with multiplicity: the $z^2$ in row $U$ contributes twice its weight to the
incoming $Z$ exponent.

Exact rational interval evaluation proves
\begin{equation}\label{kfp:eq:dual-interval}
 \frac{1000668023141660979}{10^{18}}
 \leq \log\mathcal C(\zeta_{\rm dual},m)
 \leq \frac{1000668023155325499}{10^{18}}.
\end{equation}
In particular, $1<\log\mathcal C<1001/1000$, so
$\mathcal C>1$. The slightly generous unit-sized lower margin is merely a
convenient verification statement; scaling all weights scales the logarithm.

\begin{theorem}[Certified limitation of the original system]\label{kfp:thm:bracket}
The growth $\mu_{\rm original}$ of the original equality recurrence satisfies
\[
 \frac{452349}{100000}\leq\mu_{\rm original}\leq\frac{9047}{2000}.
\]
Consequently, no supersolution of the original system can establish a growth
upper bound smaller than $4.52349$.
\end{theorem}
\begin{proof}
Balance and \eqref{kfp:eq:dual-interval} give the lower bound by
\cref{kfp:thm:dual} and \eqref{kfp:eq:dual-growth}. The repository's rational
supersolution at $\zeta=2000/9047$, replayed by our verifier without a defect
correction, gives the upper bound. Finally, a supersolution at a smaller
claimed growth value would imply a radius larger than
$100000/452349$, contradicting the certified lower bound.
\end{proof}

The last statement can also be read directly: if a supersolution existed at
some $\zeta'>\zeta_{\rm dual}$, monotonicity of $\Phi$ in its first argument
would make it a supersolution at $\zeta_{\rm dual}$, which the dual
certificate forbids. This is an exact obstruction, not a conclusion drawn
from a numerical optimizer failing to improve the bound.

\subsection{Exact verification without enormous integer products}
Computing the rational product \eqref{kfp:eq:dual-product} literally is possible
in principle, but raising many rational factors to large integer weights is
unnecessary. Instead, enclose its logarithm using rational arithmetic.
For $1\leq q\leq2$, let $u=(q-1)/(q+1)$, so $0\leq u\leq1/3$. The identity
\[
 \log q=2\sum_{k\geq0}\frac{u^{2k+1}}{2k+1}
\]
follows by integrating the geometric series for $2/(1-u^2)$ from $0$ to $u$.
The remainder after $M$ terms satisfies
\begin{equation}\label{kfp:eq:log-error}
 0\leq \log q-2\sum_{k=0}^{M-1}\frac{u^{2k+1}}{2k+1}
 \leq\frac{2u^{2M+1}}{(2M+1)(1-u^2)}.
\end{equation}
All endpoints are rational when $q$ is rational.

For a general positive rational $q$, write $q=2^e w$ with integer $e$ and
$1\leq w<2$. Evaluate $\log w$ and $\log2$ by
\eqref{kfp:eq:log-error}; multiply the interval for $\log2$ by $e$, reversing
endpoint roles when $e<0$. This gives a rational enclosure for $\log q$.

The supplied verifier uses $M=16$. Each individual logarithm interval is
rounded outward to denominator $10^{18}$ before multiplication by its
integer weight. This optional outward rounding prevents huge denominators
when intervals from many factors are added. It cannot invalidate enclosure.
Adding the resulting weighted intervals gives exactly
\eqref{kfp:eq:dual-interval}. The same cached enclosure of $\log2$ can be reused;
there are at most $54\times16$ short series terms to evaluate.

The verification cost depends on the bit sizes of the rational monomial
factors, weights, and chosen precision, rather than on the number of iterates
required to witness numerical divergence. No logarithm supplied by a
floating-point library enters the proof. We make no universal claim that
certificate discovery or certificate size is polynomial in the reciprocal
distance to the growth threshold.

\subsection{How the weights were discovered}
The following construction is useful when a finite positive critical point
has been found numerically. Suppose
\[
 \tau=\Phi(\rho,\tau),\qquad
 J=\Phi_{\bX}(\rho,\tau),\qquad
 \ell^TJ=\ell^T,\qquad \ell>0.
\]
Normalize by $C=\sum_i\ell_i\tau_i$ and define
\[
 \beta_{ik}=
 \frac{\ell_i c_{ik}\rho^{b_{ik}}\tau^{\alpha_{ik}}}{C}.
\]
The row sums are $\ell_i\tau_i/C$, and the left-eigenvector identity shows
that the weights are balanced. At $\zeta=\rho$, their logarithmic product
is zero. At $\zeta>\rho$, it becomes
\[
 \left(\sum_{i,k}\beta_{ik}b_{ik}\right)\log(\zeta/\rho)>0
\]
provided at least one positively weighted term contains the size variable.
Rational approximation inside the balance equations then yields integer
weights. This is how the displayed certificate was found.

The approximate critical point and eigenvector are not inputs to the final
verifier. Only the resulting integers and rational evaluation point are.
The completeness proof in \cref{kfp:sec:complete} also does not require that a
numerical critical point be correct or even available.

\section{The limits of the exact-prefix hierarchy}
\label{kfp:sec:stability}

\subsection{A spectral criterion}
Fix a positive evaluation point $\zeta$ at which the true series $\bA$
converges coordinatewise. Put
\[
 a=\bA(\zeta),\quad p_N=P_N(\zeta),\quad
 J=\Phi_{\bX}(\zeta,a),\quad
 J_N=\Phi_{\bX}(\zeta,p_N),
\]
and define the nonnegative forcing vector
\[
 b_N=\bigl(\tail{N}{\Phi(\xi,P_N(\xi))}\bigr)(\zeta).
\]
At this evaluation point,
\begin{equation}\label{kfp:eq:tail-numeric}
 \Psi_N(\zeta,y)=b_N+
       \Phi(\zeta,p_N+y)-\Phi(\zeta,p_N).
\end{equation}
We have $p_N\uparrow a$ and $J_N\to J$. Moreover, $b_N\to0$:
coefficientwise $\Phi(\xi,P_N)\preceq\Phi(\xi,\bA)$, and the latter
series converges at $\zeta$. Its high-degree tail bounds $b_N$.

\begin{theorem}[Strict stability and obstruction]\label{kfp:thm:stability}
Assume the preceding setup.
\begin{enumerate}[label=(\alph*)]
\item If $\spr(J)<1$, then all sufficiently large $N$ admit a real upper
certificate at $\zeta$. If the polynomial data and $\zeta$ are rational,
the certificate can be chosen rational.
\item For any $N$, if $J_N$ is irreducible,
$\spr(J_N)\geq1$, and $b_N\ne0$, then no upper certificate at $\zeta$
exists for that $N$.
\item If $\spr(J)>1$, $J_N$ is irreducible for all sufficiently large $N$,
and $b_N\ne0$ for all sufficiently large $N$, then no finite-prefix member
of the hierarchy admits an upper certificate at $\zeta$.
\end{enumerate}
\end{theorem}
\begin{proof}
For (a), the nonnegative matrix geometric series gives
\[
 h=(I-J)^{-1}\bone>0,\qquad Jh=h-\bone.
\]
For sufficiently small $\varepsilon>0$, polynomial expansion in the
nonnegative direction $\varepsilon h$ gives, uniformly for $0\leq p_N\leq a$,
\[
 \Phi(\zeta,p_N+\varepsilon h)-\Phi(\zeta,p_N)
 \leq\varepsilon Jh+C\varepsilon^2\bone
\]
for some finite $C\geq0$. Choose $\varepsilon$ with
$C\varepsilon^2\leq\varepsilon/4$, and then $N$ sufficiently large that
$b_N\leq(\varepsilon/4)\bone$. Equation \eqref{kfp:eq:tail-numeric} yields
\[
 \Psi_N(\zeta,\varepsilon h)
 \leq\varepsilon h-(\varepsilon/2)\bone.
\]
Thus $p_N+\varepsilon h$ is a certificate with strict slack and strictly
positive tail budget. When the data and evaluation point are rational,
a sufficiently close rational vector preserves both strict inequalities.

For (b), nonnegative polynomial expansion gives
$\Psi_N(\zeta,y)\geq b_N+J_Ny$ for $y\geq0$.
Let $\ell>0$ be a left Perron vector of the irreducible matrix $J_N$,
with eigenvalue $r=\spr(J_N)\geq1$.
A certificate would give $y\geq b_N+J_Ny$, and hence
\[
 (1-r)\ell^Ty\geq\ell^Tb_N>0,
\]
which is impossible.

For (c), continuity of spectral radius gives $\spr(J_N)>1$ for all
sufficiently large $N$. Part (b) rules out certificates at all those levels.
By \cref{kfp:thm:nested}, an earlier certificate would remain a certificate at
all later levels. Therefore no earlier one exists either.
\end{proof}

The equality case $\spr(J)=1$ is deliberately not classified in general.
[Added 29 September 2026, batch 40: Part~II classifies it when $J$ is
irreducible: then, as for every $\spr(J)\geq1$, a certificate at $\zeta$
exists at level $N$ exactly when the full defect $\Phi(\xi,\bA)-\bA$ has no
coefficients above degree $N$ (\cref{kfp:sat:thm:critical}), so a nonzero
full-defect tail excludes every level even when every $J_N$ is subcritical.
The reducible critical case remains open (\cref{kfp:sat:sec:exceptions}).]
The nonzero forcing and irreducibility conditions in the obstruction are
not decorative: without them a critical linear map can have supersolutions.
For the Bui map, positive prefix coordinates make the dependency pattern
irreducible, and the $C$ row already supplies a nonzero high-degree forcing
term whenever the top $E$ coefficient is positive.

\subsection{An exact counterexample to asymptotic exactness}
\begin{theorem}[All finite prefixes may retain a growth gap]\label{kfp:thm:gap}
There is a scalar proper positive convolution inequality with a true series
of growth $1$ for which the exact-prefix majorants have growth tending to
$3$, not $1$.
\end{theorem}
\begin{proof}
Take
\[
 A(\xi)=\frac{\xi}{1-\xi}=\sum_{n\geq1}\xi^n,
 \qquad \Phi(\xi,x)=\xi+x^2.
\]
The coefficient inequality $A\preceq\xi+A^2$ holds: the coefficient of
$\xi^n$ in $A^2$ is $n-1$ for $n\geq2$.
The exact-prefix defect is
\[
 D_N(\xi)=\sum_{n=3}^N(n-2)\xi^n,
\]
with an empty sum for $N<3$. The corrected equality equation is
\begin{equation}\label{kfp:eq:toy}
 B_N=\xi+B_N^2-D_N(\xi).
\end{equation}
Let $c_N(t)=t-D_N(t)$ for real $t$.
Since
\[
 D_\infty(t)=\frac{t^3}{(1-t)^2},
\]
we have $D_\infty(1/3)=1/12$ and
$D_\infty'(t)\leq1$ on $0\leq t\leq1/3$, with equality only at the
right endpoint. A finite truncation has strictly smaller derivative there.
Thus $c_N$ is strictly increasing on this interval,
$c_N(1/4)\leq1/4$, and $c_N(1/3)>1/4$.
There is a unique
\[
 \rho_N\in[1/4,1/3)\quad\hbox{with}\quad c_N(\rho_N)=1/4.
\]
At $t=\rho_N$, the value $v=1/2$ solves \eqref{kfp:eq:toy}, and
\[
 P_N(\rho_N)<A(\rho_N)<1/2.
\]
It is therefore a valid upper certificate. For any
$t\in(\rho_N,1/3]$, the discriminant of
$x^2-x+c_N(t)=0$ is negative, so there is no finite real fixed point.
Hence the radius of $B_N$ is exactly $\rho_N$.

The truncated polynomials $D_N$ converge uniformly to $D_\infty$ on
$[0,1/3]$. Every limit point $\rho$ of the sequence $\rho_N$ therefore
satisfies
\[
 \rho-\frac{\rho^3}{(1-\rho)^2}=\frac14,
 \quad\hbox{equivalently}\quad (3\rho-1)^2=0.
\]
Thus $\rho_N\to1/3$, and the corrected growth tends to $3$.
The true series has radius $1$ and growth $1$.
\end{proof}

The same barrier is visible directly in \cref{kfp:thm:stability}:
\[
 \Phi_x(t,A(t))=\frac{2t}{1-t}
\]
crosses one at $t=1/3$, well before the true singularity at $t=1$.
Every finite coefficient of $B_N$ eventually agrees with $A$, but radii of
convergence do not pass continuously through this coefficientwise limit.
[Added 29 September 2026, batch 40: Part~II proves that this is the general
mechanism (\cref{kfp:sat:thm:limit}) and gives the rate for this example,
$\frac13-\rho_N\sim\frac29\sqrt{2N-1}\,3^{-N/2}$
(\cref{kfp:sat:sec:scalar}).]

This counterexample is not a disproof of Bui's separate Question~1 about
increasing neighborhood size with improved geometric recurrences. Here the
local polynomial map is held fixed and only a finite exact prefix is
substituted. Those are different hierarchies. It does show that an eventual
convergence claim for the present prefix method requires a stability argument
or additional geometric refinement.

\subsection{A conditional sensitivity formula for choosing improvements}
A useful guide for future searches is available when a smooth critical branch
has already been justified. Suppose an admissible defect improvement
$D+\varepsilon E$ has critical data $(\rho(\varepsilon),\tau(\varepsilon))$
satisfying
\[
 \Phi(\rho(\varepsilon),\tau(\varepsilon))
 -D(\rho(\varepsilon))-\varepsilon E(\rho(\varepsilon))
 =\tau(\varepsilon).
\]
At $\varepsilon=0$, let $\ell^TJ=\ell^T$ for
$J=\Phi_{\bX}(\rho,\tau)$. Differentiating and multiplying by $\ell^T$
cancels the unknown derivative of $\tau$, yielding
\begin{equation}\label{kfp:eq:sensitivity}
 \rho'(0)=
 \frac{\ell^TE(\rho)}
      {\ell^T\bigl(\Phi_\xi(\rho,\tau)-D'(\rho)\bigr)},
\end{equation}
provided the denominator is nonzero.
For a positive tail formulation with $\tau=P(\rho)+y$, the denominator
also equals $\ell^T\Psi_\xi(\rho,y)$, because the difference is
$\ell^T(J-I)P'(\rho)=0$. It is nonnegative, and is positive when some
size-dependent monomial contributes positively.

Formula \eqref{kfp:eq:sensitivity} suggests prioritizing verified defect
improvements by their Perron-weighted value. It is conditional on the stated
smooth critical family; it is not used in \cref{kfp:thm:main}. In particular,
it never authorizes subtracting an unproved or merely guessed defect.

\section{Verification, formalization, and reproducibility}
\label{kfp:sec:verification}

\subsection{The proof dependency chain}
The mathematical chain behind the new upper bound is
\[
\begin{gathered}
 \text{geometric recurrence lemma}
 \ +\ \text{exact marked prefix}\\
 \Longrightarrow\ \text{nonnegative corrected tail system}\\
 \Longrightarrow\ \text{rational supersolution at }500/2249\\
 \Longrightarrow\ \text{coefficient majorant}
 \Longrightarrow\ \lambda\leq2249/500.
\end{gathered}
\]
The new lower certificate for the old equality system follows a different
chain:
\[
\begin{gathered}
 \text{integer balance}\ +\ \text{rational log enclosure}\\
 \Longrightarrow\ \text{no positive supersolution}\\
 \Longrightarrow\ \mu_{\rm original}\geq452349/100000.
\end{gathered}
\]
These chains share the monomial map but not the logical direction of the
majorant comparison.

\subsection{Archive contents and commands}
The archive includes the article source and PDF, the exact marked profile,
twelve upper certificates, the dual certificate, the verification transcript,
and source code. The principal files are:
\begin{center}\small
\begin{tabular}{@{}ll@{}}
\toprule
File & Purpose\\
\midrule
\texttt{code/model.py} & Sparse map, patterns, exact polynomial operations\\
\texttt{code/verify.py} & Standard-library rational certificate replay\\
\texttt{code/enumerate.cpp} & Incremental marked enumeration through size $18$\\
\texttt{code/enumerate\_direct.cpp} & Direct-scanning comparison implementation\\
\texttt{code/research.py} & Independent Python enumeration and numerical exploration\\
\texttt{code/discover\_certificates.py} & Numerical search and rationalization\\
\texttt{data/profiles.json} & Exact marked and unmarked counts\\
\texttt{data/upper\_N18.json} & Main rational upper certificate\\
\texttt{data/dual\_original.json} & Balanced lower certificate for original model\\
\texttt{data/verification.json} & Exact arithmetic residuals and log interval\\
\bottomrule
\end{tabular}
\end{center}

From the archive root, the fast arithmetic audit is
\begin{lstlisting}
python3 code/verify.py --report data/verification.json
\end{lstlisting}
It uses only Python's standard library, arbitrary-precision integers, and
\texttt{fractions.Fraction}. It checks the monomial graph hypotheses,
nonnegative prefix defects, five occupancy partitions, all upper-certificate
residuals and tail budgets, all dual balance equations, the rational log
interval, and the repository's previous $4.5235$ supersolution.

Full regeneration of the main finite table is separate:
\begin{lstlisting}
c++ -O3 -std=c++17 code/enumerate.cpp -o enumerate
./enumerate 18 > regenerated.csv
python3 code/check_profile.py regenerated.csv
\end{lstlisting}
This visits more than two billion partial objects in total and is much slower
than certificate replay. The comparison tool checks all marked and unmarked
entries against the archived table. No external sequence database is needed
for either enumeration or certificate checking.

The numerical discovery scripts require NumPy, SciPy, and SymPy. They are
provided for further experimentation, but are outside the proof boundary.
The figures in the numerical-estimate column of \cref{kfp:tab:progress} should
not be promoted to exact assertions merely because those scripts print a
small residual or a successful solver status.

\wtag\ \emph{Layout in \repo.} The report directory keeps the delivered
\texttt{code/}, \texttt{data/} and \texttt{tex/} layout, so the commands above
run from the report root under the shipped names. Two files moved: the
delivered \texttt{Makefile} is \texttt{code/Makefile} (its recipes still use
root-relative paths, so run it as \texttt{make -f code/Makefile verify} from
the report root), and \texttt{requirements-discovery.txt} is
\texttt{data/requirements-discovery.txt}. The delivered PDF was replaced by a
build of this text, and the delivered checksum manifest is not shipped. On
the intake machine \texttt{python3} is spelled \texttt{py}. Because
\texttt{verify.py --report} rewrites \texttt{data/verification.json},
\texttt{make\_tables.py} rewrites \texttt{tex/*.tex}, and
\texttt{discover\_certificates.py} overwrites \texttt{data/upper\_N*.json} and
\texttt{data/dual\_original.json}, run them on a copy of the directory. At
intake, \texttt{verify.py} passed on such a copy and its output equalled
\texttt{data/verification.txt}.

\subsection{A focused Lean extension plan}
The most direct formal extension reuses
\texttt{geometricPublishedBuiRecurrences} and the existing domination of the
fixed-polyomino count by the $G$ coordinate. The genuinely new interfaces are
an exact finite prefix for the marked profile, the defect-corrected
weighted-prefix inequality, and the new rational vector.

A practical route can avoid formalizing complex singularity theory entirely.
First prove the finite convolution identity underlying
\eqref{kfp:eq:prefix-def}. Next prove a weighted-tail version of the repository's
monotone-prefix certificate theorem. Instantiate it with the new vector,
using ordinary exact rational reduction for every inequality. The existing
growth endpoint can then carry the pointwise exponential bound to $\lambda$.
The expensive part is a formally verified enumeration or a small proof object
certifying the marked prefix, not the seventeen final inequalities.
\wtag\ In the Lean files at the pin, the ``monotone-prefix certificate
theorem'' exists in $\zeta$-generic form only as
\lean{buiIterate_le_supersolution} and \lean{PrefixRecurrence.le_supersolution};
the endpoint theorems that carry a supersolution to a pointwise bound are
specialized to $\zeta=2000/9047$ (\cref{kfp:sec:formal-status}). The final
passage to $\lambda$ can use the generic \lean{growthSup_le_of_le_pow}
(\lean{Growth.lean:38}), since $v_G<1$ gives $A_n\leq(2249/500)^n$.

For dual certificates, the balance identity is finite integer arithmetic.
One can either formalize weighted arithmetic--geometric mean directly or
expand rational weights into an unweighted finite product argument. A
logarithmic checker with the remainder bound \eqref{kfp:eq:log-error} avoids
large powers and is a natural reusable component. The abstract completeness
theorem need not be formalized before the particular $53$-weight certificate
can be checked.

No Lean source claiming these new results is included. In particular, the
presence of a Lean proof of the old bound in the repository must not be
confused with a Lean proof of the new finite-prefix theorem or of the
size-$18$ data.

\section{Proposed research questions and next projects}
\label{kfp:sec:questions}

The following projects are intended as concrete continuations, not claims
already established in this report.

\begin{enumerate}[label=\textbf{\arabic*.},leftmargin=2em]
\item \textbf{Kernel-check the $4.498$ bound.}
Formalize the finite-prefix certificate lemma and a verified marked
polyomino enumerator, or design a compact reflective certificate for its
output. The target should be an unconditional endpoint with precisely the
same counting conventions as \texttt{GeometricComplete.lean}. This is the
highest-priority validation project.
\wtag\ \Cref{kfp:sec:formal-status} lists the Lean pieces that exist, the
$\zeta$-specialization that must be generalized, and a smaller first target
(the $N=7$ certificate, $\lambda\leq4.5233$, needs counts only through size
$7$).

\item \textbf{Generate longer marked prefixes without listing every object.}
Adapt a transfer-matrix or connectivity-state enumeration to accumulate all
seventeen occurrence counts, not merely $A_n$. The relevant challenge is
preserving marked local information while merging equivalent boundary
states. Exact counts of larger objects could strengthen the present bound
without changing the imported geometric proof.

\item \textbf{Certify lower bounds for the corrected envelopes.}
Expand or represent the positive maps $\Psi_N$ and construct balanced-monomial
certificates for them. This would replace the middle column of
\cref{kfp:tab:progress} by narrow rigorous two-sided brackets. Sparsity or an
arithmetic-circuit representation may avoid a large explicit expansion.

\item \textbf{Locate the actual stability barrier of the fixed Bui map.}
Bound $\spr(\Phi_{\bX}(\zeta,\bA(\zeta)))$ using certified lower and upper
bounds on the true marked generating functions. Determine whether the
prefix hierarchy's limiting obstruction lies substantially below the true
polyomino radius. The strict criteria in \cref{kfp:thm:stability} give concrete
sufficient tests, while the equality case remains open.
[Added 29 September 2026, batch 40: Part~II closes the irreducible equality
case (\cref{kfp:sat:thm:critical}) and proves that the limiting majorant
growth of the fixed map is at least $4.3149$ (\cref{kfp:sat:thm:bui}); where
the barrier actually lies, and whether it is below the true radius, remain
open (\cref{kfp:sat:sec:future}).]

\item \textbf{Optimize which geometric deficiencies to repair.}
Use Perron-weighted sensitivities as a search heuristic, but accept a repair
only after proving the corresponding injection or coefficient improvement.
A small new neighborhood resolving a heavily weighted overcount may be more
valuable than many additional unweighted states.

\item \textbf{Study Bui's neighborhood-refinement convergence question.}
Develop a hierarchy in which neighborhood size and geometric recurrences
are refined together. Prove convergence to the true growth constant, or
identify a rigorous obstruction. The scalar counterexample here rules out
one naive prefix-only argument; it does not settle the neighborhood question.

\item \textbf{Quantify certificate bit complexity.}
Given a rational target separated from the critical value by a known amount,
bound the sizes of balanced integer weights and the number of logarithm
terms required. Separate complexity of discovery from complexity of replay.
The present million-unit total weight is an example, not a worst-case bound.

\item \textbf{Extend the alternative beyond the irreducible nonlinear case.}
Decompose reducible systems into strongly connected blocks and analyze
linear, unproductive, and boundary blocks separately. A comprehensive theorem
should explain when weak infeasibility can occur and when a strict rational
obstruction still exists. Maximum operators require a different or enriched
certificate language.

\item \textbf{Obtain certified critical asymptotics.}
For selected corrected systems, prove uniqueness of the relevant positive
critical point, nondegeneracy, and the resulting coefficient asymptotics.
Such a result would justify numerical critical continuation as mathematics,
not merely as a source of candidate certificates.
[Added 29 September 2026, batch 40: Part~II proves general radius, height
and local-amplitude asymptotics as $N\to\infty$ at an interior nonlinear
crossing (\cref{kfp:sat:thm:sharp,kfp:sat:thm:branch}). It proves no
coefficient asymptotics and does not verify its hypotheses for the Bui
system, so this project stays open.]

\item \textbf{Transfer the method to other lattice animals and local grammars.}
Identify an existing positive recurrence system on a different lattice,
compute its first exact defects, and test whether a modest prefix correction
produces a useful improvement. The general theorems apply without change
once properness and the underlying combinatorial domination are proved.
\end{enumerate}

\section{Conclusion}
\label{kfp:sec:conclusion}
The repository's polyomino project supports more than a small numerical
retuning. Incorporating exact marked prefixes changes the majorant itself,
and the resulting rational certificate proves $\lambda\leq4.498$.
Balanced monomial weights provide the missing lower-certificate mechanism
for a natural class of convolution systems; their explicit use also proves
that the original recurrence cannot support the new bound without a
structural correction. Finally, the spectral analysis identifies why this
successful improvement must not be extrapolated automatically to asymptotic
exactness. The central next step is independent verification, especially
formal certification of the finite marked enumeration.
[Added 29 September 2026, batch 40: Part~II, which follows the appendices
of Part~I, makes that limitation exact.]

\appendix
\clearpage
\phantomsection\addcontentsline{toc}{part}{Appendices to Part I}
\section{Complete marked profiles through size eighteen}
\label{kfp:app:data}

The tables give the exact occurrence counts for the patterns in
\cref{kfp:tab:patterns}. All omitted size-zero entries are zero. The row $A_n$
counts translation classes; every other row counts marked occurrences.
These data, together with \eqref{eq:map}, determine every coefficient of
$D_{18}$ by finite integer convolution. The JSON and CSV forms contain the
same integers without typographical formatting.

\begin{table}[H]
\centering\footnotesize
\setlength{\tabcolsep}{4pt}
\renewcommand{\arraystretch}{1.12}
\begin{tabular}{@{}lrrrrrrrrr@{}}
\toprule
$n$ & 1 & 2 & 3 & 4 & 5 & 6 & 7 & 8 & 9\\
\midrule
$A_n$ & 1 & 2 & 6 & 19 & 63 & 216 & 760 & 2725 & 9910\\
\midrule
$C$ & 1 & 1 & 1 & 3 & 10 & 35 & 128 & 479 & 1819\\
$D$ & 1 & 1 & 2 & 6 & 20 & 71 & 260 & 973 & 3696\\
$E$ & 1 & 1 & 3 & 10 & 35 & 128 & 479 & 1821 & 7001\\
$F$ & 1 & 3 & 10 & 35 & 128 & 481 & 1841 & 7131 & 27856\\
$G$ & 1 & 2 & 6 & 20 & 71 & 261 & 982 & 3751 & 14481\\
$H$ & 1 & 2 & 5 & 15 & 51 & 182 & 670 & 2517 & 9592\\
$P$ & 0 & 1 & 4 & 15 & 57 & 220 & 859 & 3380 & 13375\\
$Q$ & 0 & 1 & 3 & 10 & 36 & 133 & 503 & 1930 & 7480\\
$R$ & 0 & 1 & 3 & 8 & 25 & 85 & 303 & 1110 & 4150\\
$S$ & 0 & 1 & 3 & 9 & 31 & 111 & 410 & 1544 & 5896\\
$T$ & 0 & 1 & 3 & 9 & 30 & 106 & 388 & 1451 & 5510\\
$U$ & 0 & 1 & 4 & 14 & 50 & 185 & 700 & 2686 & 10412\\
$V$ & 0 & 0 & 1 & 4 & 14 & 52 & 196 & 748 & 2884\\
$W$ & 0 & 0 & 1 & 4 & 13 & 46 & 168 & 625 & 2362\\
$X$ & 0 & 1 & 2 & 5 & 16 & 54 & 192 & 703 & 2626\\
$Y$ & 0 & 1 & 2 & 4 & 12 & 39 & 135 & 485 & 1788\\
$Z$ & 0 & 1 & 2 & 3 & 8 & 25 & 84 & 296 & 1080\\
\bottomrule
\end{tabular}

\caption{Exact profiles for sizes $1$ through $9$.}
\end{table}

\begin{table}[H]
\centering\footnotesize
\setlength{\tabcolsep}{2.25pt}
\renewcommand{\arraystretch}{1.12}
\begin{tabular}{@{}lrrrrrrrrr@{}}
\toprule
$n$ & 10 & 11 & 12 & 13 & 14 & 15 & 16 & 17 & 18\\
\midrule
$A_n$ & 36446 & 135268 & 505861 & 1903890 & 7204874 & 27394666 & 104592937 & 400795844 & 1540820542\\
\midrule
$C$ & 6979 & 26988 & 105011 & 410629 & 1612207 & 6351215 & 25091789 & 99372469 & 394381157\\
$D$ & 14190 & 54917 & 213865 & 837024 & 3289287 & 12969734 & 51285522 & 203287656 & 807484165\\
$E$ & 27144 & 105927 & 415489 & 1636445 & 6467194 & 25630675 & 101822518 & 405337049 & 1616425239\\
$F$ & 109503 & 432550 & 1715129 & 6821554 & 27199077 & 108673645 & 434959946 & 1743471996 & 6997283445\\
$G$ & 56348 & 220593 & 867716 & 3426361 & 13572506 & 53905206 & 214569185 & 855714279 & 3418228578\\
$H$ & 36929 & 143269 & 559155 & 2192740 & 8632427 & 34094245 & 135023622 & 535974728 & 2131792166\\
$P$ & 53155 & 211957 & 847413 & 3395193 & 13626571 & 54768439 & 220390761 & 887757717 & 3579054867\\
$Q$ & 29204 & 114666 & 452227 & 1789916 & 7105312 & 28274531 & 112746667 & 450377230 & 1801803339\\
$R$ & 15744 & 60366 & 233323 & 907529 & 3547802 & 13926810 & 54857224 & 216706457 & 858185665\\
$S$ & 22739 & 88352 & 345290 & 1355716 & 5343140 & 21124511 & 83738100 & 332687072 & 1324308001\\
$T$ & 21162 & 81966 & 319543 & 1252166 & 4927255 & 19455297 & 77039780 & 305805462 & 1216400095\\
$U$ & 40678 & 159874 & 631249 & 2501630 & 9943629 & 39622046 & 158205688 & 632793092 & 2534808294\\
$V$ & 11212 & 43866 & 172467 & 680772 & 2695907 & 10704713 & 42601646 & 169868638 & 678451341\\
$W$ & 9035 & 34875 & 135550 & 529733 & 2079362 & 8191718 & 32369106 & 128231694 & 509107371\\
$X$ & 9950 & 38100 & 147076 & 571394 & 2231348 & 8750584 & 34438134 & 135936824 & 537948754\\
$Y$ & 6709 & 25491 & 97773 & 377796 & 1468440 & 5735092 & 22488118 & 88474763 & 349078294\\
$Z$ & 4025 & 15214 & 58108 & 223734 & 866978 & 3377002 & 13210216 & 51861498 & 204221945\\
\bottomrule
\end{tabular}

\caption{Exact profiles for sizes $10$ through $18$.}
\end{table}
\clearpage

\section{A minimal exact certificate checker}
\label{kfp:app:checker}
The following core is sufficient to check an upper candidate once the map
and true prefix are supplied. The full verifier also checks the input
conventions, all intermediate certificates, the dual weights, and its
rational logarithm enclosure.
\begin{lstlisting}[language=Python]
from fractions import Fraction as Q

# P[i] and D[i] are integer coefficient lists, low degree first.
def evaluate(coefficients, zeta):
    value = Q(0)
    for coefficient in reversed(coefficients):
        value = value * zeta + coefficient
    return value

def check_upper(zeta, values, P, D, polynomial_map):
    prefix = [evaluate(p, zeta) for p in P]
    defect = [evaluate(d, zeta) for d in D]
    image = polynomial_map(zeta, values)
    tail = [v - p for v, p in zip(values, prefix)]
    residual = [v - f + d
                for v, f, d in zip(values, image, defect)]
    if any(t < 0 for t in tail):
        raise ValueError("Candidate lies below the known prefix")
    if any(r < 0 for r in residual):
        raise ValueError("A supersolution inequality failed")
    return tail, residual
\end{lstlisting}

The defect is computed, not guessed. If $P_i(\xi)=\sum_{n=1}^Na_{i,n}\xi^n$,
then for each row $i$ one expands the finitely many monomials in
$\Phi_i(\xi,P)$, discards degrees above $N$, and subtracts $P_i$.
Because all numbers involved are integers before evaluation, this step is
particularly suitable for independent implementations or proof-assistant
reflection.

For the dual checker, \eqref{kfp:eq:balance} is a sum of integer vectors.
The $j$th entry of a monomial exponent vector is counted with its full
multiplicity. Omitting the factor two for a square, for example, would
invalidate the certificate. The supplied verifier iterates over repeated
variable indices, so this multiplicity is explicit in its sparse format.

\section{Source audit and claim classification}
\label{kfp:app:audit}

\begin{longtable}{@{}p{0.25\textwidth}p{0.67\textwidth}@{}}
\toprule
Claim or ingredient & Status and evidence\\
\midrule
\endhead
Geometric inequalities & Imported from Bui v2, Lemma~4 and Appendix~B; corresponding
repository endpoint source inspected at the pinned commit.
\wtag\ In \repo\ a Lean theorem for the actual counts,
\lean{geometricPublishedBuiRecurrences} (\cref{kfp:sec:formal-status}).\\[0.4em]
Old $4.5235$ upper bound & Existing repository result; rational vector independently
replayed in the new verifier. No new Lean build was performed.
\wtag\ Kernel-checked in \repo; the project's Lean-verified bound.\\[0.4em]
Prefix-corrected majorant & Proved in \cref{kfp:sec:prefix}, including the nonnegative tail
map and the known-prefix lower-budget condition.\\[0.4em]
New $4.498$ bound & Computer-assisted theorem: exact marked enumeration plus a rational
certificate, with the detailed dependency chain in \cref{kfp:sec:verification}.\\[0.4em]
Marked table through $18$ & Generated by exhaustive C++ enumeration. Python agrees through
$11$; direct and incremental C++ scans agree through $16$; standard unmarked
totals agree through $18$. Not kernel-checked.
\wtag\ Regenerated through $18$ at intake from the final source, and
through $10$ by an independent enumerator
(\cref{kfp:app:provenance}).\\[0.4em]
\wtag\ New $4.498$ bound in Lean & Not formalized; the Lean endpoints are
specialized to $\zeta=2000/9047$ (\cref{kfp:sec:formal-status}).\\[0.4em]
Dual obstruction & Elementary weighted AM--GM proof; particular integer balance and
strict log inequality replayed exactly.\\[0.4em]
Certificate completeness & Proved for productive, strongly connected, nonlinear positive
polynomial systems with rational data. Not claimed for all maximum or signed
recurrences. Based on classical geometric-programming ideas.\\[0.4em]
Numerical critical values & Exploratory estimates only; not certified optimal values and
not used to establish the accepted inequalities.\\[0.4em]
Stability theorem and gap & Written proofs in \cref{kfp:sec:stability}. The case of limiting
spectral radius exactly one is not classified generally.
[Added 29 September 2026, batch 40: classified for irreducible $J$ in
Part~II, \cref{kfp:sat:thm:critical}; written proof, not formalized.]\\[0.4em]
Sensitivity formula & Conditional on an already justified differentiable critical
family. Not used for the headline bound.\\[0.4em]
Novelty & The explicit bound and application were not found in the inspected sources.
No claim of an exhaustive priority search or independent refereeing.\\
\bottomrule
\end{longtable}

\section{\wtag\ Provenance, intake verification and non-claims}
\label{kfp:app:provenance}

\subsection{Provenance}
This report was built from one manuscript: manuscript~06 of \repo's
incoming-reports batch~36, the archive \texttt{ProveIt\_Polyomino\_Research}
(delivered in commit \texttt{e13affd32}, placed in commit
\texttt{1a1396d4d}). It was written against the pin \sha; the project it
continues is unchanged between that pin and the placement commit, so its
statements about the project are current except where the delivered text
undersells the formal status of the geometric lemma
(\cref{kfp:sec:formal-status}). As a single source it required no merge
choices. Its whole text is printed; passages marked \wtag\ were added when it
was written into \repo:
\cref{kfp:sec:notation,kfp:sec:formal-status}, this appendix, the paragraph
on earlier upper bounds in \cref{kfp:sec:scope}, notes in the title-page box
and the ``three levels'' box, after \eqref{kfp:eq:partitions} and the
geometric-lemma paragraph of \cref{kfp:sec:system}, on intake regeneration in
\cref{kfp:sec:numeric}, on layout and on the Lean plan in
\cref{kfp:sec:verification}, after research question~1, three rows of the
audit table in \cref{kfp:app:audit}, and four bibliography entries
(\cite{repobound,eden,klarnerrivest,barequetshalah}). All labels carry the
prefix \texttt{kfp:} except \texttt{eq:map}, which the delivered, generated
\texttt{tex/dual\_table.tex} cites by that name.

[Added 29 September 2026, batch 40:] The report now has a second source.
Batch-40 manuscript~05 (archive \texttt{ProveIt\_Spectral\_Saturation},
arrived in \texttt{017e6c0d8}, placed in \texttt{afb2d1227}, pinned to
\texttt{74f7f5bdb}) is printed as Part~II, with its own provenance in
\cref{kfp:sat:sec:source}; its labels carry the prefix \texttt{kfp:sat:}.
The Part~I text above changed only by dated ``Added 29 September 2026''
pointers, the entries for Part~I and its appendices in the contents, the
label \texttt{kfp:sec:conclusion}, four macros for Part~II in the preamble,
and the PDF subject line.

\subsection{Intake verification}
\begin{itemize}
\item The arithmetic verifier \texttt{code/verify.py} passed on a copy of the
 package (4.3\,s) and printed exactly \texttt{data/verification.txt}; the
 profile checker accepted \texttt{data/profiles\_18.csv}.
\item The size-$18$ table was regenerated from scratch by a fresh build of
 \texttt{code/enumerate.cpp} and matched in every entry
 (\cref{kfp:sec:numeric}); fresh runs through sizes $11$--$14$, $16$ and $17$
 and a direct-scan run through $12$ matched too.
\item An independent set-based Python enumerator, with coordinates from the
 project's \texttt{Patterns.lean}, matched all entries through size $10$.
 An independent transcription of \eqref{eq:map} from the project's research
 note \cite{repobound} reproduced \eqref{kfp:eq:first-defect}, the size-$18$
 margins and $v_G$, and confirmed the old $4.5235$ vector as a supersolution.
\item The seventeen offset sets of \cref{kfp:tab:patterns} were compared
 entry by entry with \texttt{Patterns.lean}, and the unmarked column with OEIS
 A001168 through $18$.
\item A $50$-digit evaluation of $\log\mathcal C(\zeta_{\rm dual},m)$ gave
 $1.00066802314520877\ldots$, inside \eqref{kfp:eq:dual-interval}.
\end{itemize}
These checks do not verify the written proofs, and no independent proof review
was performed. At intake the proofs of \cref{kfp:lem:formal},
\cref{kfp:thm:majorant,kfp:thm:upper,kfp:thm:nested,kfp:thm:dual,kfp:thm:alternative,kfp:thm:gap}
were checked and that of \cref{kfp:thm:stability} was read; no gap was found.

\subsection{What this report does not claim}
Every limitation stated by the delivered text is kept; collected:
\begin{enumerate}[label=N\arabic*.,leftmargin=2.6em]
\item The lower bound $\mu_{\rm original}\geq452349/100000$ concerns the
 uncorrected equality recurrence and is \emph{not} a lower bound on
 $\lambda$ (\cref{kfp:sec:dual}).
\item No new Lean formalization. The Lean proof of $4.5235$ in \repo\ is not a
 proof of $4.498$, of the finite-prefix theorem, or of the size-$18$ data.
\item The arithmetic checker does not prove that the table enumerates the
 stated objects; its partition, defect and $G$-domination checks are
 necessary conditions only.
\item The enumeration cross-checks are not a formal certificate that every
 size-$18$ object was generated exactly once. The delivered record made no
 fresh size-$18$ run after the final command-line changes; the intake
 regeneration closes that gap only with the same program, and the two C++
 implementations share one generation strategy.
\item The report is an AI-assisted draft and has not been independently
 refereed.
\item Novelty is assessed only against the inspected repository snapshot and
 sources; there was no exhaustive priority search.
\item Geometric programming, weighted AM--GM and Redelmeier-style enumeration
 are prior tools, not claimed inventions.
\item The certificate alternative is proved only for productive, strongly
 connected, nonlinear systems with rational data; it gives no uniform bound
 on certificate bit length and does not cover maximum or signed recurrences.
\item The case $\spr(J)=1$ of \cref{kfp:thm:stability} is not classified.
 [Added 29 September 2026, batch 40: re-scoped. Part~II classifies it for
 irreducible $J$ (\cref{kfp:sat:thm:critical}); the reducible case is still
 not classified.]
\item The counterexample of \cref{kfp:thm:gap} does not disprove Bui's
 Question~1.
\item The sensitivity formula \eqref{kfp:eq:sensitivity} is conditional and
 not used in \cref{kfp:thm:main}.
\item The numerical critical estimates in \cref{kfp:tab:progress} are not
 certified values or bounds.
\item The OEIS comparison is a consistency check, not an input to the proof or
 the checker.
\end{enumerate}
Added in the write: (W1) the geometric lemma is kernel-checked in \repo, but
only relative to the transcription of Bui's diagrams into
\texttt{Patterns.lean}; (W2) the bound $4.498$ remains unverified by Lean,
and the project's Lean-verified bound remains $9047/2000$; (W3) the earlier
bounds of \cref{kfp:sec:scope} were not re-examined, and their sources were
not read beyond their bibliographic data.

%% ====================================================================
%% Batch-40 addition (manuscript 05, "When Exact Prefixes Stop
%% Improving"). All labels of Part II carry the prefix kfp:sat:
%% (spectral saturation).
%% ====================================================================
\clearpage
% Part II follows Part I's appendices so that no Part I number changes; its
% sections continue Part I's arabic numbering.
\setcounter{section}{11}
\renewcommand{\thesection}{\arabic{section}}
\renewcommand{\theHsection}{partII.\arabic{section}}
\phantomsection\part*{Part II.\quad Spectral saturation of the exact-prefix hierarchy}
\addcontentsline{toc}{part}{Part II.\texorpdfstring{\quad}{ } Spectral saturation of the exact-prefix hierarchy}

\noindent\emph{Sections~\ref{kfp:sec:scope}--\ref{kfp:sec:conclusion} and
Appendices~\ref{kfp:app:data}--\ref{kfp:app:provenance} form the original
report (Part~I), written into \repo\ in batch~36, unchanged except for dated
pointers to this addition. Part~II was added on 29 September 2026 in
batch~40 of \repo's incoming-report intake, from one manuscript written
against the state of Part~I printed here. It classifies the equality case
$\spr(J)=1$ that Part~I left unclassified after \cref{kfp:thm:stability}
(limitation N9 of \cref{kfp:app:provenance}), \emph{for an irreducible
true-profile Jacobian only}, determines the limiting radius of the whole
prefix hierarchy, and proves that no exact prefix can take the fixed Bui
map's majorant growth below $4.3149$. \Cref{kfp:sat:sec:source} records
the provenance and fixes the notation of the addition.}

\section{Provenance, conventions and claim boundary of Part~II}
\label{kfp:sat:sec:source}

\subsection{One manuscript, one addition}
Part~II is batch-40 manuscript~05, \emph{When Exact Prefixes Stop
Improving: Spectral saturation, sharp square-root laws, and a certified
barrier for convolution majorants}, an AI-assisted mathematical draft
prepared for Vladimir Reshetnikov and dated 28 September 2026 (27 pages as
delivered). It arrived as the archive
\texttt{ProveIt\_Spectral\_Saturation.zip} (inner directory
\texttt{ProveIt\_Spectral\_Saturation/}) in \repo\ commit
\texttt{017e6c0d8} and was placed in \texttt{afb2d1227}. It was written
against \repo\ commit \texttt{74f7f5bdb}, at which Part~I's
\texttt{article.tex} had the Git blob \texttt{4a021fb3}. That blob is
unchanged at the placement commit, and so are the project
\texttt{Combinatorics/Polyominoes/KlarnerConstant} and Part~I's
\texttt{code/model.py} and \texttt{data/profiles\_18.csv}; every statement
the manuscript makes about Part~I therefore refers to the text printed above.

Its code, data and audit files are shipped beside Part~I's with the prefix
\path{02-spectral-saturation-} (Part~I's unprefixed files count as~01):
three Python scripts and a build script in \texttt{code/}, five data files
in \texttt{data/}, and
\path{02-spectral-saturation-PROOF_STATUS.md} and
\path{02-spectral-saturation-SOURCES.md} at the root, all byte-identical to
the delivery. Its manuscript, README and PDF are not shipped and survive in
the arrival commit. Its \texttt{data/profiles\_18.csv} is not shipped
either: the manuscript copied it from Part~I, and it is byte-identical to
Part~I's \texttt{data/profiles\_18.csv} (blob \texttt{82648328}), which is
the file every mention of it in Part~II now means.

The manuscript contributes, in the order printed here: its own
introduction and claim boundary (\cref{kfp:sat:sec:intro}); a restatement
of the prefix construction (\cref{kfp:sat:sec:setup}); the critical
alternative, which classifies the irreducible case $\spr(J)\ge1$
(\cref{kfp:sat:sec:critical}); the exact limiting radius of the hierarchy
(\cref{kfp:sat:sec:saturation}); a tail estimate without coefficient
regularity (\cref{kfp:sat:sec:tails}); the square-root law at an interior
nonlinear crossing (\cref{kfp:sat:sec:sharp}); the scaled branch and the
quarter-power amplitude (\cref{kfp:sat:sec:branch}); sharp equivalents for
Part~I's scalar example (\cref{kfp:sat:sec:scalar}); four examples that
delimit the hypotheses (\cref{kfp:sat:sec:exceptions}); the all-prefix
$4.3149$ barrier for the fixed Bui map (\cref{kfp:sat:sec:bui}); its
reproducibility and proof audit (\cref{kfp:sat:sec:audit}); ten research
projects (\cref{kfp:sat:sec:future}); its conclusion; and, as the last three
sections of Part~II, its three appendices (the map, a replay algorithm, and
the source ledger). Every result, proof, example, remark, limitation and
question of the manuscript is printed. Its bibliography is merged into the
report's; the entries it shares with Part~I (Bui's paper) are printed once.

\paragraph{Material already in Part~I.} The manuscript restates, to be
self-contained, Part~I's formal comparison lemma, prefix-corrected majorant,
certificate theorem, monotonicity, strict subcritical criterion and scalar
limit~$3$, and says so itself. These restatements are kept where Part~II's
proofs cite them in its own formulation, each followed by a \wtag\ note
naming the Part~I result; where the manuscript's proof differs from Part~I's
it is kept as a second proof. The manuscript's display of the $53$-monomial
map duplicates \eqref{eq:map} and is printed once, there
(\cref{kfp:sat:app:map}).

\paragraph{The manuscript's abstract.}
\begin{quote}\small
Exact initial coefficients can strengthen a positive convolution recurrence
without changing its combinatorial proof. They need not, however, make its
majorant asymptotically exact. We study this limitation for a proper
nonnegative polynomial system $A\preceq\Phi(\xi,A)$, its defect
$D=\Phi(\xi,A)-A$, and the corrected equations
$B_N=\Phi(\xi,B_N)-[\xi^{\le N}]D$.

First, at any convergent evaluation where the true-profile Jacobian is
irreducible and has spectral radius at least one, a level-$N$ upper
certificate exists \emph{if and only if} the defect tail vanishes. This
settles the irreducible critical case left open in the repository's
finite-prefix report. It also identifies the exact limiting radius of the
hierarchy: the first Perron crossing, or the true radius if no earlier
crossing occurs.

Second, at an interior nonlinear crossing $r_*$, put
$\eta_N=\ell^T(D-[\xi^{\le N}]D)(r_*)$. With explicitly defined positive
constants $\kappa$ and $\beta$, we prove
\[
 r_*-\rho_N\sim\frac{2\sqrt\beta}{\kappa}\sqrt{\eta_N},
 \qquad
 B_N(\rho_N)-A(\rho_N)
 \sim u\sqrt{\eta_N/\beta}.
\]
A universal scaled branch and a quarter-power law for its local singular
amplitude follow. No regular variation or nonlacunarity assumption on the
defect coefficients is needed.

Finally, using the repository's size-$18$ marked polyomino profile, an
exact rational vector proves that every member of its fixed-map prefix
hierarchy has growth at least $4.3149$. This is a lower bound on the
\emph{majorants}, not on polyomino growth. The package includes a
standard-library arithmetic checker, exact scalar root brackets, source
provenance, and a research agenda. The proofs are unrefereed and not
Lean-verified; the size-$18$ enumeration is inherited rather than rerun.
\end{quote}
(``The repository's finite-prefix report'' is Part~I. The size-$18$ table
was not rerun for the manuscript, but Part~I's intake regenerated it from
the final enumerator source; see \cref{kfp:app:provenance}.)

\begin{resultbox}
\textbf{Status of Part~II.} Its analytic theorems have written proofs; it is
an AI-assisted, unrefereed draft, and independent review of the uniform fold
argument (\cref{kfp:sat:lem:localization} and \cref{kfp:sat:sec:sharp}) is
recommended before correctness or novelty is treated as settled. Novelty is
a bounded assessment against the sources inspected, not a priority claim.
\textbf{Nothing in Part~II is formalized} in any proof assistant; the
project's Lean-verified bound remains $\lambda\le9047/2000=4.5235$
(\cref{kfp:sec:formal-status}). Its arithmetic checker proves finitely many
rational inequalities, not the analytic theorems; its size-$18$ input is
Part~I's enumeration, which is not kernel-checked.

\smallskip
\textbf{What $4.3149$ is.} \Cref{kfp:sat:thm:bui} proves that every
exact-prefix majorant of the fixed seventeen-variable Bui map has growth
$\mu_N\ge4.3149$: a limit of this \emph{method}. \emph{Tempting false
reading:} ``$\lambda\ge4.3149$'', a lower bound on Klarner's constant. It is
not; the true counts lie \emph{below} the majorants. \emph{Tempting false
reading:} ``$4.3149$ improves on $4.498$''. It does not; it is not an upper
bound at all, but a floor under every upper bound this method can certify.
Part~I's $4.498$ (computer-assisted) and the Lean-verified $4.5235$ stand
unchanged above it.

\smallskip
\textbf{What N9 now is.} The equality case $\spr(J)=1$ is classified for
\emph{irreducible} $J$ (\cref{kfp:sat:thm:critical}); the reducible critical
case, the crossing at the true singularity, and the exact fixed-map
saturation constant remain open (\cref{kfp:sat:sec:exceptions,kfp:sat:sec:future}).
\end{resultbox}

\subsection{Notation of Part~II}
\label{kfp:sat:sec:notation}
Part~II keeps the manuscript's notation with two changes, made so that it
reads consistently with Part~I:
\begin{itemize}
\item the manuscript's formal size variable $z$ (its Sections~1--8) is
written $\xi$, as in Part~I and as the manuscript itself does from its
Section~9 on, where $z$ is a coordinate of \eqref{eq:map};
\item the manuscript's Perron eigenvalue $\lambda$, $\lambda(t)=\spr(J(t))$
(proof of \cref{kfp:sat:thm:critical}, \cref{kfp:sat:lem:positive-constants}
with \eqref{kfp:sat:eq:lambda-expansion}, \cref{kfp:sat:lem:localization}),
is written $\Lambda$, because $\lambda$ is Klarner's constant in Part~I.
\end{itemize}
No normalization changed. Macro-level spellings (the manuscript's
$\mathcal R_N$, $[z^{\le N}]$, $[z^{>N}]$, $\mathbf 1$) are typeset as
$\satR_N$, $\strunc N$, $\shigh N$, $\bone$. The delivered, generated table
\path{data/02-spectral-saturation-scalar_table.tex} (\cref{kfp:sat:tab:scalar})
heads its gap column $\rho_*-\rho_N$; that $\rho_*$ is the crossing $r_*$
of the text, here $1/3$.

\begin{longtable}{@{}>{\raggedright\arraybackslash}p{0.2\textwidth}>{\raggedright\arraybackslash}p{0.38\textwidth}>{\raggedright\arraybackslash}p{0.36\textwidth}@{}}
\caption{Symbols of Part~II against Part~I.}\label{kfp:sat:tab:notation}\\
\toprule
Symbol & Meaning in Part~II & In Part~I\\
\midrule
\endfirsthead
\toprule
Symbol & Meaning in Part~II & In Part~I\\
\midrule
\endhead
$\xi$ & formal size variable (the manuscript's $z$ in its Sections 1--8) & the same\\
$A$, $A_i$ & the true series (a general subsolution $A\preceq\Phi(\xi,A)$) & $\bA$; $A_n$ is the polyomino count\\
$t$, $t_0$ & a positive evaluation point; $t_0=10000/43149$ & $\zeta$; $\zeta=500/2249$ for $4.498$, $\zeta_{\rm dual}$\\
$a$ & $A(t)$; also the constants of \cref{kfp:sat:lem:tail-flatness} ($0<a<r<b$) and of the fold argument, and $a_N$ the singular amplitude & $a=\bA(\zeta)$ in \cref{kfp:sec:stability}; $a_{i,n}$ counts\\
$D$, $D_N$ & full defect $\Phi(\xi,A)-A$; $D_N=\strunc N D$ & $D_N$ the same; $D$ also the defect polynomial of \eqref{kfp:eq:sensitivity}; $D$ a coordinate\\
$\satR_N$ & the defect tail $D-D_N$ & $\tail{N}{\cdot}$ is the operator $[\xi^{>N}]$\\
$b_N(t)$ & $(\shigh N\Phi(\xi,P_N))(t)$, finite forcing & the same $b_N$ (\cref{kfp:sec:stability}); not the project's $b_N$\\
$P_N$, $Y_N$, $B_N$, $\Psi_N$ & prefix, tail, majorant, tail map & the same\\
$J$, $J(t)$, $J_*$ & $\Phi_X$ at the \emph{true} profile $A(t)$; $J_*$ at $r_*$ & $J$ the same; $J_N$ at the prefix\\
$\Lambda$, $\Lambda(t)$ & Perron value of $J$, $J(t)$ (the manuscript's $\lambda$) & $\lambda$ is Klarner's constant\\
$R$ & system radius of $A$ & required offset set; radii $R_i$\\
$\rho_N$, $\mu_N$ & system radius of $B_N$, $\mu_N=1/\rho_N$ & $\rho_N$ as in \cref{kfp:thm:gap}; $\mu_{\rm original}$ is the growth of the uncorrected recurrence\\
$r_*$ & first Perron crossing \eqref{kfp:sat:eq:barrier} & ---\\
$\ell$, $u$ & left and right Perron vectors, $\ell^Tu=1$ & $\ell$ as in \eqref{kfp:eq:sensitivity}; $u$ a coordinate\\
$\kappa$, $\beta$, $\eta_N$, $\varepsilon_N$ & invariants \eqref{kfp:sat:eq:constants}; $\varepsilon_N=\sqrt{\eta_N}$ (also $\sqrt{H_N(r)}$ in \cref{kfp:sat:lem:tail-flatness}) & $\varepsilon$ a small parameter\\
$h$ & $v-a$ or $B_N(t)-A(t)$ & $h=(I-J)^{-1}\bone$\\
$Q_t(h)$ & nonlinear remainder \eqref{kfp:sat:eq:key} & $Q$ a coordinate\\
$G(t,h,r)$ & the map \eqref{kfp:sat:eq:G} & $G$ the coordinate whose counts dominate $A_n$\\
$F_N$, $F$ & scaled reduced function \eqref{kfp:sat:eq:scaled-reduction}; $F$ also a generic vector series & $F(X)=\Phi(\zeta,X)$; $F$ a coordinate\\
$K$, $K(t)$ & scaled distance $\delta/\varepsilon_N$ and a bound in \cref{kfp:sat:lem:tail-flatness}; $K(t)=\sum_{k<d}J(t)^k$ & ---\\
$S$, $S_c$, $S_-$ & a subsolution (\cref{kfp:sat:lem:graded}); scaled height & sets $S$, $S_K$; $S$ a coordinate\\
$H$, $H_N$ & a scalar series and its tail (\cref{kfp:sat:lem:tail-flatness}); $H_N(t)$ also the scalar sign oracle (\cref{kfp:sat:app:algorithm}) & $H$ the nonlinear part of $\Phi$\\
$T(\xi,Y)$ & a generic proper map & $T_N=\strunc N\Phi(\xi,P_N)$\\
$w$ & witness vector (\cref{kfp:sat:tab:witness}); direction in \cref{kfp:sat:prop:quadratic-test}; complement component $h-su$ & coordinate $w$; $w=v-P_N(\zeta)$\\
$r$, $r_N$ & a point, a vector parameter in \eqref{kfp:sat:eq:G}, an index; $r_N=\|\satR_N(t)\|_\infty$ & $r$ a coordinate; row sums $r_i$\\
$\gamma$ & $1000001/1000000$ & ---\\
$e_M$ & the $e$-coefficient of size $M$ & $a_{E,M}$\\
\bottomrule
\end{longtable}

\section{The question, the answer, and the claim boundary}
\label{kfp:sat:sec:intro}

\subsection{Why a better prefix may still be the wrong limit}
Suppose a combinatorial counting problem gives a vector of generating
functions $A(\xi)$ and a componentwise upper recurrence
\[
 A(\xi)\preceq\Phi(\xi,A(\xi)),
\]
where $\Phi$ is a polynomial with nonnegative coefficients. The equality
solution is an upper bound, but its recurrence can overcount. A natural
repair is to compute exact coefficients through size $N$ and subtract the
corresponding overcounts. The resulting majorant agrees with the true
series through size $N$ and improves as $N$ increases.

There are two very different possible conclusions. One is that every
fixed coefficient eventually becomes correct. The other is that the
exponential growth becomes correct. The first follows directly from the
construction. The second can fail because small residual errors are
amplified by an unstable feedback loop in the recurrence. The Jacobian of
that loop, evaluated at the \emph{true} generating function, is the
quantity that decides the limiting obstruction.

This article develops that observation into a precise theorem, a sharp
asymptotic law, and a small exact certificate applicable to a concrete
polyomino recurrence. The contribution is not a further decimal improvement
to the existing upper bound; it explains both the power and the ultimate
limitation of continuing that particular improvement strategy.

\subsection{The repository problem being continued}
The inspected \repo\ snapshot contains the report \emph{Finite-Prefix
Corrections and Dual Certificates for Polyomino Growth} \cite{satprefix}.
It builds on the formal project
\texttt{Combinatorics/Polyominoes/}\allowbreak\texttt{KlarnerConstant}. The report's Section~8
proves eventual certificate existence when
\[
 \spr\bigl(\Phi_X(t,A(t))\bigr)<1,
\]
and an obstruction under strict supercriticality and additional finite-prefix
conditions. It explicitly leaves the equality case unclassified. It also
shows, in a scalar example, that prefix-majorant growth can tend to $3$
while true growth is $1$.

The same report lists locating the actual fixed-map stability barrier and
obtaining certified critical asymptotics among its proposed research
projects. These are the specific continuation points addressed here.
The numerical bounds reported by that source have distinct statuses:
$4.5235$ is the endpoint of its cited Lean development, whereas $4.498$ is
the source report's computer-assisted, nonformalized upper bound. We do not
reprove either endpoint in this article \cite{satprefix,satformal}.

\wtag\ The report meant here is Part~I above, read at the pin
\texttt{74f7f5bdb}, where its \texttt{article.tex} is the blob now printed:
its Section~8 is \cref{kfp:sec:stability}, the equality case is the sentence
after the proof of \cref{kfp:thm:stability}, the scalar example is
\cref{kfp:thm:gap}, and the two research projects are questions~4 and~9 of
\cref{kfp:sec:questions}. The Lean-verified bound is $\lambda\le9047/2000$
(\cref{kfp:sec:formal-status}); nothing in Part~II changes it.

\subsection{Results proved here}
The main results are as follows. The theorem numbers give the precise
hypotheses; in particular, irreducibility and nonlinearity must not be
silently discarded.

\begin{center}\small
\begin{tabular}{@{}p{.20\linewidth}p{.48\linewidth}p{.22\linewidth}@{}}
\toprule
Question & Answer & Location\\
\midrule
Critical certificates & At an irreducible Perron value $\ge1$, a certificate
exists exactly when the omitted defect is zero. & Theorem~\ref{kfp:sat:thm:critical}\\[4pt]
Limiting accuracy & The prefix radii tend to the first spectral crossing,
or the true radius when there is no earlier crossing. & Theorem~\ref{kfp:sat:thm:limit}\\[4pt]
Rate of approach & At an interior nonlinear crossing, the radius gap is
asymptotic to an explicit constant times the square root of the omitted
defect. & Theorem~\ref{kfp:sat:thm:sharp}\\[4pt]
Local branch shape & The scaled solution has a quadratic universal profile;
its square-root singular amplitude has a quarter-power law. &
Theorem~\ref{kfp:sat:thm:branch}\\[4pt]
Concrete limitation & The fixed $17$-variable polyomino map cannot yield
prefix-majorant growth below $4.3149$, whatever the prefix length. &
Theorem~\ref{kfp:sat:thm:bui}\\
\bottomrule
\end{tabular}
\end{center}

\subsection{Relation to prior methods and novelty}
Bui's convolutional approach supplies the polyomino recurrence and a
finite-parameter method for certifying its growth \cite{bui}. We checked
version~2, revised May~6, 2026, rather than relying on the differently
titled first version. The exact-prefix construction, its monotonicity, the
strict stability criterion, and the scalar limit $3$ are already in the
repository report \cite{satprefix}. They are restated and proved below to
make the extension self-contained, not presented as discoveries here.

Inductive upper bounds and Perron-directed certificate searches are also
established techniques for positive polynomial systems. In particular,
Winkler and Katoen characterize strict inductive upper bounds using the
Jacobian at a least fixed point \cite{satwk}; Newton methods for monotone
polynomial systems have a substantial independent literature \cite{satsey}.
Square-root singularities of positive functional systems are classical
\cite{satdrmota}. Neither Perron--Frobenius theory nor the local fold mechanism
is claimed as new.

The proposed contribution is their precise application to \emph{a moving
family of exact-prefix corrections}: the full-defect critical alternative,
the exact spectral saturation theorem, the sharp normalization in terms of
the omitted defect without a regularity assumption on its coefficient
sequence, the accompanying amplitude law, and the explicit all-prefix
polyomino obstruction. We did not find these statements in the inspected
repository continuation or the primary sources reviewed. This is a bounded
novelty assessment, not an exhaustive priority claim.

\subsection{What is not being claimed}
This article does not determine Klarner's constant, improve the source's
$4.498$ upper bound, establish a lower bound $4.3149$ for the true constant,
or resolve Bui's separate question about enlarging the geometric
neighborhood system. It does not certify the repository's entire
mathematical corpus. Its analytic results have ordinary written proofs;
its numerical application has an exact arithmetic checker, but the marked
enumeration and the new theorems are not kernel-checked here.

\section{Positive systems and exact-prefix corrections}
\label{kfp:sat:sec:setup}

\subsection{Standing definitions}
For vectors, inequalities are componentwise. For power series,
$F\preceq G$ means that every coefficient of every component of $G-F$ is
nonnegative. A series has zero constant term unless stated otherwise.

Fix a dimension $d\ge1$ and
\[
 \Phi(\xi,X)\in\R_{\ge0}[\xi,X_1,\ldots,X_d]^d,
 \qquad \Phi(0,0)=0.
\]
We call the system \emph{proper} when
\begin{equation}\label{kfp:sat:eq:proper}
 L=\Phi_X(0,0)\quad\text{is nilpotent}.
\end{equation}
This condition allows acyclic zero-size linear dependencies; it is weaker
than requiring every monomial to contain $\xi$. Throughout, let
\begin{equation}\label{kfp:sat:eq:subsolution}
 A\in \xi\R_{\ge0}[[\xi]]^d,
 \qquad A\preceq\Phi(\xi,A).
\end{equation}
The \emph{full defect} and its prefix and tail are
\begin{equation}\label{kfp:sat:eq:defects}
 D=\Phi(\xi,A)-A\succeq0,\qquad
 D_N=\strunc N D,\qquad
 \satR_N=D-D_N.
\end{equation}
The exact profile prefix is $P_N=\strunc N A$.
Since all substituted series have zero constant term,
\begin{equation}\label{kfp:sat:eq:finite-computation}
 D_N=\strunc N\Phi(\xi,P_N)-P_N.
\end{equation}
Thus $D_N$ can be computed from the prefix alone; no infinite-series
oracle is required to construct a finite correction.

Define the positive tail map
\begin{align}
 \Psi_N(\xi,Y)
 &=\Phi(\xi,P_N+Y)-\strunc N\Phi(\xi,P_N)\label{kfp:sat:eq:psi}\\
 &=\Phi(\xi,P_N+Y)-P_N-D_N.\nonumber
\end{align}
Although the second line uses subtraction, the polynomial $\Psi_N$ has
nonnegative coefficients. Indeed, its terms involving $Y$ are positive
terms of the polynomial expansion, and its constant-in-$Y$ part is
$\shigh N\Phi(\xi,P_N)$.

\subsection{Formal fixed points and comparison}
\begin{lemma}[Graded convergence]\label[lemma]{kfp:sat:lem:graded}
A proper nonnegative polynomial map $T(\xi,Y)$ with $T(0,0)=0$ has a unique
zero-constant formal fixed point $Y$. Iteration from any zero-constant
formal vector converges coefficient by coefficient to $Y$, with every
coefficient eventually constant. Iteration from zero is nondecreasing.
If $S\succeq0$ has zero constant term and $S\preceq T(\xi,S)$, then
$S\preceq Y$. If $S\succeq T(\xi,S)$, then $S\succeq Y$.
\end{lemma}
\begin{proof}
Suppose two zero-constant vectors agree below degree $m$. The degree-$m$
part of the difference of their images under $T$ is $L$ times the
corresponding difference, where $L=T_Y(0,0)$. All other contributions
contain either a positive power of $\xi$ or another positive-degree series
factor, so they only involve lower-degree differences. If $L^q=0$,
$q$ iterations remove the degree-$m$ difference once lower degrees agree.
Induction on $m$ gives eventual agreement, uniqueness, and coefficientwise
stabilization of iteration. The stabilized coefficients solve the equation.

Nonnegative coefficients make $T$ monotone. Thus its iterates from zero
increase. For a subsolution, $S\preceq T^k(S)$ for every $k$, so passage
to the stabilized coefficients gives $S\preceq Y$. The reversed argument
applies to a supersolution. This last argument is important: a formal
subsolution need not lie below a particular \emph{finite} iterate from zero.
\end{proof}

\noindent\wtag\ \emph{Restatement.} Existence, uniqueness, stabilization and
subsolution comparison are Part~I's \cref{kfp:lem:formal}; the convergence
from an arbitrary zero-constant start and the supersolution comparison are
the manuscript's additions. The proof above, by degree-wise contraction
rather than by inverting $I-L$, is kept as a second proof.

\begin{proposition}[Prefix majorants; prior construction]\label[proposition]{kfp:sat:prop:prefix}
Let $Y_N$ be the formal fixed point of $\Psi_N$ and put $B_N=P_N+Y_N$.
Then
\begin{equation}\label{kfp:sat:eq:corrected}
 B_N=\Phi(\xi,B_N)-D_N,
 \qquad
 A\preceq B_{N+1}\preceq B_N,
 \qquad \strunc N B_N=P_N.
\end{equation}
In particular, $B_N\to A$ coefficientwise.
\end{proposition}
\begin{proof}
The linear part of $\Psi_N$ at $(0,0)$ is $L$, so
Lemma~\ref{kfp:sat:lem:graded} applies. Its forcing begins above degree $N$;
iteration from zero therefore shows that $Y_N$ begins above degree $N$.
Equation~\eqref{kfp:sat:eq:corrected} follows from the definition.

The true tail $S=A-P_N$ satisfies
\[
 \Psi_N(\xi,S)=S+\satR_N\succeq S,
\]
so $A\preceq B_N$. Hence $B_N-P_{N+1}\succeq0$, and
\[
 \Psi_{N+1}(\xi,B_N-P_{N+1})
 =B_N-P_{N+1}-(D_{N+1}-D_N)
 \preceq B_N-P_{N+1}.
\]
Formal supersolution comparison gives $B_{N+1}\preceq B_N$.
The prefix identity gives coefficientwise convergence.
\end{proof}

\noindent\wtag\ \emph{Restatement.} This is Part~I's
\cref{kfp:thm:majorant,kfp:thm:nested}, as the manuscript's title
``prior construction'' says; the monotonicity proof above, by supersolution
comparison instead of degree induction, is kept as a second proof.

\subsection{Numeric certificates and the relevant radius}
For a nonnegative vector series $F$, define its \emph{system radius} as the
minimum of its component radii of convergence. This convention is useful
without first proving that all radii coincide. Let $R$ be the system radius
of $A$ and $\rho_N$ that of $B_N$, and write
\[
 \mu_N=1/\rho_N
 =\max_i\limsup_{n\to\infty}\bigl([\xi^n]B_{N,i}\bigr)^{1/n}.
\]
The conventions $1/\infty=0$ and $1/0=\infty$ apply.

\begin{definition}[Upper certificate]
A real upper certificate at a positive evaluation $t$ and level $N$ is a
finite vector $v$ satisfying
\begin{equation}\label{kfp:sat:eq:certificate}
 v\ge P_N(t),\qquad v\ge\Phi(t,v)-D_N(t).
\end{equation}
The tail condition $v\ge P_N(t)$ is part of the definition.
\end{definition}

\begin{lemma}[Exact numeric meaning]\label[lemma]{kfp:sat:lem:numeric}
A certificate at $(N,t)$ exists if and only if $B_N(t)$ is finite
coordinatewise. Every such certificate dominates $B_N(t)$ and hence
$A(t)$ whenever the latter is finite. When $B_N(t)$ is finite, it is itself
a certificate. Existence at level $N$ implies existence at every later
level, at the same evaluation point.
\end{lemma}
\begin{proof}
Write $y=v-P_N(t)\ge0$. Condition~\eqref{kfp:sat:eq:certificate} is
$\Psi_N(t,y)\le y$. Numeric iteration from zero stays below $y$ and
increases. Evaluation commutes with the nonnegative coefficientwise limit:
this follows by interchanging two increasing suprema, over iteration number
and coefficient cutoff. Therefore the limit is the finite sum $Y_N(t)$.
Polynomial continuity gives its fixed-point identity. Conversely a finite
sum $Y_N(t)$ satisfies that identity by evaluation, and so supplies a
certificate. The domination of $A$ and persistence at later levels follow
from Proposition~\ref{kfp:sat:prop:prefix}.
\end{proof}

\noindent\wtag\ \emph{Restatement.} The ``if'' direction with its
domination and the persistence are Part~I's
\cref{kfp:thm:upper,kfp:thm:nested} (whose certificate conditions
\eqref{kfp:eq:upper-conditions} are \eqref{kfp:sat:eq:certificate}, with
$\zeta$ written $t$); the converse, that a finite $B_N(t)$ is itself a
certificate, is not stated in Part~I.

A certificate consequently yields
$[\xi^n]B_{N,i}\le v_i t^{-n}$. Failure of a certificate at $t$ implies
$\rho_N\le t$, but does not by itself imply the strict inequality
$\rho_N<t$: a series can diverge exactly at its radius. This distinction
will matter in the linear example of Section~\ref{kfp:sat:sec:exceptions}.

\section{The critical spectral case: an exact alternative}
\label{kfp:sat:sec:critical}

Fix $t>0$ for which $A(t)$ is finite, and abbreviate
\[
 a=A(t),\qquad J=\Phi_X(t,a).
\]
A nonnegative matrix is called irreducible when its directed dependency
graph is strongly connected. We use the Perron--Frobenius theorem in its
standard finite-dimensional form: an irreducible nonnegative matrix has
strictly positive left and right eigenvectors for its spectral radius,
and that eigenvalue is algebraically simple \cite{satwk}.

\begin{lemma}[The full-defect inequality]\label[lemma]{kfp:sat:lem:full-defect}
Every certificate $v$ gives $h=v-a\ge0$ satisfying
\begin{equation}\label{kfp:sat:eq:key}
 h\ge Jh+Q_t(h)+\satR_N(t),
 \qquad
 Q_t(h)=\Phi(t,a+h)-\Phi(t,a)-Jh\ge0.
\end{equation}
If $v=B_N(t)$, equality holds.
\end{lemma}
\begin{proof}
Lemma~\ref{kfp:sat:lem:numeric} gives $v\ge a$. Also
$\Phi(t,a)=a+D(t)$. Subtracting $a$ from the certificate inequality and
expanding at $a$ gives~\eqref{kfp:sat:eq:key}. The remainder is nonnegative
because the polynomial has nonnegative coefficients and $a,h\ge0$.
\end{proof}

\begin{theorem}[Critical and supercritical alternative]\label{kfp:sat:thm:critical}
Suppose $J$ is irreducible and $\spr(J)\ge1$. For each $N$, the following
are equivalent:
\begin{enumerate}[label=\textup{(\roman*)},itemsep=2pt]
\item a real upper certificate exists at $(N,t)$;
\item $\satR_N(t)=0$;
\item $D=D_N$ as formal series;
\item $B_N=A$ as formal series.
\end{enumerate}
Thus an infinite full defect rules out every finite-prefix certificate at
an irreducible critical evaluation, not merely at a strictly supercritical
one.
\end{theorem}
\begin{proof}
Let $\ell>0$ be a left Perron vector, with eigenvalue $\Lambda\ge1$.
Multiplying~\eqref{kfp:sat:eq:key} by $\ell^T$ gives
\begin{equation}\label{kfp:sat:eq:perron-contradiction}
 (1-\Lambda)\ell^Th
 \ge\ell^TQ_t(h)+\ell^T\satR_N(t).
\end{equation}
The left side is nonpositive and both terms on the right are nonnegative.
Because $\ell>0$, a nonzero defect tail makes the right side strictly
positive. Hence (i) implies (ii).

All coefficients of $\satR_N$ are nonnegative and $t>0$, so its value
vanishes exactly when every coefficient vanishes. This proves (ii)$\Leftrightarrow$(iii).
When $D=D_N$, the series $A-P_N$ is a fixed point of $\Psi_N$.
Formal uniqueness gives $B_N=A$, proving (iii)$\Rightarrow$(iv).
Conversely, the equations for $A$ and $B_N$ give $D=D_N$ if $A=B_N$.
Finally, (iv) and finiteness of $A(t)$ give (i) by
Lemma~\ref{kfp:sat:lem:numeric}.
\end{proof}

\begin{corollary}[Uniqueness at a nonlinear critical point]\label[corollary]{kfp:sat:cor:unique}
Under Theorem~\ref{kfp:sat:thm:critical}, assume also $a>0$ and that $\Phi(t,X)$
contains a monomial of total $X$-degree at least two. Whenever a certificate
exists, the only one is $v=a$. In particular a rational certificate at this
critical or supercritical evaluation exists exactly when $D=D_N$ and
$a\in\Q^d$.
\end{corollary}
\begin{proof}
If $h\ne0$ in~\eqref{kfp:sat:eq:key}, then $h\ge Jh$ and irreducibility imply
$h>0$: a positive coordinate propagates along paths of powers of $J$.
The stated nonlinear monomial then makes $Q_t(h)$ nonzero. Thus
$\ell^TQ_t(h)>0$, contradicting~\eqref{kfp:sat:eq:perron-contradiction} even
when $\satR_N=0$. Consequently $h=0$.
\end{proof}

\begin{remark}[Why this is stronger than a prefix-Jacobian test]
The repository's earlier obstruction uses
$J_N=\Phi_X(t,P_N(t))$. At a true critical point, every finite $J_N$ can
still be strictly subcritical. The key improvement is to first establish
$v\ge A(t)$ and expand at the \emph{full} profile. This replaces the
finite-prefix forcing by the uncorrected full defect. No limit interchange
at a singularity is needed.
\end{remark}

\section{Exact spectral saturation of the hierarchy}
\label{kfp:sat:sec:saturation}

\subsection{Subcritical existence with explicit slack}
For $0<t<R$, set $a(t)=A(t)$ and $J(t)=\Phi_X(t,a(t))$.

\begin{proposition}[Quantitative supersolution test]\label[proposition]{kfp:sat:prop:quadratic-test}
Suppose $w>0$, $\alpha>0$, and $C\ge0$ satisfy
\[
 (I-J(t))w\ge\alpha\bone,
\]
and, for $0\le s\le s_0$,
\[
 \Phi(t,a+sw)-\Phi(t,a)-sJ(t)w\le Cs^2\bone.
\]
Let $r_N=\norm{\satR_N(t)}_\infty$. If
\begin{equation}\label{kfp:sat:eq:quadratic-test}
 0<s\le s_0,\qquad r_N+Cs^2<\alpha s,
\end{equation}
then $v=a+sw$ is a strict upper certificate. If $\Phi,D_N,t$ have rational
data, a sufficiently close rational vector is also a certificate.
\end{proposition}
\begin{proof}
The right side of the corrected equation at $a+sw$, minus $a+sw$, is
\[
 \satR_N(t)-s(I-J(t))w+Q_t(sw)
 \le (r_N-\alpha s+Cs^2)\bone<0.
\]
Also $a+sw>P_N(t)$, so both inequalities have strict slack. Polynomial
continuity preserves them in a small neighborhood, which contains rational
vectors.
\end{proof}

For $C>0$, the quadratic budget has maximum $\alpha^2/(4C)$, reached at
$s=\alpha/(2C)$ when this value is allowed by $s_0$. This gives an explicit
reason to compare the remaining defect with the \emph{square} of the
spectral gap. The test is analytic as stated: $a(t)$ need not be known
exactly. A directly computable version replaces $a,J(t),\satR_N(t)$ by
$P_N(t),\Phi_X(t,P_N(t)),b_N(t)$, where
\begin{equation}\label{kfp:sat:eq:finite-forcing}
 b_N(t)=\bigl(\shigh N\Phi(\xi,P_N(\xi))\bigr)(t).
\end{equation}
The same proof then constructs $P_N(t)+sw$ using finite polynomial data.

\begin{corollary}[Eventual subcritical certificates; prior strict criterion]
\label[corollary]{kfp:sat:cor:subcritical}
If $0<t<R$ and $\spr(J(t))<1$, all sufficiently large $N$ admit a strict
certificate at $t$. If $\Phi$ and $t$ are rational, and the profile
coefficients are rational, the certificate may be chosen rational.
\end{corollary}
\begin{proof}
Take $w=(I-J(t))^{-1}\bone=\sum_{k\ge0}J(t)^k\bone>0$ and $\alpha=1$.
A polynomial Taylor remainder provides a finite $C$ for $0\le s\le1$.
Choose a fixed small $s>0$ with $Cs^2<s/2$. Since the nonnegative series
$D$ converges at $t$, its tail $r_N$ tends to zero. For large $N$, $r_N<s/2$, giving the strict
inequality~\eqref{kfp:sat:eq:quadratic-test}. Rationality follows from strict slack.
\end{proof}

\noindent\wtag\ \emph{Restatement.} This is Part~I's
\cref{kfp:thm:stability}(a), as the manuscript's title ``prior strict
criterion'' says. The proof above, through
Proposition~\ref{kfp:sat:prop:quadratic-test} and expansion at the true
profile, is kept as a second proof; Part~I expands at the prefix profile.

\subsection{The exact limiting radius}
\begin{theorem}[Spectral saturation]\label{kfp:sat:thm:limit}
Assume $R>0$, $D$ is not a polynomial, and $J(t)$ is irreducible for
$0<t<R$. Define
\begin{equation}\label{kfp:sat:eq:barrier}
 r_* = \sup\{t\in(0,R):\spr(J(t))<1\},
\end{equation}
with $r_*=R$ if the spectral radius remains below one throughout $(0,R)$.
Then
\begin{equation}\label{kfp:sat:eq:radius-limit}
 \rho_N\uparrow r_*,\qquad \mu_N\downarrow1/r_*.
\end{equation}
If $r_*<R$, then $\spr(J(r_*))=1$ and no finite-prefix series converges
coordinatewise at $r_*$. In particular, exact coefficient stabilization
does not force the limiting growth to equal the growth of $A$.
\end{theorem}
\begin{proof}
All entries of $J(t)$ are nonnegative power series, so $J(t)$ and its
spectral radius are nondecreasing in $t$. At zero the Jacobian is
nilpotent, hence its spectral radius is zero. Continuity shows that the
subcritical set is a nonempty initial interval. If its supremum is
interior to $(0,R)$, continuity also gives the value one there.

For every $t<r_*$, Corollary~\ref{kfp:sat:cor:subcritical} gives a certificate for
large $N$. Thus $\liminf\rho_N\ge t$, and then
$\liminf\rho_N\ge r_*$. Proposition~\ref{kfp:sat:prop:prefix} makes the radii
nondecreasing and bounds them by $R$. If $r_*=R$, this proves the claim.
If $r_*<R$, the full defect has a nonzero tail at every level, so
Theorem~\ref{kfp:sat:thm:critical} and Lemma~\ref{kfp:sat:lem:numeric} give
$\rho_N\le r_*$ and nonconvergence at $r_*$. This proves the result in
the remaining case.
\end{proof}

\begin{remark}[The finite-defect exception is exact]
If $D$ is a polynomial, then $B_N=A$ for all sufficiently large $N$,
regardless of the Jacobian farther along the real axis. This is not a
small technical exception to Theorem~\ref{kfp:sat:thm:limit}; it is a genuinely
different regime with finite termination.
\end{remark}

\section{Tail flatness without coefficient regularity}
\label{kfp:sat:sec:tails}

A sharp rate theorem must control how the omitted defect changes when the
evaluation point moves with $N$. A naive argument might assume that its
coefficients have a smooth asymptotic form or are not lacunary. Neither
assumption is needed.

\begin{lemma}[A square-root tail estimate]\label[lemma]{kfp:sat:lem:tail-flatness}
Let $H(\xi)=\sum_{n\ge0}h_n\xi^n$ have nonnegative coefficients and radius of
convergence greater than $r>0$. Assume $H$ is not a polynomial, and write
\[
 H_N(\xi)=\sum_{n>N}h_n\xi^n,\qquad \varepsilon_N=\sqrt{H_N(r)}.
\]
For every fixed $K<\infty$, uniformly on $|t-r|\le K\varepsilon_N$,
\begin{equation}\label{kfp:sat:eq:tail-flat}
 \frac{H_N(t)}{H_N(r)}\longrightarrow1,
 \qquad H_N'(t)=o(\varepsilon_N),
 \qquad H_N''(t)=o(1).
\end{equation}
\end{lemma}
\begin{proof}
Choose $0<a<r<b$ inside the convergence disk. For $t\in[a,b]$,
Cauchy--Schwarz gives
\begin{align}
 H_N'(t)
 &=t^{-1}\sum_{n>N}n h_nt^n\nonumber\\
 &\le a^{-1}\sqrt{H_N(t)}
       \sqrt{\sum_{n>N}n^2h_nb^n}.
 \label{kfp:sat:eq:tail-cs}
\end{align}
The final sum tends to zero because the power series is analytic beyond
$b$. Every tail is strictly positive on $(0,\infty)$, so
\[
 \left|\frac{d}{dt}\sqrt{H_N(t)}\right|
 \le \frac1{2a}\sqrt{\sum_{n>N}n^2h_nb^n}=c_N\longrightarrow0.
\]
For $|t-r|\le K\varepsilon_N$, it follows that
\[
 \left|\frac{\sqrt{H_N(t)}}{\varepsilon_N}-1\right|
 \le Kc_N\longrightarrow0.
\]
This proves the first assertion. Substitution in~\eqref{kfp:sat:eq:tail-cs}
gives the second. Finally,
\[
 0\le H_N''(t)
 \le a^{-2}\sum_{n>N}n(n-1)h_nb^n\longrightarrow0.
\]
All estimates are uniform on the specified shrinking interval.
\end{proof}

\begin{corollary}[Vector version]\label[corollary]{kfp:sat:cor:vector-tails}
Let $D$ be a nonnegative vector series analytic beyond $r$, let $\ell>0$,
and set
\[
 \eta_N=\ell^T\satR_N(r)>0,\qquad \varepsilon_N=\sqrt{\eta_N}.
\]
Uniformly for $|t-r|\le K\varepsilon_N$,
\[
 \ell^T\satR_N(t)=\eta_N(1+o(1)),\quad
 \norm{\satR_N(t)}=O(\eta_N),\quad
 \norm{\satR_N'(t)}=o(\varepsilon_N),\quad
 \norm{\satR_N''(t)}=o(1).
\]
\end{corollary}
\begin{proof}
Apply Lemma~\ref{kfp:sat:lem:tail-flatness} to $H=\ell^TD$. Each component and
its nonnegative derivatives are bounded by a fixed multiple of the
corresponding weighted scalar sum. Norm equivalence in finite dimension
finishes the argument.
\end{proof}

The vector directions $\satR_N(r)/\eta_N$ need not converge. The proof
below only uses their boundedness and their fixed Perron projection. Thus
oscillating component contributions and arbitrarily long coefficient gaps
are allowed.

\section{Sharp convergence at an interior nonlinear crossing}
\label{kfp:sat:sec:sharp}

\subsection{Hypotheses and invariant constants}
Assume the hypotheses of Theorem~\ref{kfp:sat:thm:limit}, and suppose $r_*<R$.
Assume also $a_*=A(r_*)>0$ and that $\Phi$ contains a monomial of total
$X$-degree at least two. Put
\[
 J_*=\Phi_X(r_*,a_*),\qquad J_*u=u,\qquad \ell^TJ_*=\ell^T,
 \qquad u,\ell>0,\qquad \ell^Tu=1.
\]
Define
\begin{equation}\label{kfp:sat:eq:constants}
 \kappa=\ell^TJ'(r_*)u,
 \qquad
 \beta=\frac12\ell^T\Phi_{XX}(r_*,a_*)[u,u],
 \qquad
 \eta_N=\ell^T\satR_N(r_*).
\end{equation}
The derivative $J'$ is the \emph{total} derivative along the true profile:
\begin{equation}\label{kfp:sat:eq:Jderivative}
 J'(t)=\Phi_{\xi X}(t,A(t))+
       \Phi_{XX}(t,A(t))[A'(t),\,\cdot\,].
\end{equation}
For a vector-valued Hessian, the bracket means componentwise evaluation of
its bilinear form.

\begin{lemma}[Nondegeneracy is automatic here]\label[lemma]{kfp:sat:lem:positive-constants}
Under these assumptions, $\kappa>0$, $\beta>0$, and $\eta_N>0$ for every
$N$. If $\Lambda(t)=\spr(J(t))$ near $r_*$, then
\begin{equation}\label{kfp:sat:eq:lambda-expansion}
 1-\Lambda(t)=\kappa(r_*-t)+O((r_*-t)^2).
\end{equation}
\end{lemma}
\begin{proof}
The nonnegative analytic matrix $J(t)$ has nonnegative derivative. Some
entry is nonconstant: otherwise $J_*$ would equal the nilpotent $J(0)$,
contradicting its Perron value one. A nonconstant power series with
nonnegative coefficients has strictly positive derivative at every positive
point of convergence. Since $u,\ell>0$, this gives $\kappa>0$.

A nonlinear monomial evaluated at positive $r_*,a_*,u$ contributes
positively to at least one component of the Hessian evaluation; all other
contributions are nonnegative. Hence $\beta>0$. The infinite nonnegative
defect gives $\eta_N>0$. The simple Perron eigenvalue varies analytically
near $r_*$, and differentiation of its eigenvalue equation gives
$\Lambda'(r_*)=\ell^TJ'(r_*)u=\kappa$.
\end{proof}

Rescaling $u$ to $cu$ and $\ell$ to $\ell/c$ changes $\beta$ to $c\beta$
and $\eta_N$ to $\eta_N/c$. All the formulas below are invariant under
this change. Thus the normalization $\ell^Tu=1$ leaves no ambiguity in
any observable prediction.

\subsection{Statement of the sharp law}
\begin{theorem}[Square-root saturation law]\label{kfp:sat:thm:sharp}
Under the assumptions just stated,
\begin{equation}\label{kfp:sat:eq:sharp-radius}
 \boxed{\displaystyle
 r_*-\rho_N\sim\frac{2\sqrt\beta}{\kappa}\sqrt{\eta_N}.}
\end{equation}
For all sufficiently large $N$, $\rho_N<r_*$, $B_N(\rho_N)$ is finite,
and
\begin{equation}\label{kfp:sat:eq:sharp-height}
 \boxed{\displaystyle
 B_N(\rho_N)-A(\rho_N)
 =u\sqrt{\eta_N/\beta}+o(\sqrt{\eta_N}).}
\end{equation}
The vector error is in any fixed finite-dimensional norm. Consequently
\begin{equation}\label{kfp:sat:eq:sharp-growth}
 \mu_N-\frac1{r_*}
 \sim \frac{2\sqrt\beta}{\kappa r_*^2}\sqrt{\eta_N}.
\end{equation}
\end{theorem}

The proof has two parts. A local reduction identifies the fold. A separate
localization argument proves that it is the actual convergence boundary,
not a spurious branch of the polynomial equations. Omitting the second
part would leave a genuine gap.

\subsection{Uniform localization of every possible certificate}
\begin{lemma}[No remote fixed point can evade the fold]\label[lemma]{kfp:sat:lem:localization}
There exist constants $C,\delta_0>0$, independent of $N$, such that, for
$r_*-\delta_0<t<r_*$, every certificate at $(N,t)$ implies
\begin{equation}\label{kfp:sat:eq:localization}
 \norm{B_N(t)-A(t)}\le C(r_*-t).
\end{equation}
\end{lemma}
\begin{proof}
Let $h=B_N(t)-A(t)\ge0$. The full-defect equality gives
\begin{equation}\label{kfp:sat:eq:h-equation}
 h=J(t)h+Q_t(h)+\satR_N(t),
\end{equation}
so $h\ge J(t)h$. Irreducibility persists near $r_*$. The matrix
\[
 K(t)=I+J(t)+\cdots+J(t)^{d-1}
\]
has every entry bounded below by a positive constant on a sufficiently
small closed neighborhood of $r_*$. Since $K(t)h\le dh$, there is $c>0$
such that
\begin{equation}\label{kfp:sat:eq:h-comparable}
 h_i\ge c\norm h_1\quad(1\le i\le d)
\end{equation}
for every such $h$. This argument also covers $d=1$.

Let $\ell(t)>0$ be a continuous normalized left Perron vector. Its
coordinates are uniformly bounded above and below. The positive nonlinear
monomial in $\Phi$, together with positivity of $A(t)$, ensures that at
least one coefficient in the quadratic part of $Q_t$ is uniformly positive.
Using~\eqref{kfp:sat:eq:h-comparable}, we obtain
\[
 \ell(t)^TQ_t(h)\ge c_1\norm h_1^2
\]
after decreasing the neighborhood if needed. This bound holds for all
$h\ge0$ satisfying~\eqref{kfp:sat:eq:h-comparable}, not just for small $h$:
higher-degree terms of the expansion are nonnegative.

Multiplying~\eqref{kfp:sat:eq:h-equation} by $\ell(t)^T$ now gives
\[
 (1-\Lambda(t))\ell(t)^Th
 \ge c_1\norm h_1^2.
\]
Since $\ell(t)^Th\le c_2\norm h_1$ and
$1-\Lambda(t)=O(r_*-t)$, division when $h\ne0$ proves the bound. If
$h=0$, the bound is immediate (and a nonzero residual would in fact rule
that case out).
\end{proof}

\subsection{Reduction to a scalar equation}
Set $\delta=r_*-t$ and write
\[
 h=su+w,\qquad \ell^Tw=0,\qquad
 \Pi=I-u\ell^T.
\]
On the subspace $\ell^Tw=0$, the map $I-J_*$ is invertible. Indeed,
$J_*$ preserves this subspace and its only eigenvector with eigenvalue
one is a multiple of $u$. Other peripheral eigenvalues, including $-1$,
are harmless; primitivity is not required.

For an independent vector parameter $r$, consider
\begin{equation}\label{kfp:sat:eq:G}
 G(t,h,r)=h-\Phi(t,A(t)+h)+\Phi(t,A(t))-r.
\end{equation}
The analytic implicit function theorem applied to $\Pi G=0$ gives
\[
 w=W(\delta,s,r),\qquad W(0,0,0)=0,
\]
with
\begin{equation}\label{kfp:sat:eq:W-bound}
 \norm{W(\delta,s,r)}
 =O\bigl(|\delta s|+s^2+\norm r\bigr).
\end{equation}
To see the absence of a term proportional to $\delta$ alone, note that
$G(t,0,0)=0$ identically. The absence of a term linear in $s$ at
$\delta=0$ follows from $(I-J_*)u=0$.

Substitute $r=\satR_N(r_*-\delta)$ and define the scalar reduced equation
\begin{equation}\label{kfp:sat:eq:g}
 g_N(\delta,s)=
 \ell^TG\bigl(r_*-\delta,
 su+W(\delta,s,\satR_N(r_*-\delta)),
 \satR_N(r_*-\delta)\bigr).
\end{equation}
The local fixed points are exactly its zeros.

Put $\varepsilon_N=\sqrt{\eta_N}$. On bounded sets of $K,S$, set
$\delta=\varepsilon_NK$ and $s=\varepsilon_NS$. Taylor expansion,
\eqref{kfp:sat:eq:W-bound}, and Corollary~\ref{kfp:sat:cor:vector-tails} give
\begin{equation}\label{kfp:sat:eq:scaled-reduction}
 F_N(K,S):=\frac{g_N(\varepsilon_NK,\varepsilon_NS)}{\varepsilon_N^2}
 \longrightarrow F(K,S):=\kappa KS-\beta S^2-1
\end{equation}
in $C^2$ on every bounded rectangle.

Here are the terms responsible for the limit. The projected linear part is
\[
 \ell^T(I-J(r_*-\delta))su
 =\kappa\delta s+O(\delta^2|s|).
\]
The leading nonlinear part is $-\beta s^2$. Contributions involving $W$
are $O(\varepsilon_N^3)$ in this scaling, because
$W=O(\varepsilon_N^2)$ and $\ell^T(I-J_*)W=0$. Finally the projected
forcing is $-\eta_N(1+o(1))$. For the first and second scaled derivatives,
Corollary~\ref{kfp:sat:cor:vector-tails} gives exactly the needed bounds
$\satR_N'=o(\varepsilon_N)$ and $\satR_N''=o(1)$. This verifies the
$C^2$ assertion, including the moving forcing; a fixed-parameter Taylor
expansion alone would not suffice.

\subsection{The fold and its identification with the true radius}
The limiting equations $F=0$ and $F_S=0$ have the unique solution with
$K,S>0$
\begin{equation}\label{kfp:sat:eq:fold-limit}
 K_c=\frac{2\sqrt\beta}{\kappa},\qquad S_c=\frac1{\sqrt\beta}.
\end{equation}
At this point the Jacobian of $(F,F_S)$ with respect to $(K,S)$ is
\[
 \begin{pmatrix}\kappa/\sqrt\beta&0\\ \kappa&-2\beta\end{pmatrix},
\]
which is invertible. The $C^2$ convergence in~\eqref{kfp:sat:eq:scaled-reduction}
and the implicit function theorem give a nearby fold
$(K_N,S_N)\to(K_c,S_c)$ for every sufficiently large $N$. At that fold,
\begin{equation}\label{kfp:sat:eq:fold-derivatives}
 (g_N)_\delta=\kappa\varepsilon_N/\sqrt\beta+o(\varepsilon_N)>0,
 \qquad (g_N)_{ss}=-2\beta+o(1)<0.
\end{equation}
The associated $h=\varepsilon_NS_Nu+O(\varepsilon_N^2)$ is strictly
positive and solves the corrected equation, hence gives a certificate.

We must show that this fold is the boundary of all certificates. Fix a
small $a>0$. At $K=K_c+a$, the limiting polynomial $F(K,S)$ has two
positive simple roots. The implicit function theorem gives nearby roots
of $F_N$, and their vectors $h$ are positive; thus a certificate exists.
At $K=K_c-a>0$, no limiting positive root exists. If certificates existed
there along a subsequence, Lemma~\ref{kfp:sat:lem:localization} would make their
scaled $h$ bounded. The complementary equation forces $w=O(\varepsilon_N^2)$;
a convergent subsequence of $S$ would then solve $F(K_c-a,S)=0$, a
contradiction.

The convergence boundary is therefore trapped in an arbitrarily small
fixed scaled neighborhood of $K_c$. In that neighborhood every possible
fixed point is localized by Lemma~\ref{kfp:sat:lem:localization}, and every
localized solution lies near $S_c$. The reduced function is strictly
concave in $s$ there by~\eqref{kfp:sat:eq:fold-derivatives}; its unique local maximum
crosses zero at $(K_N,S_N)$. Its derivative with respect to $K$ at this
crossing is positive. Consequently local roots exist on the larger-$K$
side and none on the smaller-$K$ side. No other roots can supply a
certificate by the localization just proved. Thus
\[
 \rho_N=r_*-\varepsilon_NK_N.
\]
At the fold the fixed point is finite. Lemma~\ref{kfp:sat:lem:numeric} identifies
it with $B_N(\rho_N)$, since no other localized nonnegative fixed point
exists there. Its height is
$\varepsilon_NS_Nu+O(\varepsilon_N^2)$.
Together with~\eqref{kfp:sat:eq:fold-limit}, these facts prove
\eqref{kfp:sat:eq:sharp-radius} and~\eqref{kfp:sat:eq:sharp-height}. The reciprocal
expansion proves~\eqref{kfp:sat:eq:sharp-growth}, completing the proof of
Theorem~\ref{kfp:sat:thm:sharp}.

\section{Universal branch shape and a quarter-power amplitude}
\label{kfp:sat:sec:branch}

\begin{theorem}[Scaled branch and local singular amplitude]\label{kfp:sat:thm:branch}
Under Theorem~\ref{kfp:sat:thm:sharp}, put
\[
 S_-(K)=\frac{\kappa K-\sqrt{\kappa^2K^2-4\beta}}{2\beta}
 \quad\left(K>\frac{2\sqrt\beta}{\kappa}\right).
\]
Uniformly on compact subsets of this interval,
\begin{equation}\label{kfp:sat:eq:branch-shape}
 \frac{B_N(r_*-K\varepsilon_N)-A(r_*-K\varepsilon_N)}{\varepsilon_N}
 \longrightarrow uS_-(K).
\end{equation}
For each sufficiently large fixed $N$, the minimal branch has a local
expansion as real $t\uparrow\rho_N$,
\begin{equation}\label{kfp:sat:eq:local-sqrt}
 B_N(t)=B_N(\rho_N)-a_N\sqrt{1-t/\rho_N}
                  +O(1-t/\rho_N),\qquad a_N>0,
\end{equation}
and
\begin{equation}\label{kfp:sat:eq:amplitude}
 \boxed{\displaystyle
 a_N=u\sqrt{r_*\kappa}\,\beta^{-3/4}\eta_N^{1/4}
       +o(\eta_N^{1/4}).}
\end{equation}
\end{theorem}
\begin{proof}
For fixed $K>K_c$, the two roots of the limiting reduced equation are
simple and positive. The corresponding local solution vectors are ordered
componentwise: their derivative with respect to $s$ is
$u+O(\varepsilon_N)>0$. The minimal fixed point selects the smaller
root. Localization excludes any other candidate, proving
\eqref{kfp:sat:eq:branch-shape}; compactness gives uniformity away from the fold.

For fixed $N$, the reduction is analytic in a neighborhood of its fold.
Since $(g_N)_\delta>0$ and $(g_N)_{ss}<0$, the analytic implicit function
theorem applied with the square-root parameter gives
\[
 s=s_N-
 \sqrt{\frac{2(g_N)_\delta}{-(g_N)_{ss}}}\,
        \sqrt{\rho_N-t}+O(\rho_N-t).
\]
The sign is negative on the minimal branch. The vector derivative with
respect to $s$ is $u+O(\varepsilon_N)$; all other changes are analytic in
$\rho_N-t$. Therefore
\[
 a_N=(u+O(\varepsilon_N))
     \sqrt{\frac{2\rho_N(g_N)_\delta}{-(g_N)_{ss}}}.
\]
Substitution of~\eqref{kfp:sat:eq:fold-derivatives}, $\rho_N\to r_*$, and
$\varepsilon_N=\eta_N^{1/2}$ proves~\eqref{kfp:sat:eq:amplitude}.
\end{proof}

The local singularity is a genuine singularity in every component because
$a_N>0$. The theorem deliberately does not claim an unqualified global
coefficient asymptotic $c_N\rho_N^{-n}n^{-3/2}$. Such a statement requires
control of other singularities on the circle of convergence and the
appropriate analytic continuation or periodicity hypotheses. A real
positive-axis expansion alone does not establish those conditions.

\subsection{Consequences for data requirements}
If a particular system has
\[
 \eta_N\sim C N^p\theta^N,\qquad 0<\theta<1,
\]
then Theorem~\ref{kfp:sat:thm:sharp} gives
\[
 r_*-\rho_N\sim\frac{2\sqrt{\beta C}}{\kappa}
                 N^{p/2}\theta^{N/2}.
\]
The exponential rate is halved. To make the radius gap of size
$\epsilon$, the leading required prefix length is
\[
 N=\frac{2\log(1/\epsilon)}{|\log\theta|}
       +O(\log\log(1/\epsilon)),
\]
when the displayed tail equivalent holds. The quarter-power amplitude
also quantifies how the singularity becomes weaker while its location
approaches the barrier. These are conditional consequences of a tail law,
not assumptions used to prove the general theorem.

\section{An exact scalar model: limit, rate, and certificates}
\label{kfp:sat:sec:scalar}

\subsection{The previously known limit}
Take
\begin{equation}\label{kfp:sat:eq:scalar-model}
 A(\xi)=\frac \xi{1-\xi},\qquad \Phi(\xi,X)=\xi+X^2.
\end{equation}
The inequality $A\preceq \xi+A^2$ holds because the coefficient of $\xi^n$
in $A^2$ is $n-1$ for $n\ge2$. The defect is
\begin{equation}\label{kfp:sat:eq:scalar-defect}
 D(\xi)=\frac{\xi^3}{(1-\xi)^2}
     =\sum_{n\ge3}(n-2)\xi^n.
\end{equation}
The true radius is $1$, but
\[
 J(t)=\frac{2t}{1-t}
\]
first reaches one at $r_*=1/3$. Thus the corrected radii tend to $1/3$
and their growth tends to $3$. This limit and its interpretation were
already established in the repository \cite{satprefix}.
\wtag\ It is Part~I's \cref{kfp:thm:gap}, whose proof is not
repeated here; what follows is new in the manuscript.

For $N\ge2$, the corrected equation and its selected root are
\begin{equation}\label{kfp:sat:eq:scalar-root}
 B_N=\xi+B_N^2-D_N(\xi),\qquad
 B_N(\xi)=\frac{1-\sqrt{1-4\xi+4D_N(\xi)}}2.
\end{equation}
The branch is the one with zero constant term. The positive-tail
construction, rather than the signed appearance of the radicand,
guarantees nonnegative coefficients.

\subsection{A sharp equivalent and an exact root equation}
Here $u=\ell=1$, $\beta=1$, and
\[
 \kappa=J'(1/3)=\frac92.
\]
A direct geometric-series computation gives
\begin{equation}\label{kfp:sat:eq:scalar-tail}
 \satR_N(t)=
 \frac{t^{N+1}\bigl((N-1)-(N-2)t\bigr)}{(1-t)^2},
 \qquad
 \eta_N=\satR_N(1/3)=\frac{2N-1}{4\cdot3^N}.
\end{equation}
Therefore the new sharp equivalents specialize to
\begin{equation}\label{kfp:sat:eq:scalar-rate}
 \boxed{\displaystyle
 \frac13-\rho_N\sim\frac29\sqrt{2N-1}\,3^{-N/2},
 \qquad
 \mu_N-3\sim2\sqrt{2N-1}\,3^{-N/2}.}
\end{equation}
The singular amplitude satisfies
\begin{equation}\label{kfp:sat:eq:scalar-amplitude}
 a_N\sim\sqrt{\frac32}\left(\frac{2N-1}{4\cdot3^N}\right)^{1/4}.
\end{equation}

These constants can also be checked without the general theorem. Since
\[
 1-4t+4D(t)=\frac{(1-3t)^2}{(1-t)^2},
\]
the exact root equation for $\rho_N$ is
\begin{equation}\label{kfp:sat:eq:scalar-polynomial}
 (1-3\rho_N)^2
 =4\rho_N^{N+1}\bigl((N-1)-(N-2)\rho_N\bigr).
\end{equation}
For $N\ge3$ there is exactly one root in $(1/4,1/3)$.
Indeed, $c_N(t)=t-D_N(t)$ is strictly increasing on $[1/4,1/3]$:
$D_N'(t)<D'(t)\le1$ there. Also $c_N(1/4)<1/4$ and
$c_N(1/3)>1/4$. The root is where $c_N=1/4$.
At that point $v=1/2>P_N(\rho_N)$ is a valid certificate.
Immediately to its right the quadratic has no real root, proving that
this is exactly the radius.

To extract the equivalent directly, first the root equation gives
$\rho_N\to1/3$. Its positive square root gives
\[
 \frac13-\rho_N
 =\frac23\rho_N^{(N+1)/2}
       \sqrt{(N-1)-(N-2)\rho_N}.
\]
The right side is $O(\sqrt N\,3^{-N/2})$, so
$N(1/3-\rho_N)\to0$. It follows that
$(3\rho_N)^{(N+1)/2}\to1$. Substitution now yields exactly the first
formula in~\eqref{kfp:sat:eq:scalar-rate}; the second follows by reciprocation.
This is an independent derivation of the leading constants.

\subsection{Exact computation and numerical illustration}
The supplied checker brackets the root of~\eqref{kfp:sat:eq:scalar-polynomial}
by rational bisection, using exact integer arithmetic in every sign test.
It performs $260$ bisections for each displayed $N$. Both rational
endpoints and their exact width are saved in
\path{data/02-spectral-saturation-verification.json}. The entries below are rounded displays
of those brackets; the last column divides the gap by the proposed leading
term in~\eqref{kfp:sat:eq:scalar-rate}.

\begin{table}[htbp]
\centering\small
\begin{tabular}{@{}rrrr@{}}
\toprule
$N$ & $\rho_N$ (rounded) & $\rho_* -\rho_N$ & ratio to leading term\\
\midrule
4 & 0.287371536944 & $4.59618\times10^{-2}$ & 0.703563009 \\
8 & 0.323938104014 & $9.39523\times10^{-3}$ & 0.884217871 \\
16 & 0.333145629046 & $1.87704\times10^{-4}$ & 0.995350175 \\
32 & 0.333333292359 & $4.09748\times10^{-8}$ & 0.999998001 \\
64 & 0.333333333333 & $1.35148\times10^{-15}$ & 1.000000000 \\
128 & 0.333333333333 & $1.03347\times10^{-30}$ & 1.000000000 \\
\bottomrule
\end{tabular}

\caption{The scalar radius approaches its barrier at the predicted half-rate.
The exact brackets, not the rounded decimal display, are the certificates.}
\label{kfp:sat:tab:scalar}
\end{table}

For $N=64$ and $128$, the rounded radius visually equals $1/3$, but the
positive gap is separately displayed. No equality with $1/3$ is being
asserted. The asymptotic theorem is proved analytically; agreement at six
finite values is a consistency check, not a proof for all $N$.

\section{Examples that delimit the hypotheses}
\label{kfp:sat:sec:exceptions}

\subsection{A supercritical system that becomes exact in finitely many steps}
Let $A(\xi)=\xi$ and $\Phi(\xi,X)=\xi+X^2$. Then $D=\xi^2$, so $B_N=A$ for
$N\ge2$. Nevertheless $\Phi_X(t,A(t))=2t$ is critical at $t=1/2$ and
supercritical afterward. The positive tail map is
\[
 \Psi_N(t,y)=2ty+y^2\qquad(N\ge2),
\]
whose least fixed point is zero even when its derivative exceeds one.
This is precisely the zero-forcing exception in
Theorem~\ref{kfp:sat:thm:critical}. It disproves any formulation that rules out
all supercritical certificates without examining the residual defect.

\subsection{Irreducibility cannot be removed}
Consider the independent two-component system
\[
 \Phi_1(\xi,X)=\xi+X_1^2,\qquad
 \Phi_2(\xi,X)=\xi+\xi X_2,
\]
with
\[
 A_1(\xi)=\xi,\qquad A_2(\xi)=\frac \xi{1-\xi/2}.
\]
The defects are
\[
 D_1=\xi^2,\qquad D_2=\frac{\xi^2}{2(1-\xi/2)}.
\]
At $t=1/2$ the true-profile Jacobian is
$\operatorname{diag}(1,1/2)$ and has spectral radius one. For every
$N\ge2$, the first component is already exact, while the second corrected
series is finite because its linear denominator is $1-\xi$. Thus
certificates exist although the total defect tail is nonzero. The critical
block receives no residual forcing. A reducible extension must track
which blocks are forced and how they feed one another.

\subsection{Linearity gives a different convergence regime}
The scalar second component of the preceding example has true radius $2$
and Jacobian $J(t)=t$, hence spectral barrier $1$. Its corrected equation is
\[
 B_N=\xi+\xi B_N-D_N(\xi).
\]
Subtracting the true equation gives the exact identity
\[
 B_N(\xi)-A(\xi)=\frac{\satR_N(\xi)}{1-\xi}.
\]
Every finite level has a nonzero positive residual at $\xi=1$, so every
radius is \emph{exactly} $1$ and the boundary value diverges. There is no
nonzero square-root radius gap. Here $\beta=0$; the nonlinear hypothesis
of Theorem~\ref{kfp:sat:thm:sharp} is essential.

\subsection{Primitivity is not required}
Let $A_1=A_2=\xi/(1-\xi)$ and
\[
 \Phi_1=\xi+X_2^2,\qquad \Phi_2=\xi+X_1^2.
\]
The minimal branch remains symmetric and equals the scalar example in
both coordinates. At $t=1/3$ the Jacobian is
$\left(\begin{smallmatrix}0&1\\1&0\end{smallmatrix}\right)$, with
eigenvalues $1$ and $-1$. The complement of the Perron direction remains
invertible for $I-J$. Thus the sharp law applies despite period two in
the dependency matrix. This matrix issue is separate from possible
periodicity of generating-function coefficients.

\section{An exact all-prefix barrier for the polyomino map}
\label{kfp:sat:sec:bui}

\subsection{The finite input and the infinite conclusion}
The repository report supplies a $17$-component marked profile and a
$53$-monomial positive map, transcribed in Section~\ref{kfp:sat:app:map}.
The variables are ordered
\[
 (c,d,e,f,g,h,p,q,r,s,t,u,v,w,x,y,z).
\]
In this section and the appendix, the \emph{size variable} is $\xi$,
so it is not confused with the last coordinate $z$.
\wtag\ (In Part~II as printed, $\xi$ is the size variable throughout;
see \cref{kfp:sat:sec:notation}. The profile is Part~I's
\cref{kfp:app:data} and \texttt{data/profiles\_18.csv}; the map is
\eqref{eq:map}.)
The source data list all marked coefficients through size $18$
\cite{satprofile,satmap}. We use those values as input. The checker replays their
partition and recurrence consistency identities, but does not independently
enumerate size-$18$ polyominoes.
\wtag\ (Part~I's intake regenerated that table through size~$18$ from the
final enumerator source, and through size~$10$ with an independent
enumerator; see \cref{kfp:app:provenance}. It is still not kernel-checked.)

A finite lower bound on the Jacobian already suffices to rule out
\emph{every} longer prefix. The following elementary lemma makes the
certificate format explicit and avoids any reliance on a floating-point
eigenvalue.

\begin{lemma}[A rational spectral obstruction]\label[lemma]{kfp:sat:lem:finite-obstruction}
Let $J\ge0$ be irreducible, let $b\ge0$ be nonzero, and suppose there
exist $w>0$ and $\gamma>1$ with $Jw\ge\gamma w$. There is no $y\ge0$
satisfying $y\ge b+Jy$.
\end{lemma}
\begin{proof}
Iteration gives
$y\ge\sum_{k=0}^{d-1}J^kb$. Strong connectivity and $b\ne0$ make this
sum strictly positive, so $y>0$. Put $m=\min_i y_i/w_i>0$ and choose an
index $i$ attaining the minimum. Then
\[
 (Jy)_i\ge m(Jw)_i\ge m\gamma w_i>mw_i=y_i,
\]
contradicting $y\ge Jy$. Every ingredient is a finite matrix or graph
inequality.
\end{proof}

\subsection{The certified witness}
Set
\begin{equation}\label{kfp:sat:eq:bui-t}
 t_0=\frac{10000}{43149},\qquad
 \gamma=\frac{1000001}{1000000},\qquad
 J_{18}=\Phi_X(t_0,P_{18}(t_0)).
\end{equation}
The positive integer vector in Table~\ref{kfp:sat:tab:witness} satisfies
\begin{equation}\label{kfp:sat:eq:bui-check}
 \boxed{J_{18}w\ge\gamma w>w.}
\end{equation}
The $17$ inequalities are checked using exact rational arithmetic. A
second implementation evaluates the directional derivative by dual-number
multiplication and agrees exactly with the matrix-vector product. The
smallest ratio $(J_{18}w)_i/w_i-1$, rounded only for display, is
$0.000004934886683313896$, safely greater than the certified
$0.000001$.

\begin{table}[htbp]
\centering\small
\begin{tabular}{@{}lr@{}}
\toprule
coordinate & $w_i$\\
\midrule
$c$ & 53,709,887,569 \\
$d$ & 112,782,022,377 \\
$e$ & 231,753,937,549 \\
$f$ & 1,000,000,000,000 \\
$g$ & 486,645,549,886 \\
$h$ & 301,166,180,537 \\
$p$ & 513,359,385,008 \\
$q$ & 254,894,013,881 \\
$r$ & 120,854,740,222 \\
$s$ & 188,385,644,383 \\
$t$ & 169,174,579,449 \\
$u$ & 356,544,026,070 \\
$v$ & 97,749,895,723 \\
$w$ & 73,183,659,458 \\
$x$ & 71,425,518,584 \\
$y$ & 47,671,677,169 \\
$z$ & 28,731,292,790 \\
\bottomrule
\end{tabular}

\caption{An integer witness for~\eqref{kfp:sat:eq:bui-check}. The column headed
$w_i$ is a vector entry, not the map coordinate named $w$.}
\label{kfp:sat:tab:witness}
\end{table}

\begin{theorem}[A $4.3149$ barrier for every prefix]\label{kfp:sat:thm:bui}
For the repository's fixed $17$-variable map and its true marked profile,
with the size-$18$ profile taken as the supplied enumerative input, no
finite-prefix upper certificate exists at $t_0=10000/43149$.
Every corrected system has growth at least
\begin{equation}\label{kfp:sat:eq:bui-wall}
 \mu_N\ge\frac{43149}{10000}=4.3149.
\end{equation}
The same radius bound applies to its $g$ component, which is the component
used to majorize the unmarked fixed-polyomino counts.
\end{theorem}
\begin{proof}
For $M\ge18$, nonnegativity of the true coefficients gives
$P_M(t_0)\ge P_{18}(t_0)$, hence
\[
 J_Mw\ge J_{18}w\ge\gamma w,
 \qquad J_M=\Phi_X(t_0,P_M(t_0)).
\]
All prefix coordinates are positive, and the map's dependency graph is
strongly connected. The checker verifies this graph directly; positivity
means its edges coincide with those of $J_M$.

The first row of the map is $\Phi_c=\xi+\xi e$. For every size $M$, the
vertical column of $M$ cells has an $e$-pattern occurrence at its bottom
cell. In the source definition of $e$, the anchor is occupied while its
horizontal neighbors and the three positions immediately below are absent;
the vertical column has exactly this local configuration. Thus the true
coefficient $e_M\ge1$. It follows that the finite-prefix forcing
\[
 b_M(t_0)=\bigl([\xi^{>M}]\Phi(\xi,P_M(\xi))\bigr)(t_0)
\]
has a strictly positive $c$ coordinate, containing $e_Mt_0^{M+1}$.

A certificate would yield a tail $y\ge0$ with
$y\ge\Psi_M(t_0,y)\ge b_M(t_0)+J_My$.
Lemma~\ref{kfp:sat:lem:finite-obstruction} rules this out. Hence no level
$M\ge18$ has a certificate. An earlier certificate would persist to level
$18$ by Lemma~\ref{kfp:sat:lem:numeric}, so every earlier level is excluded too.
This proves the system-radius and growth assertions.

For completeness, when $M\ge18$ the linear-in-$Y$ part of the positive
tail map is the formal matrix $\Phi_X(\xi,P_M(\xi))$. Each dependency
edge contains a positive monomial $c\xi^kY_j$. Therefore
$Y_{M,i}\succeq c\xi^kY_{M,j}$ along every edge. Strong connectivity
forces all nonzero tail components to have the same radius, and the
positive forcing propagates to every component. In particular the $g$
component has the system radius. For earlier levels, componentwise
nesting with $B_{18}$ gives the same upper bound on the $g$ radius.
\end{proof}

\subsection{What this numerical result means}
The theorem is an obstruction to an entire infinite hierarchy, obtained
from one finite profile and one short rational witness. It shows that,
even with arbitrarily long exact prefixes and arbitrarily sophisticated
supersolution searches, the \emph{unchanged} map cannot certify an upper
growth constant below $4.3149$ by this method.

It does not say that the true polyomino growth is at least $4.3149$.
The logical direction is
\[
 \text{true counts}\ \preceq\ \text{corrected majorants},
\]
so a lower bound for the majorants cannot be reversed into a lower bound
for the true counts. Nor does it locate the actual saturation constant
exactly: $4.3149$ is only a certified lower bound on the limiting majorant
growth. The true fixed-map barrier could be higher.

\wtag\ \emph{Tempting false readings.} (i) ``$\lambda\ge4.3149$'': false as
a reading of \cref{kfp:sat:thm:bui}, which bounds the growth $\mu_N$ of the
majorants $B_N$ from below, while $\lambda$ lies below every $\mu_N$.
(ii) ``$4.3149<4.498$, so Part~II improves Part~I's bound'': false; $4.3149$
is not an upper bound on anything. It is a floor under every upper bound that
this fixed map and exact prefixes can certify. Part~I's certificates
$4.5233,\ldots,4.498$ (\cref{kfp:tab:progress}) and the Lean-verified
$4.5235$ all lie above it, as they must. The upper bounds on $\lambda$ in
\repo\ remain $4.498$ (computer-assisted, \cref{kfp:thm:main}) and
$9047/2000$ (Lean-verified).

The calculation also does not assume that the full marked profile
converges at $t_0$. It uses only finite prefixes, a uniform nonzero-forcing
construction, and certificate persistence. This makes it valid even before
the hypotheses needed to apply the sharp interior-crossing theorem to the
actual polyomino profile have been established.

\subsection{Discovery and replay are separate}
The optional script \path{code/02-spectral-saturation-discover.py} finds a numerical positive
eigenvector, scales and rounds it to integers, and exports the witness.
That procedure is not trusted by the proof. The replay script
\path{code/02-spectral-saturation-verify.py} imports no numerical library and ignores all
floating-point discovery fields. It recomputes the exact rational Jacobian
from the supplied coefficient profile and the displayed monomial map.

The verifier additionally checks nilpotence of the zero-size linear part
(the index is $3$), the $53$-monomial count, strong connectivity, all
$323$ component/degree defect checks through degree $18$ including degree
zero, and $95$ checks for the five partition identities including degree
zero. These are useful transcription checks, not a substitute for the
combinatorial recurrence proof or the marked enumeration.

\section{Reproducibility and proof audit}
\label{kfp:sat:sec:audit}

\subsection{What is in the package}
\wtag\ The file names below are those shipped in \repo; the delivered
package used the same names without the prefix
\path{02-spectral-saturation-}. The manuscript's own \texttt{article.tex},
\texttt{article.pdf} and \texttt{README.md} are not shipped (Part~II is
their text), and neither is its copy of \texttt{data/profiles\_18.csv}, which
is byte-identical to Part~I's file of that name and is replaced by it.
The two small TeX tables in \texttt{data/} are
included by the source. The other retained files are:
\begin{itemize}[leftmargin=1.6em,itemsep=2pt]
\item \path{code/02-spectral-saturation-bui_model.py}, the transcribed map
and rational Jacobian; \path{code/02-spectral-saturation-verify.py}, the
exact checker; and optional \path{code/02-spectral-saturation-discover.py},
the numerical witness search.
\item \texttt{data/profiles\_18.csv}, the inherited marked profile (Part~I's
file); \path{data/02-spectral-saturation-bui_spectral_certificate.json}, the
integer witness; and \path{data/02-spectral-saturation-verification.json}
and \path{data/02-spectral-saturation-verification.txt}, the recorded replay
results, including exact scalar brackets. \emph{Not to be confused with}
Part~I's own \texttt{data/verification.json} and
\texttt{data/verification.txt}, which record the $4.498$ replay.
\item \path{02-spectral-saturation-SOURCES.md},
\path{02-spectral-saturation-PROOF_STATUS.md}, and a build script
\path{code/02-spectral-saturation-build.sh}.
\end{itemize}

To replay the checks without changing retained outputs, run
\begin{lstlisting}
python3 code/verify.py
\end{lstlisting}
To regenerate the verification record and table inputs, run
\begin{lstlisting}
python3 code/verify.py --write
\end{lstlisting}
\wtag\ These two commands are the delivered ones and assume the delivered
layout. In \repo\ they must not be run in place: the renamed scripts import
\texttt{bui\_model}, which is shipped as
\path{code/02-spectral-saturation-bui_model.py}, and \texttt{--write} writes
\texttt{data/verification.json}, which is Part~I's record. The report's
README gives a recipe that rebuilds the delivered layout in a scratch copy;
there the checker's output equals the shipped record and tables.
Only the optional discovery script needs NumPy. The exact checker uses
Python's standard library, principally arbitrary-precision integers and
\texttt{fractions.Fraction}. Failure of any required check raises an
exception; it does not silently downgrade to a floating-point comparison.

\subsection{Dependency structure of the mathematical claims}
The critical theorem depends only on formal comparison, positivity of the
full defect, and finite-dimensional Perron--Frobenius theory. The spectral
limit additionally uses eventual subcritical supersolutions and convergence
of nonnegative tails. The sharp law adds analyticity at an interior
crossing, a nonlinear positive Hessian direction, the tail-flatness lemma,
a finite-dimensional implicit-function reduction, and the global
localization lemma.

The numerical all-prefix obstruction follows a different chain:
\[
 \begin{gathered}
 \text{supplied marked prefix}\ +\ \text{displayed polynomial map}\\
 \Longrightarrow\ \text{exact rational matrix witness}\\
 \Longrightarrow\ \text{no certificate at every level }M\ge18\\
 \Longrightarrow\ \text{no certificate at any level, by persistence}.
 \end{gathered}
\]
The nonzero forcing at arbitrary size is proved by the vertical-column
construction, not extrapolated from the finite data. The checker verifies
arithmetic at size $18$; the infinite-level conclusion is the accompanying
mathematical argument.

\subsection{Adversarial checks on the argument}
Several potential failure modes have explicit safeguards. The tail budget
in~\eqref{kfp:sat:eq:certificate} is mandatory. The expansion at $A(t)$ is used
only after proving that every certificate dominates $A(t)$. Nonzero
forcing is distinguished from finite exact termination. Irreducibility
cannot be replaced by an unsupported claim that every nonnegative
Perron vector is strictly positive. The rate theorem uses the total
derivative~\eqref{kfp:sat:eq:Jderivative}, not the partial derivative in the size
variable alone. The local fold is identified with the actual radius using
Lemma~\ref{kfp:sat:lem:localization}. Tail-flatness is proved rather than inferred
from a guessed coefficient asymptotic. No global coefficient transfer
principle is invoked from the real-axis singularity alone.

These safeguards reduce identifiable risks; they do not constitute
independent refereeing or machine verification. In particular the
analytic implicit-function and localization arguments should receive
independent mathematical review before any publication-level novelty or
correctness claim is treated as settled.

\subsection{A feasible formalization sequence}
A Lean continuation can be divided into three independent layers. The
first layer formalizes the polynomial prefix identity, positivity of
$\Psi_N$, formal comparison, and persistence of numeric supersolutions.
It can reuse the existing recurrence and growth interfaces described in
the repository report \cite{satprefix,satformal}.

The second layer formalizes Lemma~\ref{kfp:sat:lem:finite-obstruction} and checks
the rational vector in Table~\ref{kfp:sat:tab:witness}. That lemma needs neither
complex analysis nor numerical eigenvalues. Its remaining combinatorial
input is a verified marked prefix, together with the simple vertical-column
occurrence lemma. The checker included here is not such a formal
enumerator.

The third layer develops the analytic saturation theorem and the sharp
fold analysis. Its reusable targets include a theorem for nonnegative
power-series tail square roots, Perron eigenvector perturbation, and a
one-dimensional kernel reduction with uniform parameters. It is
substantially larger than the finite rational check and need not block
formalizing the numerical obstruction first. No unfinished Lean file is
included under the appearance of a completed proof.

\section{Further research questions and projects}
\label{kfp:sat:sec:future}

\subsection{Locate the actual fixed-map polyomino saturation constant}
Theorem~\ref{kfp:sat:thm:bui} proves an obstruction but not its sharp location.
A natural next target is a rigorous interval for
$\spr(\Phi_X(t,A(t)))$ using upper and lower bounds on the marked
profile. Exact lower prefixes give one side; a validated majorant for the
remaining marked tail is needed for the other. Determine whether the
crossing is interior to the true marked radius, and hence whether
Theorem~\ref{kfp:sat:thm:sharp} applies to this particular map.

\subsection{Replace enumeration length by a better recurrence}
The $4.3149$ obstruction makes a concrete distinction between two research
investments: longer exact prefixes versus modifying the geometric
recurrence. A structural repair must change the map or its available
coefficient inequalities, not merely find a new supersolution of the old
corrected map. Identify which local overcounts most affect the Perron
crossing, prove a repair, and repeat the certificate audit.

\subsection{Resolve the reducible critical case}
The example in Section~\ref{kfp:sat:sec:exceptions} shows why the scalar value
$\spr(J)=1$ is insufficient in a reducible system. Seek a necessary and
sufficient block criterion involving residual forcing, reachability of
critical blocks, and nonlinear feedback within each block. Distinguish a
critical block that is already exact from one receiving arbitrarily small
forcing from downstream dependencies. This is a genuine extension beyond
the irreducible theorem proved here.

\subsection{Classify crossings at the true singularity}
Our sharp theorem assumes $r_*<R$. When $r_*=R$, the derivatives of $A$
may diverge and $\kappa$ need not be finite. Determine how the singular
behavior of $A$, the defect tail, and the Perron gap combine. One should
expect different exponents, not automatic reuse of the interior
square-root law.

\subsection{Obtain second-order truncation asymptotics}
The present leading law depends only on $\kappa$, $\beta$, and $\eta_N$.
A second-order expansion should involve the complementary inverse
$(I-J_*)^{-1}$ on $\ell^Tw=0$, higher derivatives of $\Phi$ and $A$,
and the vector direction of the tail. Characterize when oscillating tail
directions prevent a single second-order constant even though the leading
law remains universal.

\subsection{Uniform coefficient asymptotics in two limits}
For fixed $N$, an additional aperiodicity and continuation analysis may
turn~\eqref{kfp:sat:eq:local-sqrt} into a coefficient asymptotic. A more delicate
problem lets both $n$ and $N$ tend to infinity. The amplitude shrinks as
$\eta_N^{1/4}$ while the radius gap shrinks as $\eta_N^{1/2}$, suggesting a
transition when $n\sqrt{\eta_N}$ is of order one. Prove a uniform transfer
law and determine its relation to exact prefix agreement for $n\le N$.

\subsection{Certified optimization of the omitted defect}
For a fixed data budget, is prefix truncation the best verified correction,
or can a sparse collection of rigorously proved higher-degree defects be
more effective? The leading sensitivity is Perron-weighted, so unweighted
coefficient size is not the right objective. Extend the theory from
consecutive prefixes to arbitrary finite certified defect polynomials and
characterize the optimal allocation problem.

\subsection{Bit complexity of near-critical certificates}
The existence of a rational vector follows from strict slack, but the
smallest required denominator may deteriorate near the fold. Derive an
explicit bit-length bound in terms of the input map, an enclosure of its
Perron data, and the separation from $\rho_N$. Separate the cost of finding
a witness from the cost of checking it. The practical integer witness here
is an example, not a worst-case complexity theorem.

\subsection{Formal marked-profile certification}
A verified enumerator or compact enumeration certificate would remove the
main inherited computational assumption in the numerical application.
The size-$18$ profile contains much more information than the unmarked
polyomino sequence; matching unmarked totals is insufficient. A first
formal target could be a smaller prefix and a weaker but nontrivial
all-prefix spectral barrier, followed by an incremental extension.

\subsection{The distinct neighborhood-refinement question}
Bui's question about increasing the geometric neighborhood size remains
separate \cite{bui}. The fixed-map theorem does not rule out convergence
of a sequence of increasingly accurate maps. Formulate conditions on
$\Phi^{(k)}$ under which their spectral barriers tend to the true radius,
and construct a counterexample when those conditions fail. This would
connect the present negative limitation with a positive convergence theory
for genuinely improving geometric descriptions.

\section{Conclusion}
The exact-prefix method has an exact spectral boundary. In an irreducible
system, a nonzero uncorrected defect cannot be sustained at or beyond
Perron value one, even if every finite-prefix Jacobian remains below one.
This resolves the missing critical case and identifies the limiting radius
of the entire hierarchy.

At an interior nonlinear crossing the boundary is quantitative:
its distance is asymptotic to a square root of the omitted defect, with
an explicit invariant constant. The solution height, scaled branch, and
singular amplitude have compatible universal laws. The tail estimate
shows why coefficient gaps do not invalidate this conclusion.

For the concrete polyomino map, an exact finite certificate proves a
$4.3149$ limit on what prefix refinement alone can accomplish. The useful
conclusion is methodological: to cross that threshold, one must improve
the recurrence information itself, not merely compute a longer initial
segment of the same profile. The abstract theorems are proved here; the
remaining work is independent validation, formal certification, and the
geometric refinement needed to move the spectral barrier.

\section{The complete polynomial map used in the certificate}
\label{kfp:sat:app:map}

\wtag\ This section is the manuscript's Appendix~A. Its display of the
$53$ monomials, transcribed from Part~I's \texttt{code/model.py}
\cite{satmap}, agrees row for row with \eqref{eq:map}, and is printed once,
there. The manuscript names each row by its lowercase output coordinate
($\Phi_c,\ldots,\Phi_z$), where \eqref{eq:map} uses the uppercase type name
($\Phi_C,\ldots,\Phi_Z$); the size variable is $\xi$, the lowercase $z$ is a
profile coordinate, and every coefficient is one.

The five linear partition identities in the supplied profile are
\[
 f=g+p,\qquad g=e+q,\qquad h=d+s,\qquad r=y+w,\qquad t=x+v.
\]
The checker verifies these coefficientwise through size $18$. The
nonnegative defects for the remaining rows are checked as well.

The file \texttt{data/profiles\_18.csv} retains the complete input rather
than just a floating-point value of $P_{18}(t_0)$. As a result, the reader
can reconstruct every exact rational entry of $J_{18}$ directly from the
map and the coefficient table. The unmarked total column is retained for
provenance, but is not an input to the spectral inequality.

\section{A minimal replay algorithm}
\label{kfp:sat:app:algorithm}

The numerical core can be expressed as the following exact computation.
For each map monomial $\xi^kX_{i_1}\cdots X_{i_m}$, its contribution to
$\Phi_X(t_0,a)w$ is
\[
 t_0^k\sum_{r=1}^m w_{i_r}\prod_{s\ne r}a_{i_s}.
\]
Repeated variables occur repeatedly in the sum, giving the correct
multiplicity. For example the derivative of $\xi e^2$ contributes
$2t_0a_e w_e$, not merely $t_0a_e w_e$.

\begin{lstlisting}
t0 = Fraction(10000, 43149)
factor = Fraction(1000001, 1000000)
a = [horner(profile, t0) for profile in profiles]
J = exact_jacobian(t0, a)
image = exact_matrix_vector_product(J, witness)
assert all(image[i] >= factor*witness[i] for i in range(17))
assert all(witness[i] > 0 for i in range(17))
assert strongly_connected(J)
\end{lstlisting}

The actual checker uses explicit exceptions rather than relying on Python
\texttt{assert}, which can be disabled by an optimization flag. The
pseudocode above only describes the mathematical test. The exact
arithmetic result is independent of the floating-point approximation used
to discover the witness.

For the scalar example, the polynomial sign oracle is
\[
 H_N(t)=(1-3t)^2-4t^{N+1}\bigl((N-1)-(N-2)t\bigr).
\]
The proved uniqueness of its root in $[1/4,1/3]$ turns rational bisection
into a root certificate. The program checks both endpoint signs again
before recording each bracket. It also evaluates the identity
\[
 H_N(t)=(1-t)^2\bigl(1-4t+4D_N(t)\bigr)
\]
at three independent rational points for every displayed $N$. Those
finite identity checks are safeguards against transcription errors;
the all-$t$ identity is the algebraic derivation in
Section~\ref{kfp:sat:sec:scalar}.

\section{Source snapshot and evidence ledger}
\label{kfp:sat:app:sources}

The repository snapshot inspected for this continuation is
\begin{center}\small\texttt{74f7f5bdba1aa728b340509701a0fa4e45e8dbf3}\end{center}
The main source directory is
\begin{quote}\small\ttfamily
SetTheory/Cardinals/docs/reports/\\
enumerative-combinatorics/\\
polyomino-growth-finite-prefix-corrections/
\end{quote}
Its \texttt{article.tex}, especially Section~8 and the proposed research
questions in Section~10, is the mathematical continuation source. Its
\texttt{code/model.py} supplies the map and local pattern convention;
\texttt{data/profiles\_18.csv} supplies the finite marked counts.
The complete paths, blob identifiers returned by the connector, and
review boundaries are retained in \path{02-spectral-saturation-SOURCES.md}.
\wtag\ That directory is this report, and its \texttt{article.tex} at the
pin (blob \texttt{4a021fb3}) is Part~I as printed here: Section~8 is
\cref{kfp:sec:stability} and Section~10 is \cref{kfp:sec:questions}. The
blobs recorded for \texttt{code/model.py} (\texttt{97b241bd}) and
\texttt{data/profiles\_18.csv} (\texttt{82648328}) are also those of the
placement commit.

\begin{center}\small
\begin{tabular}{@{}>{\raggedright\arraybackslash}p{.27\linewidth}p{.59\linewidth}@{}}
\toprule
Evidence & What it establishes here\\
\midrule
Repository article and README & The exact continuation target and prior
results; not blanket verification of the repository.\\[3pt]
Repository monomial map & The finite algebraic system used by the checker;
its combinatorial interpretation is an inherited source input.\\[3pt]
Repository marked CSV & The size-$18$ coefficients used in the arithmetic
certificate; not freshly enumerated here.\\[3pt]
Exact checker output & The stated rational inequalities, graph properties,
finite consistency identities, and scalar root brackets.\\[3pt]
Written proofs & The infinite-level implications, critical classification,
spectral limit, and sharp analytic laws.\\[3pt]
Primary literature review & Context and attribution for the recurrence,
positive-system certificates, Newton theory, and singularity methods;
not an exhaustive novelty search.\\
\bottomrule
\end{tabular}
\end{center}

\clearpage
\begin{thebibliography}{20}
\small
\bibitem{bui}
Vuong Bui, \emph{A convolutional approach to bounding the number of
polyominoes}, arXiv:2511.00461v2, 6 May 2026.
\url{https://arxiv.org/abs/2511.00461v2}.
The relevant locations are Section~3, Question~2, and Section~4,
Lemma~4 with Appendix~B. Version~1 had a different title.

\bibitem{repo}
Vladimir Reshetnikov, \emph{ProveIt}, repository snapshot
\sha.
\url{https://github.com/VladimirReshetnikov/ProveIt}.
Repository sources in the next entries are read at this commit.

\bibitem{reporeadme}
\repo, \emph{Klarner's polyomino growth constant}, project README.
\url{https://github.com/VladimirReshetnikov/ProveIt/blob/e21766d04c2b8a9b2cdba0cd43e563024b3bd1b9/Combinatorics/Polyominoes/KlarnerConstant/README.md}.

\bibitem{repocert}
\repo, \texttt{Support/verify\_certificate.py}, independent rational replay
of the $4.5235$ certificate.
\url{https://github.com/VladimirReshetnikov/ProveIt/blob/e21766d04c2b8a9b2cdba0cd43e563024b3bd1b9/Combinatorics/Polyominoes/KlarnerConstant/Support/verify_certificate.py}.

\bibitem{repopatterns}
\repo, \texttt{Lean/KlarnerConstant/Patterns.lean}, the seventeen required
and forbidden offset sets.
\url{https://github.com/VladimirReshetnikov/ProveIt/blob/e21766d04c2b8a9b2cdba0cd43e563024b3bd1b9/Combinatorics/Polyominoes/KlarnerConstant/Lean/KlarnerConstant/Patterns.lean}.

\bibitem{repocomplete}
\repo, \texttt{Lean/KlarnerConstant/GeometricComplete.lean}, composition of
the geometric recurrences and the existing growth endpoint.
\url{https://github.com/VladimirReshetnikov/ProveIt/blob/e21766d04c2b8a9b2cdba0cd43e563024b3bd1b9/Combinatorics/Polyominoes/KlarnerConstant/Lean/KlarnerConstant/GeometricComplete.lean}.

\bibitem{repobound}
\repo, \emph{An exact improvement to $\lambda\leq4.5235$ for Klarner's
constant}, research note dated 12 August 2026,
\lean{Combinatorics/Polyominoes/KlarnerConstant/Research/klarner-bound-4.5235.md}.
\wtag\ Added in the write; its Sections~1 and~3 give the upper-bound history
and the weighted-prefix notation compared in \cref{kfp:sec:scope}.

\bibitem{eden}
M.~Eden, \emph{A two-dimensional growth process}, in: Proceedings of the
Fourth Berkeley Symposium on Mathematical Statistics and Probability,
Vol.~4, University of California Press, Berkeley, 1961, 223--239.
\wtag\ Added in the write.

\bibitem{klarnerrivest}
D.~A. Klarner and R.~L. Rivest, \emph{A procedure for improving the upper
bound for the number of $n$-ominoes}, Canadian Journal of Mathematics
\textbf{25} (1973), 585--602.
\wtag\ Added in the write.

\bibitem{barequetshalah}
Gill Barequet and Mira Shalah, \emph{Improved upper bounds on the growth
constants of polyominoes and polycubes}, Algorithmica \textbf{84} (2022),
3559--3586. \url{https://doi.org/10.1007/s00453-022-00948-6}.
\wtag\ Added in the write.

\bibitem{redelmeier}
D.~Hugh Redelmeier, \emph{Counting polyominoes: Yet another attack},
Discrete Mathematics \textbf{36} (1981), 191--203.
\url{https://doi.org/10.1016/0012-365X(81)90237-5}.

\bibitem{oeis}
The OEIS Foundation, \emph{A001168: Number of fixed polyominoes with $n$
cells}, consulted September 2026.
\url{https://oeis.org/A001168}.
Only the positive-size enumerative values are used for comparison; the
entry's historical prose about growth bounds is not used as current evidence.

\bibitem{duffin}
R.~J. Duffin, E.~L. Peterson, and C. Zener,
\emph{Geometric Programming---Theory and Application}, Wiley, 1967.

\bibitem{boyd}
S. Boyd, S.-J. Kim, L. Vandenberghe, and A. Hassibi,
\emph{A tutorial on geometric programming}, Optimization and Engineering
\textbf{8} (2007), 67--127.
\url{https://web.stanford.edu/~boyd/papers/gp_tutorial.html}.

\bibitem{bv}
S. Boyd and L. Vandenberghe, \emph{Convex Optimization},
Cambridge University Press, 2004.
\url{https://web.stanford.edu/~boyd/cvxbook/}.

% Part II (batch 40, manuscript 05). Its entry for Bui's paper is the
% entry bui above.
\bibitem{satprefix}
\repo\ research collection,
\emph{Finite-Prefix Corrections and Dual Certificates for Polyomino Growth},
September 2026, unrefereed AI-assisted report, especially Sections 3 and 8--10.
\href{https://github.com/VladimirReshetnikov/ProveIt/tree/74f7f5bdba1aa728b340509701a0fa4e45e8dbf3/SetTheory/Cardinals/docs/reports/enumerative-combinatorics/polyomino-growth-finite-prefix-corrections}{Pinned repository report}.
\wtag\ This is Part~I of the present report.

\bibitem{satformal}
\repo, \emph{KlarnerConstant formal project},
\texttt{Combinatorics/Polyominoes/}\allowbreak\texttt{KlarnerConstant}.
Formal status used here is the status documented by the source report;
no fresh Lean build was performed for this article.
\href{https://github.com/VladimirReshetnikov/ProveIt/tree/74f7f5bdba1aa728b340509701a0fa4e45e8dbf3/Combinatorics/Polyominoes/KlarnerConstant}{Pinned project directory}.

\bibitem{satmap}
\repo\ finite-prefix report, \texttt{code/model.py},
exact sparse map and pattern definitions, inspected at the above snapshot.
\href{https://github.com/VladimirReshetnikov/ProveIt/blob/74f7f5bdba1aa728b340509701a0fa4e45e8dbf3/SetTheory/Cardinals/docs/reports/enumerative-combinatorics/polyomino-growth-finite-prefix-corrections/code/model.py}{Pinned source file}.

\bibitem{satprofile}
\repo\ finite-prefix report, \texttt{data/profiles\_18.csv},
marked and unmarked counts through size 18.
\href{https://github.com/VladimirReshetnikov/ProveIt/blob/74f7f5bdba1aa728b340509701a0fa4e45e8dbf3/SetTheory/Cardinals/docs/reports/enumerative-combinatorics/polyomino-growth-finite-prefix-corrections/data/profiles_18.csv}{Pinned data file}.

\bibitem{satwk}
T. Winkler and J.-P. Katoen,
\emph{Certificates for Probabilistic Pushdown Automata via Optimistic Value
Iteration}, arXiv:2301.08657v2, February 26, 2023, especially Sections 2--3.
\href{https://arxiv.org/html/2301.08657v2}{Full text}.

\bibitem{satsey}
A. Stewart, K. Etessami, and M. Yannakakis,
\emph{Upper bounds for Newton's method on monotone polynomial systems,
and P-time model checking of probabilistic one-counter automata},
arXiv:1302.3741.
\href{https://arxiv.org/abs/1302.3741}{Author preprint record}.
Used as background attribution, not as an unproved input to our rate theorem.

\bibitem{satdrmota}
M. Drmota, \emph{Systems of functional equations},
Random Structures \& Algorithms \textbf{10} (1997), 103--124.
\href{https://doi.org/10.1002/(SICI)1098-2418(199701/03)10:1/2<103::AID-RSA5>3.0.CO;2-Z}{Publisher record}.
Bibliographic context for classical positive-system singularity theory;
the local analysis needed here is proved in the text.
\end{thebibliography}
\end{document}
